\small
\renewcommand{\arraystretch}{1.12}
\section{Đặc tả phương thức (Methods Description)}
\subsection{Use Case UC1 - Tìm phương tiện và Hub phù hợp}
\subsubsection{searchVehicleAndHub()}
\begin{longtable}{|L{0.24\textwidth}|L{0.70\textwidth}|}
\hline
\textbf{Thuộc tính} & \textbf{Nội dung chi tiết} \\
\hline
\endfirsthead
\hline
\textbf{Thuộc tính} & \textbf{Nội dung chi tiết} \\
\hline
\endhead
\textbf{Tên phương thức} & searchVehicleAndHub() \\
\hline
\textbf{Mục đích} & Tìm các phương tiện điện dùng chung hoặc Mobility Hub phù hợp\newline với vị trí và tiêu chí tìm kiếm của sinh viên, hỗ trợ cả 2 kiểu tìm\newline (theophươngtiện - Normal Flow, hoặc theo Hub - AF2). \\
\hline
\textbf{Tham số đầu vào}\newline \textbf{(@param)} & location: Location - vị trí tìm kiếm (GPS hoặc chọn thủ công, AF1);\newline radiusKm: double - bán kính tìm; searchType: SearchType\newline \{VEHICLE, HUB\}; vehicleType: String - loại xe mong muốn (chỉ\newline áp dụng khi searchType=VEHICLE); minBatteryLevel: double -\newline mứcpin tối thiểu(chỉ ápdụngkhi searchType=VEHICLE). \\
\hline
\textbf{Giá trị trả về (@return)} & List<SearchResultDTO> - danh sách kết quả sắp xếp theo khoảng\newline cách tăng dần, kèm trạng thái và thời điểm cập nhật; danh sách rỗng\newline nếu khôngcó kếtquảphù hợp (AF3). \\
\hline
\textbf{Ngoại lệ (@exception)} & InvalidLocationException (EX1); InvalidSearchCriteriaException\newline (EX2); MapServiceUnavailableException (EX3);\newline DataSourceUnavailableException(EX4). \\
\hline
\textbf{Tiền điều kiện (Pre-}\newline \textbf{conditions)} & Hệ thống đã cấu hình danh mục và toạ độ các Hub thuộc phạm vi\newline phục vụ; khôngbắt buộc sinh viên đăngnhập. \\
\hline
\textbf{Hậu điều kiện (Post-}\newline \textbf{conditions)} & Trả về kết quả phù hợp kèm thời điểm cập nhật, hoặc danh sách\newline rỗng; không ghi/sửa bất kỳdữ liệu nào(thao tác chỉ đọc). \\
\hline
 & function searchVehicleAndHub(location, radiusKm, searchType,\newline vehicleType, minBatteryLevel):\newline if location == null: throw InvalidLocationException\newline if radiusKm <= 0: throw InvalidSearchCriteriaException\newline try:\newline candidates = dataAccess.findWithinRadius(\newline location, radiusKm, searchType)\newline catch DataAccessError: \\
\hline
\textbf{Luồng thuật toán}\newline \textbf{(Pseudo-code)} & throw DataSourceUnavailableException\newline if searchType == VEHICLE:\newline results = filter(candidates,\newline type==vehicleType AND battery>=minBatteryLevel\newline AND status==AVAILABLE)\newline else:\newline results = candidates\newline // kèm số xe/chỗ đỗ/cổng sạc khả dụng mỗi Hub\newline sort results by distance(location) ascending\newline return results \\
\hline
\end{longtable}
\vspace{-0.5\baselineskip}
\subsubsection{resolveSearchLocation()}
\begin{longtable}{|L{0.24\textwidth}|L{0.70\textwidth}|}
\hline
\textbf{Thuộc tính} & \textbf{Nội dung chi tiết} \\
\hline
\endfirsthead
\hline
\textbf{Thuộc tính} & \textbf{Nội dung chi tiết} \\
\hline
\endhead
\textbf{Tên phương thức} & resolveSearchLocation() \\
\hline
\textbf{Mục đích} & Xác định vị trí dùng làm tâm tìm kiếm: ưu tiên vị trí do Sinh viên\newline chọn thủ công trên bản đồ (AF1), nếu không có thì dùng toạ độ GPS\newline của thiết bị(bước 3-4 của UC1). \\
\hline
\textbf{Tham số đầu vào}\newline \textbf{(@param)} & gpsLocation: Location (nullable) - toạ độ GPS lấy từ thiết bị;\newline manualLocation: Location (nullable) - vị trí, khu vực hoặc Hub Sinh\newline viên chọn trên bản đồ. \\
\hline
\textbf{Giá trị trả về (@return)} & Location - vị trí tìm kiếm hợp lệ, nằm trong phạm vi phục vụ của\newline khu đô thị ĐHQG-HCM. \\
\hline
\textbf{Ngoại lệ (@exception)} & InvalidLocationException (EX1) - khi cả 2 tham số đều null hoặc vị\newline trí nằm ngoàiphạm viphục vụ. \\
\hline
\textbf{Tiền điều kiện (Pre-}\newline \textbf{conditions)} & Giao diện tìm kiếm đang hiển thị; phạm vi phục vụ\newline (SERVICE\_AREA)đã được cấu hình. \\
\hline
\textbf{Hậu điều kiện (Post-}\newline \textbf{conditions)} & Trả về đúng 1 vị trí tìm kiếm để hiển thị trên bản đồ (bước 4); không\newline lưu vị trí của Sinh viên vào cơ sở dữ liệu. \\
\hline
\textbf{Luồng thuật toán}\newline \textbf{(Pseudo-code)} & function resolveSearchLocation(gpsLocation, manualLocation):\newline if manualLocation != null:\newline location = manualLocation        // AF1: ưu tiên chọn tay\newline else if gpsLocation != null:\newline location = gpsLocation\newline else:\newline throw InvalidLocationException   // EX1\newline if not SERVICE\_AREA.contains(location):\newline throw InvalidLocationException\newline return location \\
\hline
\end{longtable}
\vspace{-0.5\baselineskip}
\subsubsection{validateSearchCriteria()}
\begin{longtable}{|L{0.24\textwidth}|L{0.70\textwidth}|}
\hline
\textbf{Thuộc tính} & \textbf{Nội dung chi tiết} \\
\hline
\endfirsthead
\hline
\textbf{Thuộc tính} & \textbf{Nội dung chi tiết} \\
\hline
\endhead
\textbf{Tên phương thức} & validateSearchCriteria() \\
\hline
\textbf{Mục đích} & Kiểm tra tính hợp lệ của tiêu chí tìm kiếm trước khi truy vấn dữ liệu\newline (bước 6 của UC1), tránh truy vấn vô nghĩa và chỉ rõ trường cần sửa\newline cho Sinh viên(EX2). \\
\hline
\textbf{Tham số đầu vào}\newline \textbf{(@param)} & criteria: SearchCriteria - gồm radiusKm: double, searchType:\newline SearchType, vehicleType: String (nullable), minBatteryLevel:\newline double(nullable). \\
\hline
\textbf{Giá trị trả về (@return)} & void - không trả giá trị; phương thức kết thúc bình thường nghĩa là\newline tiêu chí hợplệ. \\
\hline
\textbf{Ngoại lệ (@exception)} & InvalidSearchCriteriaException (EX2) - kèm danh sách các trường\newline khônghợplệ đểgiao diện đánh dấu. \\
\hline
\textbf{Tiền điều kiện (Pre-}\newline \textbf{conditions)} & Sinh viên đã nhập tiêu chí và gửi yêu cầu tìm kiếm (bước 5). \\
\hline
\textbf{Hậu điều kiện (Post-}\newline \textbf{conditions)} & Nếu hợp lệ, luồng tiếp tục truy vấn dữ liệu; nếu không hợp lệ, không\newline có truyvấn nào đượcgửi xuốngtầngdữ liệu. \\
\hline
\end{longtable}
\vspace{-0.5\baselineskip}
\begin{longtable}{|L{0.24\textwidth}|L{0.70\textwidth}|}
\hline
\textbf{Thuộc tính} & \textbf{Nội dung chi tiết} \\
\hline
\endfirsthead
\hline
\textbf{Thuộc tính} & \textbf{Nội dung chi tiết} \\
\hline
\endhead
 & function validateSearchCriteria(criteria):\newline errors = []\newline if criteria.radiusKm <= 0 OR\newline criteria.radiusKm > MAX\_SEARCH\_RADIUS\_KM:\newline errors.add("radiusKm") \\
\hline
\textbf{Luồng thuật toán}\newline \textbf{(Pseudo-code)} & if criteria.searchType == null: errors.add("searchType")\newline if criteria.searchType == VEHICLE:\newline if criteria.minBatteryLevel != null AND\newline (minBatteryLevel < 0 OR minBatteryLevel > 100):\newline errors.add("minBatteryLevel")\newline if errors not empty:\newline throw InvalidSearchCriteriaException(errors) \\
\hline
\end{longtable}
\vspace{-0.5\baselineskip}
\subsubsection{displayResultsOnMap()}
\begin{longtable}{|L{0.24\textwidth}|L{0.70\textwidth}|}
\hline
\textbf{Thuộc tính} & \textbf{Nội dung chi tiết} \\
\hline
\endfirsthead
\hline
\textbf{Thuộc tính} & \textbf{Nội dung chi tiết} \\
\hline
\endhead
\textbf{Tên phương thức} & displayResultsOnMap() \\
\hline
\textbf{Mục đích} & Gửi danh sách kết quả tới Map API để hiển thị vị trí Hub trên bản đồ\newline (bước 7 của UC1); khi Map API không khả dụng thì chuyển sang\newline chế độ hiển thị danh sách(EX3). \\
\hline
\textbf{Tham số đầu vào}\newline \textbf{(@param)} & results: List<SearchResultDTO> - kết quả tìm kiếm đã sắp xếp;\newline searchLocation: Location - tâm tìm kiếm. \\
\hline
\textbf{Giá trị trả về (@return)} & MapViewModel - dữ liệu hiển thị gồm các điểm đánh dấu (marker)\newline trên bản đồ, hoặc chế độ LIST\_ONLY nếu bản đồ lỗi. \\
\hline
\textbf{Ngoại lệ (@exception)} & Không ném ngoại lệ ra ngoài: MapServiceUnavailableException từ\newline MapAdapter được bắt và xử lýbằngchế độ dựphòng. \\
\hline
\textbf{Tiền điều kiện (Pre-}\newline \textbf{conditions)} & Đã có kết quả từ searchVehicleAndHub (có thể là danh sách rỗng). \\
\hline
\textbf{Hậu điều kiện (Post-}\newline \textbf{conditions)} & Sinh viên luôn nhận được kết quả, dạng bản đồ hoặc dạng danh sách;\newline khôngthayđổi dữ liệu. \\
\hline
\textbf{Luồng thuật toán}\newline \textbf{(Pseudo-code)} & function displayResultsOnMap(results, searchLocation):\newline view = new MapViewModel(results)\newline try:\newline markers = mapAdapter.renderMarkers(\newline searchLocation, results.hubLocations)\newline view.setMarkers(markers)\newline view.mode = MAP\newline catch MapServiceUnavailableException:   // EX3\newline view.mode = LIST\_ONLY\newline view.notice = "Bản đồ tạm thời không khả dụng"\newline return view \\
\hline
\end{longtable}
\vspace{-0.5\baselineskip}
\subsection{Use Case UC2 - Xem trạng thái phương tiện và Hub}
\subsubsection{getResourceStatus()}
\begin{longtable}{|L{0.24\textwidth}|L{0.70\textwidth}|}
\hline
\textbf{Thuộc tính} & \textbf{Nội dung chi tiết} \\
\hline
\endfirsthead
\hline
\textbf{Thuộc tính} & \textbf{Nội dung chi tiết} \\
\hline
\endhead
\textbf{Tên phương thức} & getResourceStatus() \\
\hline
\textbf{Mục đích} & Lấy trạng thái gần thời gian thực của 1 phương tiện hoặc 1 Mobility\newline Hub cụ thể, phục vụ sinh viên ra quyết định tiếp theo (đặt\newline xe/chỗ/sạc). \\
\hline
\textbf{Tham số đầu vào}\newline \textbf{(@param)} & resourceType: ResourceType \{VEHICLE, HUB\}; resourceId: String\newline - mã định danh tài nguyên. \\
\hline
\textbf{Giá trị trả về (@return)} & ResourceStatusDTO - trạng thái chi tiết kèm thời điểm cập nhật gần\newline nhất; đánh dấu "stale" nếu dữ liệuquá cũ(EX4). \\
\hline
\textbf{Ngoại lệ (@exception)} & ResourceNotFoundException (EX2);\newline DataSourceUnavailableException(EX3). \\
\hline
\textbf{Tiền điều kiện (Pre-}\newline \textbf{conditions)} & resourceId hợp lệ, đã chọn từ kết quả UC1 hoặc nhập trực tiếp; hệ\newline thốngcó hoặc truyxuất được dữ liệu trạngthái. \\
\hline
\textbf{Hậu điều kiện (Post-}\newline \textbf{conditions)} & Trả về đúng trạng thái và thời điểm cập nhật của tài nguyên; không\newline ghi/sửa dữ liệu. \\
\hline
\textbf{Luồng thuật toán}\newline \textbf{(Pseudo-code)} & function getResourceStatus(resourceType, resourceId):\newline if resourceId is empty: throw InvalidArgumentException\newline try:\newline resource = dataAccess.findById(resourceType, resourceId)\newline catch DataAccessError:\newline throw DataSourceUnavailableException\newline if resource == null: throw ResourceNotFoundException\newline status = buildStatusDTO(resource)\newline if resource.lastUpdatedAt older than STALE\_THRESHOLD:\newline status.markAsStale()\newline return status \\
\hline
\end{longtable}
\vspace{-0.5\baselineskip}
\subsubsection{buildVehicleStatus()}
\begin{longtable}{|L{0.24\textwidth}|L{0.70\textwidth}|}
\hline
\textbf{Thuộc tính} & \textbf{Nội dung chi tiết} \\
\hline
\endfirsthead
\hline
\textbf{Thuộc tính} & \textbf{Nội dung chi tiết} \\
\hline
\endhead
\textbf{Tên phương thức} & buildVehicleStatus() \\
\hline
\textbf{Mục đích} & Tổng hợp thông tin trạng thái chi tiết của một phương tiện để hiển\newline thị cho Sinh viên(bước 3 của UC2 - nhánh 1a). \\
\hline
\textbf{Tham số đầu vào}\newline \textbf{(@param)} & vehicle: Vehicle - thực thể phương tiện đã truy xuất từ cơ sở dữ liệu. \\
\hline
\textbf{Giá trị trả về (@return)} & VehicleStatusDTO - gồm loại xe, trạng thái, mức pin, Hub hiện tại,\newline thời điểm cậpnhật và cờ isAvailable. \\
\hline
\textbf{Ngoại lệ (@exception)} & Khôngcó. \\
\hline
\textbf{Tiền điều kiện (Pre-}\newline \textbf{conditions)} & vehicle khác null (đã kiểm tra tồn tại trong getResourceStatus). \\
\hline
\textbf{Hậu điều kiện (Post-}\newline \textbf{conditions)} & Trả về DTO chỉ đọc; không thay đổi thực thể Vehicle. \\
\hline
\end{longtable}
\vspace{-0.5\baselineskip}
\begin{longtable}{|L{0.24\textwidth}|L{0.70\textwidth}|}
\hline
\textbf{Thuộc tính} & \textbf{Nội dung chi tiết} \\
\hline
\endfirsthead
\hline
\textbf{Thuộc tính} & \textbf{Nội dung chi tiết} \\
\hline
\endhead
\textbf{Luồng thuật toán}\newline \textbf{(Pseudo-code)} & function buildVehicleStatus(vehicle):\newline dto = new VehicleStatusDTO()\newline dto.vehicleType = vehicle.vehicleType\newline dto.status = vehicle.status\newline dto.batteryLevel = vehicle.batteryLevel\newline dto.currentHub = vehicle.currentHub.name\newline dto.lastUpdatedAt = vehicle.lastUpdatedAt\newline dto.isAvailable = (vehicle.status == AVAILABLE)\newline return dto \\
\hline
\end{longtable}
\vspace{-0.5\baselineskip}
\subsubsection{buildHubStatus()}
\begin{longtable}{|L{0.24\textwidth}|L{0.70\textwidth}|}
\hline
\textbf{Thuộc tính} & \textbf{Nội dung chi tiết} \\
\hline
\endfirsthead
\hline
\textbf{Thuộc tính} & \textbf{Nội dung chi tiết} \\
\hline
\endhead
\textbf{Tên phương thức} & buildHubStatus() \\
\hline
\textbf{Mục đích} & Tổng hợp trạng thái của một Mobility Hub: số phương tiện khả\newline dụng, chỗ đỗ trống, cổng sạc khả dụng và mức độ sử dụng (bước 3\newline của UC2 - nhánh 1b). \\
\hline
\textbf{Tham số đầu vào}\newline \textbf{(@param)} & hub: MobilityHub - thực thể Hub kèm danh sách phương tiện, chỗ\newline đỗ và cổngsạc. \\
\hline
\textbf{Giá trị trả về (@return)} & HubStatusDTO - gồm availableVehicles, emptySpots,\newline availablePorts, usageRate(0-1)và lastUpdatedAt. \\
\hline
\textbf{Ngoại lệ (@exception)} & Khôngcó. \\
\hline
\textbf{Tiền điều kiện (Pre-}\newline \textbf{conditions)} & hub khác null; hub.capacity đã được cấu hình. \\
\hline
\textbf{Hậu điều kiện (Post-}\newline \textbf{conditions)} & Trả về DTO chỉ đọc; không thay đổi thực thể MobilityHub. \\
\hline
\textbf{Luồng thuật toán}\newline \textbf{(Pseudo-code)} & function buildHubStatus(hub):\newline dto = new HubStatusDTO()\newline dto.availableVehicles = count(hub.vehicles,\newline status==AVAILABLE)\newline dto.emptySpots = count(hub.parkingSpots, status==EMPTY)\newline dto.availablePorts = count(hub.chargingPorts,\newline status==AVAILABLE)\newline dto.usageRate = (hub.capacity - dto.emptySpots) / hub.capacity\newline dto.lastUpdatedAt = hub.lastUpdatedAt\newline return dto \\
\hline
\end{longtable}
\vspace{-0.5\baselineskip}
\subsubsection{isStatusStale()}
\begin{longtable}{|L{0.24\textwidth}|L{0.70\textwidth}|}
\hline
\textbf{Thuộc tính} & \textbf{Nội dung chi tiết} \\
\hline
\endfirsthead
\hline
\textbf{Thuộc tính} & \textbf{Nội dung chi tiết} \\
\hline
\endhead
\textbf{Tên phương thức} & isStatusStale() \\
\hline
\textbf{Mục đích} & Xác định dữ liệu trạng thái của tài nguyên đã quá cũ hay chưa, làm\newline cơ sở để đánh dấu "chưa đáng tin cậy" và không khẳng định tài\newline nguyên còn khả dụng (EX4 của UC2; liênquan NFR02). \\
\hline
\textbf{Tham số đầu vào}\newline \textbf{(@param)} & lastUpdatedAt: DateTime (nullable) - thời điểm cập nhật gần nhất\newline của tài nguyên. \\
\hline
\textbf{Giá trị trả về (@return)} & boolean - true nếu dữ liệu quá cũ hoặc chưa xác định, false nếu còn\newline mới. \\
\hline
\end{longtable}
\vspace{-0.5\baselineskip}
\begin{longtable}{|L{0.24\textwidth}|L{0.70\textwidth}|}
\hline
\textbf{Thuộc tính} & \textbf{Nội dung chi tiết} \\
\hline
\endfirsthead
\hline
\textbf{Thuộc tính} & \textbf{Nội dung chi tiết} \\
\hline
\endhead
\textbf{Ngoại lệ (@exception)} & Khôngcó. \\
\hline
\textbf{Tiền điều kiện (Pre-}\newline \textbf{conditions)} & Hằng số STALE\_THRESHOLD đã được cấu hình (ví dụ 15 giây\newline theo NFR02, hoặc ngưỡngdo nhómquyđịnh). \\
\hline
\textbf{Hậu điều kiện (Post-}\newline \textbf{conditions)} & Không thay đổi dữ liệu; kết quả được dùng để gắn cờ stale cho DTO. \\
\hline
\textbf{Luồng thuật toán} & function isStatusStale(lastUpdatedAt):\newline if lastUpdatedAt == null: \\
\hline
\textbf{(Pseudo-code)} & return true          // trạng thái chưa xác định\newline return(now()- lastUpdatedAt)> STALE\_THRESHOLD \\
\hline
\end{longtable}
\vspace{-0.5\baselineskip}
\subsection{Use Case UC3, UC4, UC5, UC17 - Đặt chỗ và Xác thực}
\subsubsection{BookingConfirmModal.confirmBooking()}
\begin{longtable}{|L{0.24\textwidth}|L{0.70\textwidth}|}
\hline
\textbf{Thuộc tính} & \textbf{Nội dung chi tiết} \\
\hline
\endfirsthead
\hline
\textbf{Thuộc tính} & \textbf{Nội dung chi tiết} \\
\hline
\endhead
\textbf{Tên phương thức} & + confirmBooking(sessionId: String, request: BookingRequest):\newline Booking \\
\hline
\textbf{Mục đích} & Gửiyêu cầu xác nhận đặt xe hoặc chỗ đỗ. \\
\hline
\textbf{Tham số đầu vào}\newline \textbf{(@param)} & sessionId: String – mã phiên đăng nhập.\newline request: BookingRequest – dữ liệu yêu cầu đặt chỗ, gồm tài nguyên,\newline Hub, thờigian và mãyêu cầu để đối soát khi thử lại. \\
\hline
\textbf{Giá trị trả về(@return)} & Booking– Đơn đặt chỗ sau khi xử lý. \\
\hline
\textbf{Ngoại lệ (@exception)} & AuthenticationRequired: phiên không hợp lệ;\newline ValidationError: dữ liệu sai;\newline BookingConflict: tài nguyên không còn khả dụng;\newline StorageError: lỗi lưu hoặc chưa xác định kếtquả. \\
\hline
\textbf{Tiền điều kiện (Pre-}\newline \textbf{conditions)} & Người dùng đã chọn tài nguyên và đăng nhập. \\
\hline
\textbf{Hậu điều kiện (Post-}\newline \textbf{conditions)} & Hiển thị đơn đã lưu hoặc thông báo lỗi. \\
\hline
\textbf{Luồng thuật toán}\newline \textbf{(Pseudo-code)} & 1. Đọc lựa chọn của người dùng.\newline 2. Gửi yêu cầu đến phương thức đặt xe hoặc đặt chỗ đỗ tương ứng.\newline 3. Hiển thị kếtquả; chỉ báo thành côngkhi đơn đã lưu. \\
\hline
\end{longtable}
\vspace{-0.5\baselineskip}
\subsubsection{BookingConfirmModal.cancelBooking()}
\begin{longtable}{|L{0.24\textwidth}|L{0.70\textwidth}|}
\hline
\textbf{Thuộc tính} & \textbf{Nội dung chi tiết} \\
\hline
\endfirsthead
\hline
\textbf{Thuộc tính} & \textbf{Nội dung chi tiết} \\
\hline
\endhead
\textbf{Tên phương thức} & + cancelBooking(): void \\
\hline
\textbf{Mục đích} & Hủylựa chọn trước khi xác nhận. \\
\hline
\textbf{Tham số đầu vào}\newline \textbf{(@param)} & Không có. \\
\hline
\textbf{Giá trị trả về(@return)} & void – Khôngtrả vềgiá trị. \\
\hline
\textbf{Ngoại lệ (@exception)} & Không. \\
\hline
\textbf{Tiền điều kiện (Pre-}\newline \textbf{conditions)} & Yêu cầu chưa được gửi để lưu. \\
\hline
\textbf{Hậu điều kiện (Post-}\newline \textbf{conditions)} & Đóng hộp thoại; không tạo hay hủy đơn đã lưu. \\
\hline
\end{longtable}
\vspace{-0.5\baselineskip}
\begin{longtable}{|L{0.24\textwidth}|L{0.70\textwidth}|}
\hline
\textbf{Thuộc tính} & \textbf{Nội dung chi tiết} \\
\hline
\endfirsthead
\hline
\textbf{Thuộc tính} & \textbf{Nội dung chi tiết} \\
\hline
\endhead
\textbf{Luồng thuật toán} & 1. Xóa lựa chọn tạm. \\
\hline
\textbf{(Pseudo-code)} & 2. Đónghộpthoại và trở về màn hình trước. \\
\hline
\end{longtable}
\vspace{-0.5\baselineskip}
\subsubsection{BookingController.bookVehicle()}
\begin{longtable}{|L{0.24\textwidth}|L{0.70\textwidth}|}
\hline
\textbf{Thuộc tính} & \textbf{Nội dung chi tiết} \\
\hline
\endfirsthead
\hline
\textbf{Thuộc tính} & \textbf{Nội dung chi tiết} \\
\hline
\endhead
\textbf{Tên phương thức} & + bookVehicle(sessionId: String, request: BookingRequest):\newline Booking \\
\hline
\textbf{Mục đích} & Kiểm tra và lưu đơn đặt xe. \\
\hline
\textbf{Tham số đầu vào}\newline \textbf{(@param)} & sessionId: String – mã phiên đăng nhập.\newline request: BookingRequest – dữ liệu yêu cầu đặt chỗ, gồm tài nguyên,\newline Hub, thờigian và mãyêu cầu để đối soát khi thử lại. \\
\hline
\textbf{Giá trị trả về(@return)} & Booking– Đơn đặt chỗ sau khi xử lý. \\
\hline
\textbf{Ngoại lệ (@exception)} & AuthenticationRequired: chưa xác thực;\newline BookingConflict: xe không khả dụng;\newline StorageError: lỗi lưu hoặc chưa xác định kếtquả. \\
\hline
\textbf{Tiền điều kiện (Pre-}\newline \textbf{conditions)} & Có yêu cầu đặt xe; người dùng có phiên sinh viên hợp lệ. \\
\hline
\textbf{Hậu điều kiện (Post-}\newline \textbf{conditions)} & Tạo đơn CONFIRMED, giữ xe RESERVED và ghi lịch sử cùng\newline giao dịch. \\
\hline
\textbf{Luồng thuật toán}\newline \textbf{(Pseudo-code)} & 1. Kiểm tra phiên bằng Session.isValid và quyền bằng User.hasRole;\newline kiểm tra và quyền sinh viên. Nếu hết phiên, giữ lựa chọn tạm và yêu\newline cầu SSO; xóa lựa chọn nếu xác thực thất bại hoặc bị hủy.\newline 2. Khóa xe và gọi Vehicle.isAvailable và kiểm tra lại trạng thái khả\newline dụng, Hub và sự cố.\newline 3. Trong cùng giao dịch, tạo đơn, giữ xe và ghi AuditLog; nếu lỗi\newline trước commit thì hoàn tác.\newline 4. Nếu mất phản hồi, đối soát requestId trước khi thử lại.\newline 5. Trả đơn đã lưu vàgửi thôngbáo xác nhận khi cần. \\
\hline
\end{longtable}
\vspace{-0.5\baselineskip}
\subsubsection{BookingController.bookParking()}
\begin{longtable}{|L{0.24\textwidth}|L{0.70\textwidth}|}
\hline
\textbf{Thuộc tính} & \textbf{Nội dung chi tiết} \\
\hline
\endfirsthead
\hline
\textbf{Thuộc tính} & \textbf{Nội dung chi tiết} \\
\hline
\endhead
\textbf{Tên phương thức} & + bookParking(sessionId: String, request: BookingRequest):\newline Booking \\
\hline
\textbf{Mục đích} & Kiểm tra và lưu đơngiữ chỗ đỗ. \\
\hline
\textbf{Tham số đầu vào}\newline \textbf{(@param)} & sessionId: String – mã phiên đăng nhập.\newline request: BookingRequest – dữ liệu yêu cầu đặt chỗ, gồm tài nguyên,\newline Hub, thờigian và mãyêu cầu để đối soát khi thử lại. \\
\hline
\textbf{Giá trị trả về(@return)} & Booking– Đơn đặt chỗ sau khi xử lý. \\
\hline
\textbf{Ngoại lệ (@exception)} & AuthenticationRequired;\newline ValidationError;\newline BookingConflict: chỗ đỗ không còn phù hợp;\newline StorageError. \\
\hline
\textbf{Tiền điều kiện (Pre-}\newline \textbf{conditions)} & Sinh viên đã đăng nhập và chọn Hub, chỗ đỗ, khoảng thời gian. \\
\hline
\textbf{Hậu điều kiện (Post-}\newline \textbf{conditions)} & Lưu đơn và khoảng giữ chỗ không trùng; ghi AuditLog. \\
\hline
\end{longtable}
\vspace{-0.5\baselineskip}
\begin{longtable}{|L{0.24\textwidth}|L{0.70\textwidth}|}
\hline
\textbf{Thuộc tính} & \textbf{Nội dung chi tiết} \\
\hline
\endfirsthead
\hline
\textbf{Thuộc tính} & \textbf{Nội dung chi tiết} \\
\hline
\endhead
\textbf{Luồng thuật toán}\newline \textbf{(Pseudo-code)} & 1. Kiểm tra phiên bằng Session.isValid và quyền bằng User.hasRole;\newline kiểm tra, quyền và dữ liệu; thời điểm kết thúc phải sau bắt đầu.\newline 2. Khóa chỗ đỗ và kiểm tra lại Hub, khả dụng và các khoảng giữ đã\newline lưu.\newline 3. Nếu có khoảng giao nhau, từ chối yêu cầu.\newline 4. Lưu đơn, giữ khoảng thời gian và ghi lịch sử trong cùng giao dịch.\newline 5. Nếu lỗi, hoàn tác hoặc đối soát requestId khi chưa rõ kết quả; trả\newline thôngbáophù hợp. \\
\hline
\end{longtable}
\vspace{-0.5\baselineskip}
\subsubsection{ActiveRideScreen.confirmPickup()}
\begin{longtable}{|L{0.24\textwidth}|L{0.70\textwidth}|}
\hline
\textbf{Thuộc tính} & \textbf{Nội dung chi tiết} \\
\hline
\endfirsthead
\hline
\textbf{Thuộc tính} & \textbf{Nội dung chi tiết} \\
\hline
\endhead
\textbf{Tên phương thức} & + confirmPickup(sessionId: String, bookingId: String, hubId:\newline String): Booking \\
\hline
\textbf{Mục đích} & Tiếpnhận xác nhận nhận xe môphỏng. \\
\hline
\textbf{Tham số đầu vào}\newline \textbf{(@param)} & sessionId: String – mã phiên đăng nhập.\newline bookingId: String – mã đơn đặt xe.\newline hubId: String– mã Hub nhận xe. \\
\hline
\textbf{Giá trị trả về(@return)} & Booking– Đơn đặt chỗ sau khi xử lý. \\
\hline
\textbf{Ngoại lệ (@exception)} & AuthenticationRequired;\newline InvalidBooking;\newline ResourceUnavailable;\newline StorageError. \\
\hline
\textbf{Tiền điều kiện (Pre-}\newline \textbf{conditions)} & Người dùng đã chọn đơn xe và Hub nhận. \\
\hline
\textbf{Hậu điều kiện (Post-}\newline \textbf{conditions)} & Hiển thị kết quả nhận xe hoặc lỗi. \\
\hline
\textbf{Luồng thuật toán}\newline \textbf{(Pseudo-code)} & 1. Gửi mã đơn và Hub đến TripController.\newline 2. Chờ kết quả xác nhận.\newline 3. Hiển thị thành công khi đơn chuyển PICKED\_UP; nếu chưa hoàn\newline tất, thôngbáo liên hệ vận hành. \\
\hline
\end{longtable}
\vspace{-0.5\baselineskip}
\subsubsection{ActiveRideScreen.cancelPickup()}
\begin{longtable}{|L{0.24\textwidth}|L{0.70\textwidth}|}
\hline
\textbf{Thuộc tính} & \textbf{Nội dung chi tiết} \\
\hline
\endfirsthead
\hline
\textbf{Thuộc tính} & \textbf{Nội dung chi tiết} \\
\hline
\endhead
\textbf{Tên phương thức} & + cancelPickup(): void \\
\hline
\textbf{Mục đích} & Dừngnhận xe trước xác nhận. \\
\hline
\textbf{Tham số đầu vào}\newline \textbf{(@param)} & Không có. \\
\hline
\textbf{Giá trị trả về(@return)} & void – Khôngtrả vềgiá trị. \\
\hline
\textbf{Ngoại lệ (@exception)} & Không. \\
\hline
\textbf{Tiền điều kiện (Pre-}\newline \textbf{conditions)} & Chưa gửi xác nhận nhận xe. \\
\hline
\textbf{Hậu điều kiện (Post-}\newline \textbf{conditions)} & Đơn đặt xe và trạng thái xe không đổi. \\
\hline
\textbf{Luồng thuật toán}\newline \textbf{(Pseudo-code)} & 1. Đóng màn hình xác nhận.\newline 2. Trở về danh sách đơn đặt xe. \\
\hline
\end{longtable}
\vspace{-0.5\baselineskip}
\subsubsection{TripController.confirmPickup()}
\begin{longtable}{|L{0.24\textwidth}|L{0.70\textwidth}|}
\hline
\textbf{Thuộc tính} & \textbf{Nội dung chi tiết} \\
\hline
\endfirsthead
\hline
\textbf{Thuộc tính} & \textbf{Nội dung chi tiết} \\
\hline
\endhead
\textbf{Tên phương thức} & + confirmPickup(sessionId: String, bookingId: String, hubId:\newline String): Booking \\
\hline
\textbf{Mục đích} & Xác minh vàghi nhận việc nhận xe môphỏng. \\
\hline
\textbf{Tham số đầu vào}\newline \textbf{(@param)} & sessionId: String – mã phiên đăng nhập.\newline bookingId: String – mã đơn đặt xe.\newline hubId: String– mã Hub nhận xe. \\
\hline
\textbf{Giá trị trả về(@return)} & Booking– Đơn đặt chỗ sau khi xử lý. \\
\hline
\textbf{Ngoại lệ (@exception)} & AuthenticationRequired;\newline InvalidBooking: sai chủ đơn, Hub hoặc hết hạn;\newline ResourceUnavailable: xe có sự cố;\newline StorageError. \\
\hline
\textbf{Tiền điều kiện (Pre-}\newline \textbf{conditions)} & Có phiên sinh viên và đơn đặt xe. \\
\hline
\textbf{Hậu điều kiện (Post-}\newline \textbf{conditions)} & Thành công: xe IN\_USE, đơn PICKED\_UP, lưu người và thời gian\newline nhận. \\
\hline
\textbf{Luồng thuật toán}\newline \textbf{(Pseudo-code)} & 1. Kiểm tra phiên bằng Session.isValid và quyền bằng User.hasRole;\newline kiểm tra, chủ đơn, hiệu lực đặt xe bằng Booking.isValid, Hub và tình\newline trạng xe dưới khóa.\newline 2. Ghi nhận thao tác mô phỏng và chặn đặt lại xe.\newline 3. Cập nhật xe, đơn và AuditLog; trả đơn khi hoàn tất.\newline 4. Nếu lỗi sau khi ghi nhận, giữ thông tin nhận và chặn đặt lại,\newline chuyển vận hành kiểm tra; không báo thành công.\newline 5. Khi gửi lại yêu cầu, kiểm tra kết quả đã lưu để không nhận xe hai\newline lần. \\
\hline
\end{longtable}
\vspace{-0.5\baselineskip}
\subsubsection{LoginScreen.login()}
\begin{longtable}{|L{0.24\textwidth}|L{0.70\textwidth}|}
\hline
\textbf{Thuộc tính} & \textbf{Nội dung chi tiết} \\
\hline
\endfirsthead
\hline
\textbf{Thuộc tính} & \textbf{Nội dung chi tiết} \\
\hline
\endhead
\textbf{Tên phương thức} & + login(returnTarget: String): String \\
\hline
\textbf{Mục đích} & Mở trangđăngnhậpSSO. \\
\hline
\textbf{Tham số đầu vào}\newline \textbf{(@param)} & returnTarget: String – vị trí nội bộ trở về sau đăng nhập. \\
\hline
\textbf{Giá trị trả về(@return)} & String– URL đăngnhậpSSO. \\
\hline
\textbf{Ngoại lệ (@exception)} & SsoUnavailable: khôngkết nối được SSO. \\
\hline
\textbf{Tiền điều kiện (Pre-}\newline \textbf{conditions)} & Người dùng chọn đăng nhập hoặc cần xác thực để đặt xe. \\
\hline
\textbf{Hậu điều kiện (Post-}\newline \textbf{conditions)} & Trình duyệt chuyển đến SSO. \\
\hline
\textbf{Luồng thuật toán}\newline \textbf{(Pseudo-code)} & 1. Lưu vị trí cần trở về.\newline 2. Lấy URL đăng nhập từ cấu hình SSO.\newline 3. Chuyển người dùng sang SSO để nhập thông tin; không lưu mật\newline khẩu tại Hub. \\
\hline
\end{longtable}
\vspace{-0.5\baselineskip}
\subsubsection{LoginScreen.completeLogin()}
\begin{longtable}{|L{0.24\textwidth}|L{0.70\textwidth}|}
\hline
\textbf{Thuộc tính} & \textbf{Nội dung chi tiết} \\
\hline
\endfirsthead
\hline
\textbf{Thuộc tính} & \textbf{Nội dung chi tiết} \\
\hline
\endhead
\textbf{Tên phương thức} & + completeLogin(credential: String): Session \\
\hline
\textbf{Mục đích} & Tiếpnhận kếtquả SSO và hiển thịgiao diệnphù hợp. \\
\hline
\textbf{Tham số đầu vào}\newline \textbf{(@param)} & credential: String – kết quả xác thực do SSO trả về. \\
\hline
\textbf{Giá trị trả về(@return)} & Session – Phiên đăngnhậphợplệ. \\
\hline
\textbf{Ngoại lệ (@exception)} & AuthenticationFailed;\newline AccessDenied;\newline SsoUnavailable;\newline StorageError. \\
\hline
\textbf{Tiền điều kiện (Pre-}\newline \textbf{conditions)} & SSO trả kết quả xác thực. \\
\hline
\textbf{Hậu điều kiện (Post-}\newline \textbf{conditions)} & Có phiên hợp lệ hoặc thông báo lỗi. \\
\hline
\textbf{Luồng thuật toán}\newline \textbf{(Pseudo-code)} & 1. Gửi kết quả SSO đến AuthController.\newline 2. Khi thành công, trở về chức năng đang dùng hoặc trang theo vai\newline trò.\newline 3. Khi thất bại, báo lỗi; nếu người dùng hủy thì về trang chủ và xóa\newline lựa chọn tạm. \\
\hline
\end{longtable}
\vspace{-0.5\baselineskip}
\subsubsection{AuthController.authenticate()}
\begin{longtable}{|L{0.24\textwidth}|L{0.70\textwidth}|}
\hline
\textbf{Thuộc tính} & \textbf{Nội dung chi tiết} \\
\hline
\endfirsthead
\hline
\textbf{Thuộc tính} & \textbf{Nội dung chi tiết} \\
\hline
\endhead
\textbf{Tên phương thức} & + authenticate(credential: String): Session \\
\hline
\textbf{Mục đích} & Xác minh kếtquả SSO và tạophiên người dùng. \\
\hline
\textbf{Tham số đầu vào}\newline \textbf{(@param)} & credential: String – kết quả xác thực do SSO trả về. \\
\hline
\textbf{Giá trị trả về(@return)} & Session – Phiên đăngnhậphợplệ. \\
\hline
\textbf{Ngoại lệ (@exception)} & AuthenticationFailed: SSO từ chối hoặc kết quả không hợp lệ;\newline AccessDenied: thiếu quyền;\newline SsoUnavailable;\newline StorageError. \\
\hline
\textbf{Tiền điều kiện (Pre-}\newline \textbf{conditions)} & Nhận kết quả từ luồng SSO hợp lệ. \\
\hline
\textbf{Hậu điều kiện (Post-}\newline \textbf{conditions)} & Tạo phiên gắn người dùng và quyền; ghi lịch sử đăng nhập. \\
\hline
\textbf{Luồng thuật toán}\newline \textbf{(Pseudo-code)} & 1. Xác minh kết quả với SSO, bao gồm nguồn trả về và hiệu lực.\newline 2. Tìm tài khoản Hub tương ứng và kiểm tra quyền STUDENT hoặc\newline OPERATOR.\newline 3. Tạo và lưu phiên cùng lịch sử đăng nhập.\newline 4. Trả phiên; nếu bị từ chối, hết thời gian chờ hoặc lỗi lưu thì không\newline cấp phiên mới. \\
\hline
\end{longtable}
\vspace{-0.5\baselineskip}
\subsubsection{Booking.isValid()}
\begin{longtable}{|L{0.24\textwidth}|L{0.70\textwidth}|}
\hline
\textbf{Thuộc tính} & \textbf{Nội dung chi tiết} \\
\hline
\endfirsthead
\hline
\textbf{Thuộc tính} & \textbf{Nội dung chi tiết} \\
\hline
\endhead
\textbf{Tên phương thức} & + isValid(now: DateTime): Boolean \\
\hline
\textbf{Mục đích} & Kiểm tra đơn còn hiệu lực. \\
\hline
\textbf{Tham số đầu vào}\newline \textbf{(@param)} & now: DateTime – Thời gian máy chủ. \\
\hline
\textbf{Giá trị trả về(@return)} & Boolean – true khi thỏa điều kiện; false khi khôngthỏa. \\
\hline
\textbf{Ngoại lệ (@exception)} & Không. \\
\hline
\textbf{Tiền điều kiện (Pre-}\newline \textbf{conditions)} & Đơn được đọc từ dữ liệu hiện tại. \\
\hline
\textbf{Hậu điều kiện (Post-}\newline \textbf{conditions)} & Trả true khi đơn đã xác nhận và thời gian hiện tại còn trong khoảng\newline hiệu lực. \\
\hline
\textbf{Luồng thuật toán}\newline \textbf{(Pseudo-code)} & 1. Kiểm tra trạng thái đơn đã xác nhận.\newline 2. So sánh thời gian hiện tại với thời gian bắt đầu và kết thúc.\newline 3. Trả true nếu hợplệ, ngược lại trả false. \\
\hline
\end{longtable}
\vspace{-0.5\baselineskip}
\subsubsection{Vehicle.isAvailable()}
\begin{longtable}{|L{0.24\textwidth}|L{0.70\textwidth}|}
\hline
\textbf{Thuộc tính} & \textbf{Nội dung chi tiết} \\
\hline
\endfirsthead
\hline
\textbf{Thuộc tính} & \textbf{Nội dung chi tiết} \\
\hline
\endhead
\textbf{Tên phương thức} & + isAvailable(): Boolean \\
\hline
\textbf{Mục đích} & Kiểm tra xe có thể được đặt. \\
\hline
\textbf{Tham số đầu vào}\newline \textbf{(@param)} & Không có. \\
\hline
\textbf{Giá trị trả về(@return)} & Boolean – true khi thỏa điều kiện; false khi khôngthỏa. \\
\hline
\textbf{Ngoại lệ (@exception)} & Không. \\
\hline
\textbf{Tiền điều kiện (Pre-}\newline \textbf{conditions)} & Trạng thái xe đã được cập nhật; kiểm tra dưới khóa khi lưu đơn. \\
\hline
\textbf{Hậu điều kiện (Post-}\newline \textbf{conditions)} & Không thay đổi dữ liệu xe. \\
\hline
\textbf{Luồng thuật toán}\newline \textbf{(Pseudo-code)} & 1. Kiểm tra trạng thái xe là AVAILABLE và xe không bị chặn đặt.\newline 2. Trả true nếu thỏa mãn, ngược lại trả false. \\
\hline
\end{longtable}
\vspace{-0.5\baselineskip}
\subsubsection{Session.isValid()}
\begin{longtable}{|L{0.24\textwidth}|L{0.70\textwidth}|}
\hline
\textbf{Thuộc tính} & \textbf{Nội dung chi tiết} \\
\hline
\endfirsthead
\hline
\textbf{Thuộc tính} & \textbf{Nội dung chi tiết} \\
\hline
\endhead
\textbf{Tên phương thức} & + isValid(now: DateTime): Boolean \\
\hline
\textbf{Mục đích} & Kiểm tra thời hạnphiên. \\
\hline
\textbf{Tham số đầu vào}\newline \textbf{(@param)} & now: DateTime – Thời gian máy chủ. \\
\hline
\textbf{Giá trị trả về(@return)} & Boolean – true khi thỏa điều kiện; false khi khôngthỏa. \\
\hline
\textbf{Ngoại lệ (@exception)} & Không. \\
\hline
\textbf{Tiền điều kiện (Pre-}\newline \textbf{conditions)} & Phiên tồn tại; dùng thời gian máy chủ. \\
\hline
\textbf{Hậu điều kiện (Post-}\newline \textbf{conditions)} & Không thay đổi phiên. \\
\hline
\end{longtable}
\vspace{-0.5\baselineskip}
\begin{longtable}{|L{0.24\textwidth}|L{0.70\textwidth}|}
\hline
\textbf{Thuộc tính} & \textbf{Nội dung chi tiết} \\
\hline
\endfirsthead
\hline
\textbf{Thuộc tính} & \textbf{Nội dung chi tiết} \\
\hline
\endhead
\textbf{Luồng thuật toán} & 1. So sánh thời gian hiện tại với expiresAt. \\
\hline
\textbf{(Pseudo-code)} & 2. Trả true nếu chưa đến thời điểm hết hạn; ngược lại trả false. \\
\hline
\end{longtable}
\vspace{-0.5\baselineskip}
\subsubsection{User.hasRole()}
\begin{longtable}{|L{0.24\textwidth}|L{0.70\textwidth}|}
\hline
\textbf{Thuộc tính} & \textbf{Nội dung chi tiết} \\
\hline
\endfirsthead
\hline
\textbf{Thuộc tính} & \textbf{Nội dung chi tiết} \\
\hline
\endhead
\textbf{Tên phương thức} & + hasRole(requiredRole: String): Boolean \\
\hline
\textbf{Mục đích} & Kiểm tra vai trò cần thiết. \\
\hline
\textbf{Tham số đầu vào}\newline \textbf{(@param)} & requiredRole: String – Vai trò cần kiểm tra. \\
\hline
\textbf{Giá trị trả về(@return)} & Boolean – true khi thỏa điều kiện; false khi khôngthỏa. \\
\hline
\textbf{Ngoại lệ (@exception)} & Không. \\
\hline
\textbf{Tiền điều kiện (Pre-}\newline \textbf{conditions)} & Thông tin vai trò lấy từ hệ thống, không lấy từ dữ liệu tự khai của\newline giao diện. \\
\hline
\textbf{Hậu điều kiện (Post-}\newline \textbf{conditions)} & Không thay đổi quyền. \\
\hline
\textbf{Luồng thuật toán}\newline \textbf{(Pseudo-code)} & 1. Đối chiếu role với vai trò yêu cầu.\newline 2. Trả true nếuphù hợp, ngược lại trả false. \\
\hline
\end{longtable}
\vspace{-0.5\baselineskip}
\subsection{Use Case UC6 - Trả phương tiện}
\subsubsection{returnVehicle()}
\begin{longtable}{|L{0.24\textwidth}|L{0.70\textwidth}|}
\hline
\textbf{Thuộc tính} & \textbf{Nội dung chi tiết} \\
\hline
\endfirsthead
\hline
\textbf{Thuộc tính} & \textbf{Nội dung chi tiết} \\
\hline
\endhead
\textbf{Tên phương thức} & returnVehicle \\
\hline
\textbf{Mục đích} & Thực hiện giao dịch trả phương tiện: kiểm tra hợp lệ, ghi audit log,\newline kết thúc lượt mượn, cậpnhật trạngthái xe vàgửi xác nhận. \\
\hline
\textbf{Tham số đầu vào}\newline \textbf{(@param)} & studentId: string, mã định danh sinh viên trả xe.\newline vehicleId: string, mã định danh phương tiện được trả.\newline hubId: string, mã định danh Hub nơi được trảphươngtiện. \\
\hline
\textbf{Giá trị trả về (@return)} & ReturnRecord: lớp thực thể lưu bản ghi trả xe, bao gồm mã bản ghi\newline thuê và trả xe, sinh viên, phương tiện, mã Hub nhận phương tiện trả\newline về, thờigian trả và trạngthái củaphươngtiện sau khi trả. \\
\hline
\textbf{Ngoại lệ (@exception)} & InvalidReturnException: không nhận diện được xe hoặc lượt mượn,\newline hoặc Hub không thể tiếp nhận.\newline ReturnRecordException: lỗi kết nối hoặc lưu, giao dịch bị hủy. Lỗi\newline gửi thôngbáo khôngném ngoại lệ. \\
\hline
\textbf{Tiền điều kiện (Pre-}\newline \textbf{conditions)} & Sinh viên có phiên đăng nhập hợp lệ và đang sử dụng phương tiện.\newline Phương tiện phải có RentalRecord và đang trong trạng thái IN\_USE.\newline Hub nhận có khả năng tiếp nhận (hasCapacityFor trả về True) hoặc\newline sinh viên đã đặt chỗ đỗ. \\
\hline
\textbf{Hậu điều kiện (Post-}\newline \textbf{conditions)} & Một bản ghi ReturnRecord và AuditLog được tạo.\newline RentalRecord.status = ENDED; Vehicle.status thuộc\newline \{‘AVAILABLE’, ‘WAITING\_CHARGE’, ‘MAINTENANCE’\},\newline Vehicle.currentHubId = hubId. \\
\hline
\end{longtable}
\vspace{-0.5\baselineskip}
\begin{longtable}{|L{0.24\textwidth}|L{0.70\textwidth}|}
\hline
\textbf{Thuộc tính} & \textbf{Nội dung chi tiết} \\
\hline
\endfirsthead
\hline
\textbf{Thuộc tính} & \textbf{Nội dung chi tiết} \\
\hline
\endhead
\textbf{Luồng thuật toán}\newline \textbf{(Pseudo-code)} & rental = validateReturn(studentId, vehicleId);\newline if (!canHubAcceptReturn)\newline throw InvalidReturnException;\newline writeAuditLog(...);\newline rental.endRental(now);\newline status = determineVehicleStatus(vehicleId);\newline vehicle.updateStatus(status);\newline record = new ReturnRecord(...);\newline if(!canCreateReturnRecord)\newline throw ReturnRecordException;\newline if(!ReturnRecordException)\{\newline try \{ NotificationService.sendNotification(...) \};\newline catch \{ logError() \};\newline \}\newline return record; \\
\hline
\end{longtable}
\vspace{-0.5\baselineskip}
\subsection{Use Case UC7 - Đăng ký sạc}
\subsubsection{registerCharging()}
\begin{longtable}{|L{0.24\textwidth}|L{0.70\textwidth}|}
\hline
\textbf{Thuộc tính} & \textbf{Nội dung chi tiết} \\
\hline
\endfirsthead
\hline
\textbf{Thuộc tính} & \textbf{Nội dung chi tiết} \\
\hline
\endhead
\textbf{Tên phương thức} & registerCharging \\
\hline
\textbf{Mục đích} & Thực hiện giao dịch đăng ký chỗ sạc: kiểm tra khung giờ, tạo bản\newline ghi đăngký,giữ cổngsạc, lưu AuditLog. \\
\hline
\textbf{Tham số đầu vào}\newline \textbf{(@param)} & studentId: string, mã sinh viên.\newline vehicleId: string, mã phương tiện đang sử dụng.\newline portId: string, mã cổng sạc muốn mượn.\newline startTime: LocalDateTime, thời gian đăng ký bắt đầu sạc.\newline endTime: LocalDateTime, thờigian đăngkýkết thúc sạc. \\
\hline
\textbf{Giá trị trả về (@return)} & ChargingRegistration: lớp bản ghi đăng ký sạc, bao gồm thông tin\newline đầu vào, mã định danh bản đăng ký, trạng thái đăng ký để ở\newline CONFIRMED và thờigian tạo bản đăngký. \\
\hline
\textbf{Ngoại lệ (@exception)} & IllegalArgumentException: end không sau start.\newline SlotConflictException: khung giờ vừa được người khác đặt.\newline RegistrationSaveException: khônglưu được đăngký. \\
\hline
\textbf{Tiền điều kiện (Pre-}\newline \textbf{conditions)} & Sinh viên có phiên đăng nhập hợp lệ.\newline Phương tiện đang có trạng thái sử dụng bởi sinh viên\newline Cổngsạc tồn tại và khôngở trạngthái FAULTY. \\
\hline
\textbf{Hậu điều kiện (Post-}\newline \textbf{conditions)} & Bản ghi ChargingRegistration và AuditLog được tạo.\newline Cổng được giữ trong khung giờ đã chọn.\newline Yêu cầu thông báo được tạo.\newline Thất bại: mọi thay đổi tạm thời bị hủy, không có đăng ký nào được\newline tạo. \\
\hline
\end{longtable}
\vspace{-0.5\baselineskip}
\begin{longtable}{|L{0.24\textwidth}|L{0.70\textwidth}|}
\hline
\textbf{Thuộc tính} & \textbf{Nội dung chi tiết} \\
\hline
\endfirsthead
\hline
\textbf{Thuộc tính} & \textbf{Nội dung chi tiết} \\
\hline
\endhead
\textbf{Luồng thuật toán}\newline \textbf{(Pseudo-code)} & if (end <= start)\newline throw IllegalArgumentException;\newline if (!port.isAvailableDuring(start, end))\newline throw SlotConflictException;\newline reg = new ChargingRegistration(...);\newline port.reserve();\newline writeAuditLog(...);\newline return reg; \\
\hline
\end{longtable}
\vspace{-0.5\baselineskip}
\subsection{Use Case UC8 - Giám sát mạng lưới Mobility Hub}
\subsubsection{getNetworkOverview()}
\begin{longtable}{|L{0.24\textwidth}|L{0.70\textwidth}|}
\hline
\textbf{Thuộc tính} & \textbf{Nội dung chi tiết} \\
\hline
\endfirsthead
\hline
\textbf{Thuộc tính} & \textbf{Nội dung chi tiết} \\
\hline
\endhead
\textbf{Tên phương thức} & getNetworkOverview() \\
\hline
\textbf{Mục đích} & Cung cấp cho Đơn vị vận hành cái nhìn tổng quan, gần thời gian\newline thực về toàn bộ mạng lưới Mobility Hub để quyết định có cần theo\newline dõi chi tiết hoặc điềuphối. \\
\hline
\textbf{Tham số đầu vào}\newline \textbf{(@param)} & filter: OverviewFilter (nullable) - lọc theo khu vực/trạng thái hoạt\newline động/mức độ sử dụng (AF1). \\
\hline
\textbf{Giá trị trả về (@return)} & NetworkOverviewDTO - số Hub hoạt động, số phương tiện/chỗ\newline đỗ/cổng sạc khả dụng, mức độ sử dụng từng Hub, thời điểm cập nhật\newline gần nhất. \\
\hline
\textbf{Ngoại lệ (@exception)} & UnauthorizedException (EX1); DataSourceUnavailableException\newline (EX2). \\
\hline
\textbf{Tiền điều kiện (Pre-}\newline \textbf{conditions)} & Người dùng đã xác thực qua SSO với vai trò Đơn vị vận hành; các\newline Hub đã được cấu hình. \\
\hline
\textbf{Hậu điều kiện (Post-}\newline \textbf{conditions)} & Trả về dashboard tổng quan; Hub/tài nguyên dữ liệu chưa xác định\newline được đánh dấu riêng (EX3); khôngthayđổi dữ liệu vận hành. \\
\hline
\textbf{Luồng thuật toán}\newline \textbf{(Pseudo-code)} & function getNetworkOverview(filter):\newline if !session.valid OR session.role != OPERATOR:\newline throw UnauthorizedException\newline try:\newline hubs = dataAccess.getAllHubsWithLatestState()\newline catch DataAccessError:\newline throw DataSourceUnavailableException\newline if filter != null: hubs = applyFilter(hubs, filter)\newline for hub in hubs:\newline if hub.config invalid:\newline mark hub "needs review", skip usage calc\newline else:\newline hub.usageRate = computeUsage(hub)\newline return buildOverviewDTO(hubs) \\
\hline
\end{longtable}
\vspace{-0.5\baselineskip}
\subsubsection{validateOperatorAccess()}
\begin{longtable}{|L{0.24\textwidth}|L{0.70\textwidth}|}
\hline
\textbf{Thuộc tính} & \textbf{Nội dung chi tiết} \\
\hline
\endfirsthead
\hline
\textbf{Thuộc tính} & \textbf{Nội dung chi tiết} \\
\hline
\endhead
\textbf{Tên phương thức} & validateOperatorAccess() \\
\hline
\textbf{Mục đích} & Kiểm tra phiên làm việc còn hiệu lực và người dùng có vai trò Đơn\newline vị vận hành trước khi cho phép truy cập dashboard (bước 2 của\newline UC8; NFR10). \\
\hline
\textbf{Tham số đầu vào}\newline \textbf{(@param)} & session: UserSession - phiên đăng nhập hiện tại lấy từ SSO. \\
\hline
\textbf{Giá trị trả về(@return)} & void - kết thúc bình thườngnghĩa là đượcphéptruycập. \\
\hline
\textbf{Ngoại lệ (@exception)} & UnauthorizedException (EX1) - phiên hết hạn, không tồn tại, hoặc\newline vai trò không phải OPERATOR. \\
\hline
\textbf{Tiền điều kiện (Pre-}\newline \textbf{conditions)} & Người dùng đã gửi yêu cầu mở chức năng giám sát. \\
\hline
\textbf{Hậu điều kiện (Post-}\newline \textbf{conditions)} & Nếu bị từ chối, không có dữ liệu giám sát nào được truy xuất; lần\newline truycậpbị từ chối đượcghi vào AuditLog (NFR11). \\
\hline
\textbf{Luồng thuật toán}\newline \textbf{(Pseudo-code)} & function validateOperatorAccess(session):\newline if session == null OR session.isExpired():\newline auditLog.record("ACCESS\_DENIED", session)\newline throw UnauthorizedException\newline if session.role != OPERATOR:\newline auditLog.record("ACCESS\_DENIED", session)\newline throw UnauthorizedException \\
\hline
\end{longtable}
\vspace{-0.5\baselineskip}
\subsubsection{applyOverviewFilter()}
\begin{longtable}{|L{0.24\textwidth}|L{0.70\textwidth}|}
\hline
\textbf{Thuộc tính} & \textbf{Nội dung chi tiết} \\
\hline
\endfirsthead
\hline
\textbf{Thuộc tính} & \textbf{Nội dung chi tiết} \\
\hline
\endhead
\textbf{Tên phương thức} & applyOverviewFilter() \\
\hline
\textbf{Mục đích} & Lọc danh sách Hub trên dashboard theo khu vực, trạng thái hoạt\newline động hoặc mức độ sử dụng (AF1 của UC8); hỗ trợ trường hợp không\newline có Hubphù hợp (AF5). \\
\hline
\textbf{Tham số đầu vào}\newline \textbf{(@param)} & hubs: List<HubOverview> - danh sách Hub đã tổng hợp; filter:\newline OverviewFilter - gồm area: String (nullable), operationStatus:\newline HubStatus(nullable), minUsageRate: double(nullable). \\
\hline
\textbf{Giá trị trả về (@return)} & List<HubOverview> - danh sách Hub thoả bộ lọc; danh sách rỗng\newline nếu khôngcó Hubphù hợp (AF5). \\
\hline
\textbf{Ngoại lệ (@exception)} & Khôngcó. \\
\hline
\textbf{Tiền điều kiện (Pre-}\newline \textbf{conditions)} & Dashboard đã tải dữ liệu thành công (bước 5). \\
\hline
\textbf{Hậu điều kiện (Post-}\newline \textbf{conditions)} & Bộ lọc được giữ nguyên trên giao diện để Đơn vị vận hành điều\newline chỉnh hoặc xoá; khôngthayđổi dữ liệu vận hành. \\
\hline
\end{longtable}
\vspace{-0.5\baselineskip}
\begin{longtable}{|L{0.24\textwidth}|L{0.70\textwidth}|}
\hline
\textbf{Thuộc tính} & \textbf{Nội dung chi tiết} \\
\hline
\endfirsthead
\hline
\textbf{Thuộc tính} & \textbf{Nội dung chi tiết} \\
\hline
\endhead
\textbf{Luồng thuật toán}\newline \textbf{(Pseudo-code)} & function applyOverviewFilter(hubs, filter):\newline result = hubs\newline if filter.area != null:\newline result = filter(result, hub.area == filter.area)\newline if filter.operationStatus != null:\newline result = filter(result,\newline hub.status == filter.operationStatus)\newline if filter.minUsageRate != null:\newline result = filter(result,\newline hub.usageRate >= filter.minUsageRate)\newline return result     // rỗng->giao diện báo AF5 \\
\hline
\end{longtable}
\vspace{-0.5\baselineskip}
\subsubsection{computeHubUsage()}
\begin{longtable}{|L{0.24\textwidth}|L{0.70\textwidth}|}
\hline
\textbf{Thuộc tính} & \textbf{Nội dung chi tiết} \\
\hline
\endfirsthead
\hline
\textbf{Thuộc tính} & \textbf{Nội dung chi tiết} \\
\hline
\endhead
\textbf{Tên phương thức} & computeHubUsage() \\
\hline
\textbf{Mục đích} & Tính mức độ sử dụng và trạng thái hoạt động của từng Hub (bước 4\newline của UC8); xử lý Hub có cấu hình không hợp lệ (EX4) và Hub mất\newline kết nối hoặc dữ liệuquá cũ(EX3, NFR05). \\
\hline
\textbf{Tham số đầu vào}\newline \textbf{(@param)} & hub: MobilityHub - Hub kèm trạng thái mới nhất của phương tiện,\newline chỗ đỗ và cổngsạc. \\
\hline
\textbf{Giá trị trả về (@return)} & HubOverview - gồm usageRate, availableVehicles, emptySpots,\newline availablePorts và cờ trạng thái dữ liệu (OK / NEEDS\_REVIEW /\newline UNKNOWN). \\
\hline
\textbf{Ngoại lệ (@exception)} & Không ném ngoại lệ ra ngoài: Hub lỗi được đánh dấu để không làm\newline gián đoạn các Hub còn lại(NFR05). \\
\hline
\textbf{Tiền điều kiện (Pre-}\newline \textbf{conditions)} & Dữ liệu Hub đã được truy xuất (bước 3). \\
\hline
\textbf{Hậu điều kiện (Post-}\newline \textbf{conditions)} & Hub lỗi cấu hình được đánh dấu NEEDS\_REVIEW và không tham\newline gia tính tổng; Hub mất kết nối được đánh dấu UNKNOWN, hiển thị\newline dữ liệugần nhất. \\
\hline
 & function computeHubUsage(hub):\newline overview = new HubOverview(hub.id, hub.name)\newline if hub.capacity == null OR hub.capacity <= 0:   // EX4\newline overview.dataFlag = NEEDS\_REVIEW\newline return overview\newline if isStatusStale(hub.lastUpdatedAt):             // EX3\newline overview.dataFlag = UNKNOWN \\
\hline
\textbf{Luồng thuật toán}\newline \textbf{(Pseudo-code)} & else:\newline overview.dataFlag = OK\newline overview.availableVehicles =\newline count(hub.vehicles, status==AVAILABLE)\newline overview.emptySpots = count(hub.parkingSpots, status==EMPTY)\newline overview.availablePorts =\newline count(hub.chargingPorts, status==AVAILABLE)\newline overview.usageRate =\newline (hub.capacity - overview.emptySpots) / hub.capacity\newline return overview \\
\hline
\end{longtable}
\vspace{-0.5\baselineskip}
\subsection{Use Case UC9 - Theo dõi tài nguyên}
\subsubsection{getResourceStatus()}
\begin{longtable}{|L{0.24\textwidth}|L{0.70\textwidth}|}
\hline
\textbf{Thuộc tính} & \textbf{Nội dung chi tiết} \\
\hline
\endfirsthead
\hline
\textbf{Thuộc tính} & \textbf{Nội dung chi tiết} \\
\hline
\endhead
\textbf{Tên phương thức} & getResourceStatus \\
\hline
\textbf{Mục đích} & Truy xuất trạng thái tài nguyên của một Hub (số xe, chỗ đỗ, cổng sạc\newline khả dụng và danh sách từng tài nguyên) cùng thời điểm cập nhật;\newline đánh dấu dữ liệuquá cũ hoặc chưa xác định. \\
\hline
\textbf{Tham số đầu vào}\newline \textbf{(@param)} & hubId: string, mã hub cần được theo dõi \\
\hline
\textbf{Giá trị trả về (@return)} & HubResource: lớp bao gồm số tài nguyên khả dụng mỗi loại, thời\newline gian cậpnhật, và trạngthái dữ liệu. \\
\hline
\textbf{Ngoại lệ (@exception)} & HubNotFoundException: Hub không tồn tại.\newline DataUnavailableException: Dữ liệu tài nguyên khôngthể truyxuất. \\
\hline
\textbf{Tiền điều kiện (Pre-}\newline \textbf{conditions)} & UC8 đang chạy, người dùng có phiên hợp lệ với quyền Đơn vị vận\newline hành.\newline Hub có dữ liệu trạngthái hoặc có thể truyxuất. \\
\hline
\textbf{Hậu điều kiện (Post-}\newline \textbf{conditions)} & Trả về trạng thái hiện tại của Hub. Không có dữ liệu vận hành nào bị\newline thayđổi. Dòngcó lastUpdatedquá cũ được đánh dấu stale. \\
\hline
\textbf{Luồng thuật toán}\newline \textbf{(Pseudo-code)} & hub = hubRepo.find(hubId);\newline if(hub = null) throw HubNotFoundException;\newline items = resourceRepo.findByHub(hubId);\newline if(item=null) throw DataUnavailableException;\newline status = new HubResource(hub, items);\newline for(ResourceItem in HubResource.resources)\{\newline availableResourceItems = count(item.status == AVAILABLE for\newline each ResourceType);\newline status.lastUpdated = max(item.lastUpdated);\newline dataState = status.evaluateData(now);\newline return status; \\
\hline
\end{longtable}
\vspace{-0.5\baselineskip}
\subsection{Use Case UC10 - Xem cảnh báo quá tải}
\subsubsection{evaluateHubLoad()}
\begin{longtable}{|L{0.24\textwidth}|L{0.70\textwidth}|}
\hline
\textbf{Thuộc tính} & \textbf{Nội dung chi tiết} \\
\hline
\endfirsthead
\hline
\textbf{Thuộc tính} & \textbf{Nội dung chi tiết} \\
\hline
\endhead
\textbf{Tên phương thức} & evaluateHubLoad \\
\hline
\textbf{Mục đích} & So sánh mức sử dụng từng loại tài nguyên (chỗ đỗ, cổng sạc) của\newline Hub với ngưỡng; tạo hoặc cập nhật cảnh báo quá tải với mức nguy\newline hiểm Medium / High. \\
\hline
\textbf{Tham số đầu vào}\newline \textbf{(@param)} & hubId: string, mã hub cần được đánh giá. \\
\hline
\textbf{Giá trị trả về(@return)} & List<OverloadAlert>, các cảnh báo đangACTIVE của Hub. \\
\hline
\textbf{Ngoại lệ (@exception)} & HubNotFoundException: Hub không tồn tại.\newline DataUnavailableException: Dữ liệu tài nguyên không thể truy xuất\newline hoặc ở trạngthái stale. \\
\hline
\end{longtable}
\vspace{-0.5\baselineskip}
\begin{longtable}{|L{0.24\textwidth}|L{0.70\textwidth}|}
\hline
\textbf{Thuộc tính} & \textbf{Nội dung chi tiết} \\
\hline
\endfirsthead
\hline
\textbf{Thuộc tính} & \textbf{Nội dung chi tiết} \\
\hline
\endhead
\textbf{Tiền điều kiện (Pre-}\newline \textbf{conditions)} & UC8 đang chạy, người dùng có phiên hợp lệ với quyền Đơn vị vận\newline hành.\newline Hub tồn tại và HubResource có DataState = FRESH.\newline Đã có OverloadThreshold cho từngUtilizationType. \\
\hline
\textbf{Hậu điều kiện (Post-}\newline \textbf{conditions)} & Mỗi loại vượt ngưỡng có đúng một cảnh báo ACTIVE (không tạo\newline trùng). Cảnh báo có utilization, threshold, severity, expectedCause,\newline detectedAt. Yêu cầu thôngbáo đượcgửi khi cần ngoài ứngdụng. \\
\hline
\textbf{Luồng thuật toán}\newline \textbf{(Pseudo-code)} & hub = hubRepo.find(hubId);\newline if(hub=null) throw HubNotFoundException;\newline status = getResourceStatus(hubId);\newline if(status.dataState != FRESH) throw DataUnavailableException;\newline for(type in [PARKING\_SPOT, CHARGING\_PORT])\{\newline rate = status.getUtilization(type);\newline t = thresholdRepo.find(type);\newline if(t.isExceeded(rate))\{\newline existing = alertRepo.findActive(hubId, type);\newline if(existing = null)\{\newline existing = new OverloadAlert( hubId, type, rate,\newline t.threshold, t.classify(rate), now);\newline notifyOperator(alert);\newline \} else existing.update(rate, t.classify(rate));\newline return alertRepo.findActiveByHub(hubId); \\
\hline
\end{longtable}
\vspace{-0.5\baselineskip}
\subsection{Use Case UC11 – Điều phối phương tiện (Vehicle Dispatch)}
\subsubsection{loadDispatchView()}
\begin{longtable}{|L{0.24\textwidth}|L{0.70\textwidth}|}
\hline
\textbf{Thuộc tính} & \textbf{Nội dung chi tiết} \\
\hline
\endfirsthead
\hline
\textbf{Thuộc tính} & \textbf{Nội dung chi tiết} \\
\hline
\endhead
\textbf{Tên phương thức} & loadDispatchView() \\
\hline
\textbf{Mục đích} & Khởi tạo và hiển thị màn hình điều phối phương tiện, kiểm tra quyền\newline truy cập Đơn vị vận hành và tải bản đồ/bảng cân bằng trạng thái các\newline Hub. \\
\hline
\textbf{Tham số đầu vào}\newline \textbf{(@param)} & Không có. \\
\hline
\textbf{Giá trị trả về(@return)} & void – Khôngtrả vềgiá trị; cậpnhậtgiao diện màn hình điềuphối. \\
\hline
\textbf{Ngoại lệ (@exception)} & UnauthorizedException – Xảy ra khi người dùng không có vai trò\newline Operator (NFR10). DataSourceUnavailableException – Xảy ra khi\newline khôngthể tải danh sách Hub vàphươngtiện. \\
\hline
\textbf{Tiền điều kiện (Pre-}\newline \textbf{conditions)} & Người dùng đã xác thực qua SSO và chọn chức năng 'Điều phối\newline phươngtiện' trên Web Dashboard. \\
\hline
\textbf{Hậu điều kiện (Post-}\newline \textbf{conditions)} & Giao diện bản đồ/bảng cân bằng Hub và danh sách phương tiện khả\newline dụngtại Hub nguồn được hiển thị. \\
\hline
\textbf{Luồng thuật toán}\newline \textbf{(Pseudo-code)} & 1. Gọi DispatchController.getDispatchHubData(null, null) để kiểm\newline tra quyền và tải trạng thái các Hub.\newline 2. Nếu kiểm tra quyền thất bại -> gọi\newline displayDispatchFeedback("Không đủ quyền truy cập", true).\newline 3. Nếu tải thành công -> hiển thị bản đồ mạng lưới Hub, danh sách\newline Hub nguồn/đích và chophépOperator thao tác. \\
\hline
\end{longtable}
\vspace{-0.5\baselineskip}
\subsubsection{prefillFromSimulation(sourceHubId, destHubId, vehicleIds)}
\begin{longtable}{|L{0.24\textwidth}|L{0.70\textwidth}|}
\hline
\textbf{Thuộc tính} & \textbf{Nội dung chi tiết} \\
\hline
\endfirsthead
\hline
\textbf{Thuộc tính} & \textbf{Nội dung chi tiết} \\
\hline
\endhead
\textbf{Tên phương thức} & prefillFromSimulation(sourceHubId, destHubId, vehicleIds) \\
\hline
\textbf{Mục đích} & Tự động điền sẵn thông tin Hub nguồn, Hub đích và danh sách xe đề\newline xuất khi Đơn vị vận hành chuyển tiếp phương án từ mô phỏng What-\newline If(UC14 – AF1). \\
\hline
\textbf{Tham số đầu vào}\newline \textbf{(@param)} & sourceHubId: String – Mã Hub xuất phát được đề xuất.\newline destHubId: String – Mã Hub đích được đề xuất.\newline vehicleIds: List<String> – Danh sách mãphươngtiện cần điềuphối. \\
\hline
\textbf{Giá trị trả về(@return)} & void – Khôngtrả vềgiá trị. \\
\hline
\textbf{Ngoại lệ (@exception)} & InvalidArgumentException – Xảy ra khi tham số đề xuất từ UC14 bị\newline thiếu hoặc rỗng. \\
\hline
\textbf{Tiền điều kiện (Pre-}\newline \textbf{conditions)} & Đơn vị vận hành chọn 'Áp dụng vào thực tế' từ báo cáo mô phỏng\newline UC14 với đề xuất điềuphối xe. \\
\hline
\textbf{Hậu điều kiện (Post-}\newline \textbf{conditions)} & Các thuộc tính selectedSourceHubId, selectedDestHubId,\newline selectedVehicleIds được gán sẵn trên form để Operator kiểm tra\newline trước khi xác nhận. \\
\hline
\textbf{Luồng thuật toán}\newline \textbf{(Pseudo-code)} & 1. Kiểm tra tính hợp lệ của sourceHubId, destHubId và vehicleIds.\newline 2. Gán selectedSourceHubId = sourceHubId; selectedDestHubId =\newline destHubId; selectedVehicleIds = vehicleIds.\newline 3. Gọi DispatchController.getDispatchHubData(sourceHubId,\newline destHubId) để nạp trạng thái thực tế mới nhất.\newline 4. Hiển thị biểu mẫu đã điền sẵn để Operator rà soát hoặc chỉnh sửa. \\
\hline
\end{longtable}
\vspace{-0.5\baselineskip}
\subsubsection{submitDispatchOrder(sourceHubId, destHubId, vehicleIds)}
\begin{longtable}{|L{0.24\textwidth}|L{0.70\textwidth}|}
\hline
\textbf{Thuộc tính} & \textbf{Nội dung chi tiết} \\
\hline
\endfirsthead
\hline
\textbf{Thuộc tính} & \textbf{Nội dung chi tiết} \\
\hline
\endhead
\textbf{Tên phương thức} & submitDispatchOrder(sourceHubId, destHubId, vehicleIds) \\
\hline
\textbf{Mục đích} & Tiếp nhận thao tác xác nhận lệnh điều phối của Đơn vị vận hành và\newline gửiyêu cầu xử lý giao dịch nguyên tử tới DispatchController. \\
\hline
\textbf{Tham số đầu vào}\newline \textbf{(@param)} & sourceHubId: String – Mã Hub nguồn.\newline destHubId: String – Mã Hub đích.\newline vehicleIds: List<String> – Danh sách mã các phương tiện được chọn\newline để điềuphối. \\
\hline
\textbf{Giá trị trả về(@return)} & void – Khôngtrả vềgiá trị. \\
\hline
\textbf{Ngoại lệ (@exception)} & VehicleConflictException – Xe đã bị Sinh viên đặt trước (EX1).\newline HubCapacityExceededException – Hub đích không đủ chỗ đỗ trống\newline (EX2). \\
\hline
\textbf{Tiền điều kiện (Pre-}\newline \textbf{conditions)} & Đơn vị vận hành đã chọn Hub nguồn, Hub đích hợp lệ và ít nhất một\newline phươngtiện khả dụng. \\
\hline
\textbf{Hậu điều kiện (Post-}\newline \textbf{conditions)} & Yêu cầu điều phối được xử lý; giao diện hiển thị kết quả thành công\newline hoặc cảnh báo xungđột/quá tải. \\
\hline
\end{longtable}
\vspace{-0.5\baselineskip}
\begin{longtable}{|L{0.24\textwidth}|L{0.70\textwidth}|}
\hline
\textbf{Thuộc tính} & \textbf{Nội dung chi tiết} \\
\hline
\endfirsthead
\hline
\textbf{Thuộc tính} & \textbf{Nội dung chi tiết} \\
\hline
\endhead
\textbf{Luồng thuật toán}\newline \textbf{(Pseudo-code)} & 1. Kiểm tra danh sách vehicleIds khác rỗng và sourceHubId !=\newline destHubId.\newline 2. Gọi order =\newline DispatchController.processDispatch(currentOperatorId,\newline sourceHubId, destHubId, vehicleIds).\newline 3. Nếu thành công -> gọi displayDispatchFeedback("Tạo lệnh điều\newline phối " + order.id + " thành công", false).\newline 4. Nếu có ngoại lệ xung đột hoặc hết chỗ -> gọi\newline displayDispatchFeedback(err.message, true) và làm mới danh sách\newline xe/chỗ đỗ. \\
\hline
\end{longtable}
\vspace{-0.5\baselineskip}
\subsubsection{cancelPendingOrder(orderId)}
\begin{longtable}{|L{0.24\textwidth}|L{0.70\textwidth}|}
\hline
\textbf{Thuộc tính} & \textbf{Nội dung chi tiết} \\
\hline
\endfirsthead
\hline
\textbf{Thuộc tính} & \textbf{Nội dung chi tiết} \\
\hline
\endhead
\textbf{Tên phương thức} & cancelPendingOrder(orderId) \\
\hline
\textbf{Mục đích} & Tiếp nhận yêu cầu hủy lệnh điều phối đang ở trạng thái chờ xử lý\newline trước khi vận chuyển(UC11 – AF2). \\
\hline
\textbf{Tham số đầu vào}\newline \textbf{(@param)} & orderId: String – Mã định danh lệnh điều phối cần hủy. \\
\hline
\textbf{Giá trị trả về(@return)} & void – Khôngtrả vềgiá trị. \\
\hline
\textbf{Ngoại lệ (@exception)} & InvalidStateException – Xảy ra khi lệnh điều phối không ở trạng thái\newline PENDING hoặc đã bắt đầu vận chuyển. \\
\hline
\textbf{Tiền điều kiện (Pre-}\newline \textbf{conditions)} & Đơn vị vận hành chọn một lệnh điều phối có trạng thái PENDING\newline và nhấn 'Hủylệnh'. \\
\hline
\textbf{Hậu điều kiện (Post-}\newline \textbf{conditions)} & Yêu cầu hủy được chuyển đến DispatchController; giao diện cập\newline nhật trạngthái lệnh và các xe liênquan. \\
\hline
\textbf{Luồng thuật toán}\newline \textbf{(Pseudo-code)} & 1. Gọi result =\newline DispatchController.cancelDispatchOrder(currentOperatorId,\newline orderId).\newline 2. Nếu result == true -> gọi displayDispatchFeedback("Đã hủy lệnh\newline điều phối " + orderId, false).\newline 3. Nếu phát sinh lỗi -> gọi displayDispatchFeedback("Không thể\newline hủylệnh: " + err.message, true). \\
\hline
\end{longtable}
\vspace{-0.5\baselineskip}
\subsubsection{displayDispatchFeedback(message, isError)}
\begin{longtable}{|L{0.24\textwidth}|L{0.70\textwidth}|}
\hline
\textbf{Thuộc tính} & \textbf{Nội dung chi tiết} \\
\hline
\endfirsthead
\hline
\textbf{Thuộc tính} & \textbf{Nội dung chi tiết} \\
\hline
\endhead
\textbf{Tên phương thức} & displayDispatchFeedback(message, isError) \\
\hline
\textbf{Mục đích} & Hiển thị thông báo kết quả tạo/hủy lệnh điều phối hoặc thông báo lỗi\newline (xungđột xe, Hub đích đầychỗ)lêngiao diện. \\
\hline
\textbf{Tham số đầu vào}\newline \textbf{(@param)} & message: String – Nội dung thông báo cần hiển thị.\newline isError: boolean – Cờ phân biệt thông báo lỗi (true) hay thành công\newline (false). \\
\hline
\textbf{Giá trị trả về(@return)} & void – Khôngtrả vềgiá trị. \\
\hline
\textbf{Ngoại lệ (@exception)} & Khôngcó. \\
\hline
\textbf{Tiền điều kiện (Pre-}\newline \textbf{conditions)} & Hệ thống đã hoàn tất xử lý nghiệp vụ hoặc bắt được ngoại lệ từ\newline DispatchController. \\
\hline
\end{longtable}
\vspace{-0.5\baselineskip}
\begin{longtable}{|L{0.24\textwidth}|L{0.70\textwidth}|}
\hline
\textbf{Thuộc tính} & \textbf{Nội dung chi tiết} \\
\hline
\endfirsthead
\hline
\textbf{Thuộc tính} & \textbf{Nội dung chi tiết} \\
\hline
\endhead
\textbf{Hậu điều kiện (Post-}\newline \textbf{conditions)} & Thông báo rõ ràng được hiển thị tới Đơn vị vận hành trên màn hình\newline DispatchScreen. \\
\hline
\textbf{Luồng thuật toán} & 1. Nếu isError == true -> hiển thị banner cảnh báo màu đỏ kèm\newline message và giữ nguyên lựa chọn hiện tại để thử lại. \\
\hline
\textbf{(Pseudo-code)} & 2. Ngược lại -> hiển thị thông báo thành công màu xanh kèm\newline message và cậpnhật lại bảngdữ liệu Hub. \\
\hline
\end{longtable}
\vspace{-0.5\baselineskip}
\subsubsection{getDispatchHubData(sourceHubId, destHubId)}
\begin{longtable}{|L{0.24\textwidth}|L{0.70\textwidth}|}
\hline
\textbf{Thuộc tính} & \textbf{Nội dung chi tiết} \\
\hline
\endfirsthead
\hline
\textbf{Thuộc tính} & \textbf{Nội dung chi tiết} \\
\hline
\endhead
\textbf{Tên phương thức} & getDispatchHubData(sourceHubId, destHubId) \\
\hline
\textbf{Mục đích} & Kiểm tra quyền Operator (NFR10) và truy xuất dữ liệu cân bằng\newline phươngtiện, mứcpin tại Hub nguồn và sức chứa trốngtại Hub đích. \\
\hline
\textbf{Tham số đầu vào}\newline \textbf{(@param)} & sourceHubId: String (nullable) – Mã Hub nguồn cần xem xe khả\newline dụng.\newline destHubId: String (nullable)– Mã Hub đích cần xem chỗ đỗ trống. \\
\hline
\textbf{Giá trị trả về (@return)} & DispatchHubDataDTO – Đối tượng tổng hợp danh sách Hub, các\newline SharedVehicle khả dụng kèm \% pin tại Hub nguồn và availableSpots\newline của Hub đích. \\
\hline
\textbf{Ngoại lệ (@exception)} & UnauthorizedException – Người dùng không có quyền Operator.\newline HubNotFoundException – Khôngtìm thấyHub theo mã cungcấp. \\
\hline
\textbf{Tiền điều kiện (Pre-}\newline \textbf{conditions)} & Phiên làm việc SSO đang hoạt động; dữ liệu trạng thái các Hub đã\newline được đồngbộ từ UC16. \\
\hline
\textbf{Hậu điều kiện (Post-}\newline \textbf{conditions)} & Trả về dữ liệu chỉ đọc phục vụ hiển thị; không thay đổi trạng thái\newline phươngtiện hayHub. \\
\hline
\textbf{Luồng thuật toán}\newline \textbf{(Pseudo-code)} & 1. Xác thực phiên SSO và vai trò OPERATOR; nếu không hợp lệ ->\newline ném UnauthorizedException.\newline 2. Nạp danh sách toàn bộ MobilityHub đang ở trạng thái ACTIVE.\newline 3. Nếu sourceHubId != null -> truy vấn các SharedVehicle có\newline currentHubId == sourceHubId và status == AVAILABLE.\newline 4. Nếu destHubId != null -> đọc availableSpots hiện tại của Hub\newline đích.\newline 5. Đóng gói vào DispatchHubDataDTO và trả về. \\
\hline
\end{longtable}
\vspace{-0.5\baselineskip}
\subsubsection{processDispatch(operatorId, sourceHubId, destHubId, vehicleIds)}
\begin{longtable}{|L{0.24\textwidth}|L{0.70\textwidth}|}
\hline
\textbf{Thuộc tính} & \textbf{Nội dung chi tiết} \\
\hline
\endfirsthead
\hline
\textbf{Thuộc tính} & \textbf{Nội dung chi tiết} \\
\hline
\endhead
\textbf{Tên phương thức} & processDispatch(operatorId, sourceHubId, destHubId, vehicleIds) \\
\hline
\textbf{Mục đích} & Thực thi quy trình điều phối phương tiện nguyên tử: khóa giao dịch\newline (NFR09), kiểm tra sức chứa Hub đích, chống race condition với đặt\newline xe của sinh viên, tạo DispatchOrder, ghi AuditLog (NFR11) và gửi\newline thôngbáo tới đội vận chuyển. \\
\hline
\textbf{Tham số đầu vào}\newline \textbf{(@param)} & operatorId: String – Mã định danh Đơn vị vận hành.\newline sourceHubId: String – Mã Hub xuất phát.\newline destHubId: String – Mã Hub đích.\newline vehicleIds: List<String> – Danh sách mã xe cần điềuphối. \\
\hline
\textbf{Giá trị trả về (@return)} & DispatchOrder – Bản ghi lệnh điều phối mới được tạo với trạng thái\newline PENDING. \\
\hline
\end{longtable}
\vspace{-0.5\baselineskip}
\begin{longtable}{|L{0.24\textwidth}|L{0.70\textwidth}|}
\hline
\textbf{Thuộc tính} & \textbf{Nội dung chi tiết} \\
\hline
\endfirsthead
\hline
\textbf{Thuộc tính} & \textbf{Nội dung chi tiết} \\
\hline
\endhead
\textbf{Ngoại lệ (@exception)} & VehicleConflictException (EX1) – Một hoặc nhiều xe vừa được\newline Sinh viên đặt trước (status != AVAILABLE).\newline HubCapacityExceededException (EX2) – Hub đích không đủ chỗ đỗ\newline trống.\newline TransactionException – Lỗi khóa hoặc lưugiao dịch. \\
\hline
\textbf{Tiền điều kiện (Pre-}\newline \textbf{conditions)} & Đơn vị vận hành có quyền hợp lệ; sourceHubId != destHubId;\newline vehicleIds khôngrỗng. \\
\hline
\textbf{Hậu điều kiện (Post-}\newline \textbf{conditions)} & Thành công: Các xe chuyển sang DISPATCHING; sức chứa dự kiến\newline của Hub nguồn/đích được điều chỉnh; DispatchOrder và AuditLog\newline được lưu; NotificationService nhận yêu cầu gửi thông báo.\newline Thất bại: Rollback toàn bộ thayđổi, khôngkhóa xe nào. \\
\hline
\textbf{Luồng thuật toán}\newline \textbf{(Pseudo-code)} & 1. Bắt đầu giao dịch nguyên tử (Transaction Lock với timeout =\newline LOCK\_TIMEOUT\_SECONDS).\newline 2. Kiểm tra destHub.hasAvailableCapacity(vehicleIds.size()); nếu\newline false -> rollback và ném HubCapacityExceededException.\newline 3. Với mỗi vId trong vehicleIds:\newline a. Khóa hàng bản ghi SharedVehicle(vId) (SELECT FOR\newline UPDATE).\newline b. Nếu !vehicle.isAvailable() hoặc vehicle.currentHubId !=\newline sourceHubId -> thêm vId vào conflictList.\newline 4. Nếu conflictList không rỗng -> rollback giao dịch và ném\newline VehicleConflictException(conflictList).\newline 5. Cập nhật từng SharedVehicle sang status = DISPATCHING; gọi\newline sourceHub.adjustExpectedCapacity(+count) và\newline destHub.adjustExpectedCapacity(-count).\newline 6. Tạo DispatchOrder(operatorId, sourceHubId, destHubId,\newline vehicleIds, PENDING) và lưu CSDL.\newline 7. Gọi AuditLog.record(operatorId, "CREATE\_DISPATCH",\newline order.id, "Dispatch " + vehicleIds.size() + " vehicles").\newline 8. Commit giao dịch.\newline 9. Gọi NotificationService.sendNotification("FIELD\_TEAM",\newline "Lệnh điều phối mới: " + order.id).\newline 10. Trả về order. \\
\hline
\end{longtable}
\vspace{-0.5\baselineskip}
\subsubsection{cancelDispatchOrder(operatorId, orderId)}
\begin{longtable}{|L{0.24\textwidth}|L{0.70\textwidth}|}
\hline
\textbf{Thuộc tính} & \textbf{Nội dung chi tiết} \\
\hline
\endfirsthead
\hline
\textbf{Thuộc tính} & \textbf{Nội dung chi tiết} \\
\hline
\endhead
\textbf{Tên phương thức} & cancelDispatchOrder(operatorId, orderId) \\
\hline
\textbf{Mục đích} & Hủy lệnh điều phối đang chờ xử lý, mở khóa và khôi phục trạng thái\newline các xe về AVAILABLE tại Hub nguồn, hoàn lại sức chứa dự kiến và\newline ghi AuditLog (UC11 – AF2). \\
\hline
\textbf{Tham số đầu vào}\newline \textbf{(@param)} & operatorId: String – Mã Đơn vị vận hành thực hiện hủy.\newline orderId: String– Mã lệnh điềuphối cần hủy. \\
\hline
\textbf{Giá trị trả về(@return)} & boolean – Trả về true nếu hủythành công. \\
\hline
\textbf{Ngoại lệ (@exception)} & ResourceNotFoundException – Không tìm thấy DispatchOrder.\newline InvalidStateException – Lệnh điều phối đã chuyển sang\newline IN\_TRANSIT hoặc COMPLETED nên khôngthể hủy. \\
\hline
\textbf{Tiền điều kiện (Pre-}\newline \textbf{conditions)} & Đơn vị vận hành có quyền hợp lệ; DispatchOrder tồn tại và đang ở\newline trạngthái PENDING. \\
\hline
\end{longtable}
\vspace{-0.5\baselineskip}
\begin{longtable}{|L{0.24\textwidth}|L{0.70\textwidth}|}
\hline
\textbf{Thuộc tính} & \textbf{Nội dung chi tiết} \\
\hline
\endfirsthead
\hline
\textbf{Thuộc tính} & \textbf{Nội dung chi tiết} \\
\hline
\endhead
\textbf{Hậu điều kiện (Post-}\newline \textbf{conditions)} & DispatchOrder.status = CANCELLED; các xe trong lệnh trở về trạng\newline thái AVAILABLE; sức chứa dự kiến của 2 Hub được hoàn tác;\newline AuditLogđượcghi nhận. \\
\hline
\textbf{Luồng thuật toán}\newline \textbf{(Pseudo-code)} & 1. Bắt đầu giao dịch nguyên tử.\newline 2. Tìm DispatchOrder theo orderId; nếu null -> ném\newline ResourceNotFoundException.\newline 3. Nếu order.status != PENDING -> rollback và ném\newline InvalidStateException.\newline 4. Gọi order.markCancelled().\newline 5. Với mỗi vId trong order.vehicleIds: cập nhật\newline SharedVehicle.updateStatus(AVAILABLE, now()).\newline 6. Hoàn tác sức chứa dự kiến: sourceHub.adjustExpectedCapacity(-\newline count); destHub.adjustExpectedCapacity(+count).\newline 7. Gọi AuditLog.record(operatorId, "CANCEL\_DISPATCH",\newline orderId, "Order cancelled before transit").\newline 8. Commitgiao dịch và trả về true. \\
\hline
\end{longtable}
\vspace{-0.5\baselineskip}
\subsubsection{markCancelled()}
\begin{longtable}{|L{0.24\textwidth}|L{0.70\textwidth}|}
\hline
\textbf{Thuộc tính} & \textbf{Nội dung chi tiết} \\
\hline
\endfirsthead
\hline
\textbf{Thuộc tính} & \textbf{Nội dung chi tiết} \\
\hline
\endhead
\textbf{Tên phương thức} & markCancelled() \\
\hline
\textbf{Mục đích} & Chuyển trạng thái của đối tượng lệnh điều phối sang CANCELLED\newline khi lệnh bị hủytrước khi vận chuyển. \\
\hline
\textbf{Tham số đầu vào}\newline \textbf{(@param)} & Không có. \\
\hline
\textbf{Giá trị trả về(@return)} & void – Khôngtrả vềgiá trị. \\
\hline
\textbf{Ngoại lệ (@exception)} & InvalidStateException – Xảyra nếu status hiện tại khác PENDING. \\
\hline
\textbf{Tiền điều kiện (Pre-}\newline \textbf{conditions)} & Đối tượng DispatchOrder đang ở trạng thái PENDING. \\
\hline
\textbf{Hậu điều kiện (Post-}\newline \textbf{conditions)} & Thuộc tính status được cập nhật thành\newline DispatchStatus.CANCELLED. \\
\hline
\textbf{Luồng thuật toán}\newline \textbf{(Pseudo-code)} & 1. Nếu this.status != DispatchStatus.PENDING -> ném\newline InvalidStateException.\newline 2. Gán this.status = DispatchStatus.CANCELLED. \\
\hline
\end{longtable}
\vspace{-0.5\baselineskip}
\subsubsection{markCompleted()}
\begin{longtable}{|L{0.24\textwidth}|L{0.70\textwidth}|}
\hline
\textbf{Thuộc tính} & \textbf{Nội dung chi tiết} \\
\hline
\endfirsthead
\hline
\textbf{Thuộc tính} & \textbf{Nội dung chi tiết} \\
\hline
\endhead
\textbf{Tên phương thức} & markCompleted() \\
\hline
\textbf{Mục đích} & Đánh dấu lệnh điều phối đã hoàn tất sau khi đội vận chuyển đưa xe\newline đến Hub đích an toàn. \\
\hline
\textbf{Tham số đầu vào}\newline \textbf{(@param)} & Không có. \\
\hline
\textbf{Giá trị trả về(@return)} & void – Khôngtrả vềgiá trị. \\
\hline
\textbf{Ngoại lệ (@exception)} & InvalidStateException – Xảy ra nếu status hiện tại là CANCELLED\newline hoặc COMPLETED. \\
\hline
\textbf{Tiền điều kiện (Pre-}\newline \textbf{conditions)} & Đối tượng DispatchOrder đang ở trạng thái PENDING hoặc\newline IN\_TRANSIT. \\
\hline
\end{longtable}
\vspace{-0.5\baselineskip}
\begin{longtable}{|L{0.24\textwidth}|L{0.70\textwidth}|}
\hline
\textbf{Thuộc tính} & \textbf{Nội dung chi tiết} \\
\hline
\endfirsthead
\hline
\textbf{Thuộc tính} & \textbf{Nội dung chi tiết} \\
\hline
\endhead
\textbf{Hậu điều kiện (Post-}\newline \textbf{conditions)} & Thuộc tính status được cập nhật thành\newline DispatchStatus.COMPLETED. \\
\hline
\textbf{Luồng thuật toán}\newline \textbf{(Pseudo-code)} & 1. Nếu this.status == CANCELLED hoặc this.status ==\newline COMPLETED -> ném InvalidStateException.\newline 2. Gán this.status = DispatchStatus.COMPLETED. \\
\hline
\end{longtable}
\vspace{-0.5\baselineskip}
\subsection{Use Case UC12 – Xử lý các sự cố (Incident Handling)}
\subsubsection{displayIncidentList(incidents)}
\begin{longtable}{|L{0.24\textwidth}|L{0.70\textwidth}|}
\hline
\textbf{Thuộc tính} & \textbf{Nội dung chi tiết} \\
\hline
\endfirsthead
\hline
\textbf{Thuộc tính} & \textbf{Nội dung chi tiết} \\
\hline
\endhead
\textbf{Tên phương thức} & displayIncidentList(incidents) \\
\hline
\textbf{Mục đích} & Hiển thị danh sách các sự cố mới ghi nhận kèm vị trí Hub, loại tài\newline nguyên, trạng thái và mức độ ưu tiên để Đơn vị vận hành lựa chọn\newline xử lý. \\
\hline
\textbf{Tham số đầu vào}\newline \textbf{(@param)} & incidents: List<Incident> – Danh sách sự cố đã được sắp xếp theo\newline mức độ ưu tiên. \\
\hline
\textbf{Giá trị trả về(@return)} & void – Khôngtrả vềgiá trị. \\
\hline
\textbf{Ngoại lệ (@exception)} & DisplayException – Xảy ra khi giao diện không thể render danh sách\newline sự cố. \\
\hline
\textbf{Tiền điều kiện (Pre-}\newline \textbf{conditions)} & Đơn vị vận hành mở chức năng Xử lý sự cố trên Web Dashboard. \\
\hline
\textbf{Hậu điều kiện (Post-}\newline \textbf{conditions)} & Danh sách sự cố hiển thị rõ ràng, phân biệt mức ưu tiên\newline (LOW/MEDIUM/HIGH/CRITICAL)và trạngthái. \\
\hline
\textbf{Luồng thuật toán}\newline \textbf{(Pseudo-code)} & 1. Kiểm tra danh sách incidents; nếu rỗng -> hiển thị trạng thái\newline 'Không có sự cố tồn đọng'.\newline 2. Với mỗi Incident, hiển thị hubId, resourceId, type, priority, status\newline và reportedAt.\newline 3. Gắn sự kiện chọn dòngđể mở chi tiết sự cố. \\
\hline
\end{longtable}
\vspace{-0.5\baselineskip}
\subsubsection{submitPauseResource(incidentId, reason, fixingMethod)}
\begin{longtable}{|L{0.24\textwidth}|L{0.70\textwidth}|}
\hline
\textbf{Thuộc tính} & \textbf{Nội dung chi tiết} \\
\hline
\endfirsthead
\hline
\textbf{Thuộc tính} & \textbf{Nội dung chi tiết} \\
\hline
\endhead
\textbf{Tên phương thức} & submitPauseResource(incidentId, reason, fixingMethod) \\
\hline
\textbf{Mục đích} & Tiếp nhận thao tác chọn phương án 'Tạm ngưng hoạt động tài\newline nguyên' kèm mô tả nguyên nhân và hướng khắc phục từ Đơn vị vận\newline hành. \\
\hline
\textbf{Tham số đầu vào}\newline \textbf{(@param)} & incidentId: String – Mã sự cố cần xử lý.\newline reason: String – Mô tả nguyên nhân sự cố.\newline fixingMethod: String– Phươngán kỹthuật khắcphục. \\
\hline
\textbf{Giá trị trả về(@return)} & void – Khôngtrả vềgiá trị. \\
\hline
\textbf{Ngoại lệ (@exception)} & InvalidArgumentException – Xảy ra khi reason hoặc fixingMethod\newline rỗng. \\
\hline
\textbf{Tiền điều kiện (Pre-}\newline \textbf{conditions)} & Đơn vị vận hành đang xem chi tiết một sự cố có trạng thái\newline UNCHECKED hoặc IN\_PROGRESS. \\
\hline
\textbf{Hậu điều kiện (Post-}\newline \textbf{conditions)} & Yêu cầu tạm ngưng tài nguyên được chuyển tới IncidentController\newline vàgiao diện cậpnhật kếtquả. \\
\hline
\end{longtable}
\vspace{-0.5\baselineskip}
\begin{longtable}{|L{0.24\textwidth}|L{0.70\textwidth}|}
\hline
\textbf{Thuộc tính} & \textbf{Nội dung chi tiết} \\
\hline
\endfirsthead
\hline
\textbf{Thuộc tính} & \textbf{Nội dung chi tiết} \\
\hline
\endhead
\textbf{Luồng thuật toán} & 1. Kiểm tra reason và fixingMethod không được để trống; nếu trống\newline -> gọi displayIncidentAlert("Vui lòng nhập nguyên nhân và hướng\newline khắc phục", HIGH) và dừng.\newline 2. Gọi updated = \\
\hline
\textbf{(Pseudo-code)} & IncidentController.handlePauseResource(currentOperatorId,\newline incidentId, reason, fixingMethod).\newline 3. Hiển thị thông báo 'Đã khóa tài nguyên và chuyển sang Bảo trì' và\newline làm mới chi tiết sự cố. \\
\hline
\end{longtable}
\vspace{-0.5\baselineskip}
\subsubsection{submitResolveIncident(incidentId)}
\begin{longtable}{|L{0.24\textwidth}|L{0.70\textwidth}|}
\hline
\textbf{Thuộc tính} & \textbf{Nội dung chi tiết} \\
\hline
\endfirsthead
\hline
\textbf{Thuộc tính} & \textbf{Nội dung chi tiết} \\
\hline
\endhead
\textbf{Tên phương thức} & submitResolveIncident(incidentId) \\
\hline
\textbf{Mục đích} & Tiếp nhận thao tác nhấn 'Hoàn tất xử lý' sau khi đội kỹ thuật đã sửa\newline chữa xongtài nguyên bị sự cố(UC12 – AF1). \\
\hline
\textbf{Tham số đầu vào}\newline \textbf{(@param)} & incidentId: String – Mã sự cố đã khắc phục xong. \\
\hline
\textbf{Giá trị trả về(@return)} & void – Khôngtrả vềgiá trị. \\
\hline
\textbf{Ngoại lệ (@exception)} & InvalidStateException – Xảy ra khi sự cố chưa ở trạng thái\newline IN\_PROGRESS. \\
\hline
\textbf{Tiền điều kiện (Pre-}\newline \textbf{conditions)} & Tài nguyên liên quan đang ở trạng thái MAINTENANCE và đội kỹ\newline thuật đã hoàn tất sửa chữa. \\
\hline
\textbf{Hậu điều kiện (Post-}\newline \textbf{conditions)} & Yêu cầu khôi phục tài nguyên được gửi tới IncidentController và\newline giao diện hiển thị xác nhận. \\
\hline
\textbf{Luồng thuật toán}\newline \textbf{(Pseudo-code)} & 1. Gọi resolved =\newline IncidentController.completeIncident(currentOperatorId, incidentId).\newline 2. Cập nhật giao diện hiển thị trạng thái sự cố RESOLVED và trạng\newline thái tài nguyên AVAILABLE.\newline 3. Hiển thị thôngbáo 'Hoàn tất xử lýsự cố thành công'. \\
\hline
\end{longtable}
\vspace{-0.5\baselineskip}
\subsubsection{displayIncidentAlert(message, priority)}
\begin{longtable}{|L{0.24\textwidth}|L{0.70\textwidth}|}
\hline
\textbf{Thuộc tính} & \textbf{Nội dung chi tiết} \\
\hline
\endfirsthead
\hline
\textbf{Thuộc tính} & \textbf{Nội dung chi tiết} \\
\hline
\endhead
\textbf{Tên phương thức} & displayIncidentAlert(message,priority) \\
\hline
\textbf{Mục đích} & Hiển thị cảnh báo sự cố khẩn cấp (đặc biệt khi một Hub mất kết nối\newline hoàn toàn – EX1)hoặc thôngbáo lỗi thao tác trêngiao diện. \\
\hline
\textbf{Tham số đầu vào}\newline \textbf{(@param)} & message: String – Nội dung cảnh báo.\newline priority: PriorityLevel – Mức độ nghiêm trọngcủa cảnh báo. \\
\hline
\textbf{Giá trị trả về(@return)} & void – Khôngtrả vềgiá trị. \\
\hline
\textbf{Ngoại lệ (@exception)} & Khôngcó. \\
\hline
\textbf{Tiền điều kiện (Pre-}\newline \textbf{conditions)} & Hệ thống phát hiện lỗi thao tác hoặc kích hoạt cơ chế cô lập Hub\newline (NFR05). \\
\hline
\textbf{Hậu điều kiện (Post-}\newline \textbf{conditions)} & Cảnh báo nổi bật đúng mức ưu tiên được hiển thị cho Đơn vị vận\newline hành. \\
\hline
\textbf{Luồng thuật toán}\newline \textbf{(Pseudo-code)} & 1. Ánh xạ priority sang màu sắc/nhãn cảnh báo (CRITICAL/HIGH\newline -> đỏ, MEDIUM -> cam, LOW -> xanh).\newline 2. Hiển thị hộpthoại/banner cảnh báo kèm nội dungmessage. \\
\hline
\end{longtable}
\vspace{-0.5\baselineskip}
\subsubsection{getPrioritizedIncidents()}
\begin{longtable}{|L{0.24\textwidth}|L{0.70\textwidth}|}
\hline
\textbf{Thuộc tính} & \textbf{Nội dung chi tiết} \\
\hline
\endfirsthead
\hline
\textbf{Thuộc tính} & \textbf{Nội dung chi tiết} \\
\hline
\endhead
\textbf{Tên phương thức} & getPrioritizedIncidents() \\
\hline
\textbf{Mục đích} & Kiểm tra quyền Operator (NFR10) và truy xuất danh sách sự cố hiện\newline có, ưu tiên các sự cố CRITICAL (như mất kết nối toàn Hub) lên đầu\newline danh sách. \\
\hline
\textbf{Tham số đầu vào}\newline \textbf{(@param)} & Không có. \\
\hline
\textbf{Giá trị trả về (@return)} & List<Incident> – Danh sách các sự cố được sắp xếp theo mức độ ưu\newline tiêngiảm dần và thờigianphát hiện mới nhất. \\
\hline
\textbf{Ngoại lệ (@exception)} & UnauthorizedException – Người dùng không có quyền Operator.\newline DataSourceUnavailableException – Lỗi truyxuất CSDL. \\
\hline
\textbf{Tiền điều kiện (Pre-}\newline \textbf{conditions)} & Đơn vị vận hành đã đăng nhập hợp lệ qua SSO. \\
\hline
\textbf{Hậu điều kiện (Post-}\newline \textbf{conditions)} & Trả về danh sách sự cố đã sắp xếp; không làm thay đổi dữ liệu trạng\newline thái tài nguyên. \\
\hline
\textbf{Luồng thuật toán}\newline \textbf{(Pseudo-code)} & 1. Kiểm tra quyền OPERATOR (NFR10); nếu không hợp lệ -> ném\newline UnauthorizedException.\newline 2. Truy xuất các Incident từ kho dữ liệu.\newline 3. Sắp xếp theo priority giảm dần (CRITICAL > HIGH > MEDIUM\newline > LOW), sau đó theo reportedAt giảm dần.\newline 4. Trả về danh sách kếtquả. \\
\hline
\end{longtable}
\vspace{-0.5\baselineskip}
\subsubsection{handlePauseResource(operatorId, incidentId, reason, fixingMethod)}
\begin{longtable}{|L{0.24\textwidth}|L{0.70\textwidth}|}
\hline
\textbf{Thuộc tính} & \textbf{Nội dung chi tiết} \\
\hline
\endfirsthead
\hline
\textbf{Thuộc tính} & \textbf{Nội dung chi tiết} \\
\hline
\endhead
\textbf{Tên phương thức} & handlePauseResource(operatorId, incidentId, reason, fixingMethod) \\
\hline
\textbf{Mục đích} & Khóa tài nguyên bị lỗi sang trạng thái MAINTENANCE, hủy tự\newline động các đơn đặt xe/lịch sạc đang hoạt động của Sinh viên liên quan\newline tới tài nguyên đó,gửi thôngbáo xin lỗi và lưu AuditLog (NFR11). \\
\hline
\textbf{Tham số đầu vào}\newline \textbf{(@param)} & operatorId: String – Mã Đơn vị vận hành.\newline incidentId: String – Mã sự cố.\newline reason: String – Nguyên nhân sự cố.\newline fixingMethod: String– Hướngkhắcphục. \\
\hline
\textbf{Giá trị trả về (@return)} & Incident – Bản ghi sự cố sau khi đã cập nhật sang trạng thái\newline IN\_PROGRESS. \\
\hline
\textbf{Ngoại lệ (@exception)} & UnauthorizedException – Thiếu quyền Operator.\newline ResourceNotFoundException – Không tìm thấy Incident hoặc\newline Resource.\newline TransactionException – Lỗi cậpnhật nguyên tử. \\
\hline
\textbf{Tiền điều kiện (Pre-}\newline \textbf{conditions)} & Sự cố tồn tại; tài nguyên liên quan chưa ở trạng thái RESOLVED. \\
\hline
\textbf{Hậu điều kiện (Post-}\newline \textbf{conditions)} & Resource.status = MAINTENANCE (loại khỏi tìm kiếm/đặt chỗ);\newline các VehicleBooking hoặc ChargingRegistration đang hiệu lực bị\newline hủy; Sinh viên nhận thông báo; AuditLog được lưu tối thiểu 90 ngày\newline (AUDIT\_RETENTION\_DAYS). \\
\hline
\end{longtable}
\vspace{-0.5\baselineskip}
\begin{longtable}{|L{0.24\textwidth}|L{0.70\textwidth}|}
\hline
\textbf{Thuộc tính} & \textbf{Nội dung chi tiết} \\
\hline
\endfirsthead
\hline
\textbf{Thuộc tính} & \textbf{Nội dung chi tiết} \\
\hline
\endhead
 & 1. Xác thực quyền OPERATOR. Bắt đầu giao dịch nguyên tử.\newline 2. Tìm Incident(incidentId) và Resource(incident.resourceId); nếu\newline null -> ném ResourceNotFoundException.\newline 3. Gọi incident.pauseAndAssign(reason, fixingMethod).\newline 4. Gọi resource.updateStatus(ResourceStatus.MAINTENANCE,\newline now()).\newline 5. Kiểm tra các VehicleBooking đang ACTIVE cho resource.id ->\newline gọi booking.cancelByIncident(incidentId) và lưu danh sách\newline studentId bị ảnh hưởng. \\
\hline
\textbf{Luồng thuật toán}\newline \textbf{(Pseudo-code)} & 6. Kiểm tra các ChargingRegistration đang CONFIRMED cho\newline resource.id -> gọi reg.cancelByIncident(incidentId) và lưu danh sách\newline studentId bị ảnh hưởng.\newline 7. Gọi AuditLog.record(operatorId,\newline "PAUSE\_RESOURCE\_INCIDENT", resource.id, reason + " \textbar\{\} " +\newline fixingMethod).\newline 8. Commit giao dịch.\newline 9. Với mỗi studentId bị ảnh hưởng -> gọi\newline NotificationService.sendNotification(studentId, "Lượt đặt bị hủy do\newline sự cố bảo trì tài nguyên").\newline 10. Trả về incident. \\
\hline
\end{longtable}
\vspace{-0.5\baselineskip}
\subsubsection{completeIncident(operatorId, incidentId)}
\begin{longtable}{|L{0.24\textwidth}|L{0.70\textwidth}|}
\hline
\textbf{Thuộc tính} & \textbf{Nội dung chi tiết} \\
\hline
\endfirsthead
\hline
\textbf{Thuộc tính} & \textbf{Nội dung chi tiết} \\
\hline
\endhead
\textbf{Tên phương thức} & completeIncident(operatorId, incidentId) \\
\hline
\textbf{Mục đích} & Xác nhận hoàn tất khắc phục sự cố, khôi phục trạng thái tài nguyên\newline từ MAINTENANCE về AVAILABLE và ghi AuditLog thời điểm\newline hoàn tất(UC12 – AF1). \\
\hline
\textbf{Tham số đầu vào}\newline \textbf{(@param)} & operatorId: String – Mã Đơn vị vận hành.\newline incidentId: String– Mã sự cố cần đóng. \\
\hline
\textbf{Giá trị trả về (@return)} & Incident – Bản ghi sự cố với trạng thái RESOLVED và resolvedAt\newline đã cậpnhật. \\
\hline
\textbf{Ngoại lệ (@exception)} & UnauthorizedException – Thiếu quyền Operator.\newline ResourceNotFoundException – Không tìm thấy sự cố.\newline InvalidStateException – Sự cố đã được đóngtrước đó. \\
\hline
\textbf{Tiền điều kiện (Pre-}\newline \textbf{conditions)} & Sự cố đang ở trạng thái IN\_PROGRESS và tài nguyên đang ở trạng\newline thái MAINTENANCE. \\
\hline
\textbf{Hậu điều kiện (Post-}\newline \textbf{conditions)} & Incident.status = RESOLVED; Resource.status = AVAILABLE;\newline AuditLog ghi nhận thao tác khôiphục tài nguyên. \\
\hline
\textbf{Luồng thuật toán}\newline \textbf{(Pseudo-code)} & 1. Xác thực quyền OPERATOR. Bắt đầu giao dịch.\newline 2. Tìm Incident(incidentId) và Resource(incident.resourceId).\newline 3. Nếu incident.status == RESOLVED -> ném\newline InvalidStateException.\newline 4. Gọi incident.resolve().\newline 5. Gọi resource.updateStatus(ResourceStatus.AVAILABLE, now()).\newline 6. Gọi AuditLog.record(operatorId, "RESOLVE\_INCIDENT",\newline resource.id, "Restored to AVAILABLE").\newline 7. Commitgiao dịch và trả về incident. \\
\hline
\end{longtable}
\vspace{-0.5\baselineskip}
\subsubsection{isolateDisconnectedHub(hubId)}
\begin{longtable}{|L{0.24\textwidth}|L{0.70\textwidth}|}
\hline
\textbf{Thuộc tính} & \textbf{Nội dung chi tiết} \\
\hline
\endfirsthead
\hline
\textbf{Thuộc tính} & \textbf{Nội dung chi tiết} \\
\hline
\endhead
\textbf{Tên phương thức} & isolateDisconnectedHub(hubId) \\
\hline
\textbf{Mục đích} & Kích hoạt kiến trúc cô lập lỗi (NFR05) khi phát hiện một Hub mất\newline kết nối hoàn toàn: tạm ngưng giao dịch riêng tại Hub đó, gán sự cố\newline mức CRITICAL và khônglàmgián đoạn các Hub khác. \\
\hline
\textbf{Tham số đầu vào}\newline \textbf{(@param)} & hubId: String – Mã MobilityHub bị mất kết nối. \\
\hline
\textbf{Giá trị trả về(@return)} & void – Khôngtrả vềgiá trị. \\
\hline
\textbf{Ngoại lệ (@exception)} & ResourceNotFoundException – Không tìm thấy MobilityHub theo\newline hubId. \\
\hline
\textbf{Tiền điều kiện (Pre-}\newline \textbf{conditions)} & Hệ thống giám sát/đồng bộ phát hiện Hub ngừng phản hồi quá\newline ngưỡng quyđịnh. \\
\hline
\textbf{Hậu điều kiện (Post-}\newline \textbf{conditions)} & MobilityHub.status = ISOLATED; một Incident mới với priority =\newline CRITICAL được tạo (nếu chưa có); các Hub khác vẫn hoạt động\newline bình thường; AuditLog ghi nhận sự kiện cô lập. \\
\hline
\textbf{Luồng thuật toán}\newline \textbf{(Pseudo-code)} & 1. Tìm MobilityHub(hubId); nếu null -> ném\newline ResourceNotFoundException.\newline 2. Gọi hub.isolateHub() (chỉ khóa giao dịch tại hubId, không chặn\newline luồng của các Hub khác).\newline 3. Tạo Incident mới với type = "HUB\_DISCONNECTED", priority\newline = CRITICAL, status = UNCHECKED.\newline 4. Gọi AuditLog.record("SYSTEM", "ISOLATE\_HUB", hubId,\newline "Fault isolation triggered (NFR05)").\newline 5. Gọi\newline NotificationService.sendNotification("OPERATOR\_ON\_DUTY",\newline "Cảnh báo CRITICAL: Hub " + hubId + " mất kết nối và đã được cô\newline lập"). \\
\hline
\end{longtable}
\vspace{-0.5\baselineskip}
\subsubsection{Incident.pauseAndAssign(reason, resolution)}
\begin{longtable}{|L{0.24\textwidth}|L{0.70\textwidth}|}
\hline
\textbf{Thuộc tính} & \textbf{Nội dung chi tiết} \\
\hline
\endfirsthead
\hline
\textbf{Thuộc tính} & \textbf{Nội dung chi tiết} \\
\hline
\endhead
\textbf{Tên phương thức} & Incident.pauseAndAssign(reason, resolution) \\
\hline
\textbf{Mục đích} & Cập nhật nguyên nhân, hướng khắc phục và chuyển trạng thái sự cố\newline sangIN\_PROGRESS. \\
\hline
\textbf{Tham số đầu vào}\newline \textbf{(@param)} & reason: String – Nguyên nhân gây ra sự cố.\newline resolution: String– Biệnphápkhắcphục. \\
\hline
\textbf{Giá trị trả về(@return)} & void – Khôngtrả vềgiá trị. \\
\hline
\textbf{Ngoại lệ (@exception)} & InvalidArgumentException – Khi reason hoặc resolution bị rỗng. \\
\hline
\textbf{Tiền điều kiện (Pre-}\newline \textbf{conditions)} & Incident chưa ở trạng thái RESOLVED. \\
\hline
\textbf{Hậu điều kiện (Post-}\newline \textbf{conditions)} & rootCause = reason; resolutionNote = resolution; status =\newline IncidentStatus.IN\_PROGRESS. \\
\hline
\textbf{Luồng thuật toán}\newline \textbf{(Pseudo-code)} & 1. Kiểm tra reason và resolution khác rỗng.\newline 2. Gán this.rootCause = reason; this.resolutionNote = resolution.\newline 3. Gán this.status = IncidentStatus.IN\_PROGRESS. \\
\hline
\end{longtable}
\vspace{-0.5\baselineskip}
\subsubsection{Incident.resolve()}
\begin{longtable}{|L{0.24\textwidth}|L{0.70\textwidth}|}
\hline
\textbf{Thuộc tính} & \textbf{Nội dung chi tiết} \\
\hline
\endfirsthead
\hline
\textbf{Thuộc tính} & \textbf{Nội dung chi tiết} \\
\hline
\endhead
\textbf{Tên phương thức} & Incident.resolve() \\
\hline
\textbf{Mục đích} & Đánh dấu sự cố đã được khắc phục hoàn toàn và ghi nhận mốc thời\newline giangiảiquyết. \\
\hline
\textbf{Tham số đầu vào}\newline \textbf{(@param)} & Không có. \\
\hline
\textbf{Giá trị trả về(@return)} & void – Khôngtrả vềgiá trị. \\
\hline
\textbf{Ngoại lệ (@exception)} & InvalidStateException – Nếu sự cố đã ở trạngthái RESOLVED. \\
\hline
\textbf{Tiền điều kiện (Pre-}\newline \textbf{conditions)} & Incident đang ở trạng thái IN\_PROGRESS hoặc UNCHECKED. \\
\hline
\textbf{Hậu điều kiện (Post-}\newline \textbf{conditions)} & status = IncidentStatus.RESOLVED; resolvedAt = thời điểm hiện\newline tại. \\
\hline
\textbf{Luồng thuật toán}\newline \textbf{(Pseudo-code)} & 1. Nếu this.status == IncidentStatus.RESOLVED -> ném\newline InvalidStateException.\newline 2. Gán this.status = IncidentStatus.RESOLVED; this.resolvedAt =\newline now(). \\
\hline
\end{longtable}
\vspace{-0.5\baselineskip}
\subsubsection{VehicleBooking.cancelByIncident(incidentId)}
\begin{longtable}{|L{0.24\textwidth}|L{0.70\textwidth}|}
\hline
\textbf{Thuộc tính} & \textbf{Nội dung chi tiết} \\
\hline
\endfirsthead
\hline
\textbf{Thuộc tính} & \textbf{Nội dung chi tiết} \\
\hline
\endhead
\textbf{Tên phương thức} & VehicleBooking.cancelByIncident(incidentId) \\
\hline
\textbf{Mục đích} & Hủy tự động lượt đặt xe dùng chung của Sinh viên khi xe gặp sự cố\newline phải đưa vào bảo trì(UC12 – Bước 6). \\
\hline
\textbf{Tham số đầu vào}\newline \textbf{(@param)} & incidentId: String – Mã sự cố làm cơ sở hủy lượt đặt. \\
\hline
\textbf{Giá trị trả về (@return)} & boolean – Trả về true nếu lượt đặt đang hiệu lực được hủy thành\newline công, false nếu lượt đặt đã kết thúc trước đó. \\
\hline
\textbf{Ngoại lệ (@exception)} & Khôngcó. \\
\hline
\textbf{Tiền điều kiện (Pre-}\newline \textbf{conditions)} & Lượt đặt xe đang gắn với phương tiện bị tạm ngưng hoạt động. \\
\hline
\textbf{Hậu điều kiện (Post-}\newline \textbf{conditions)} & Nếu đang PENDING/CONFIRMED -> status chuyển sang\newline CANCELLED\_BY\_INCIDENT. \\
\hline
\textbf{Luồng thuật toán}\newline \textbf{(Pseudo-code)} & 1. Nếu this.status không nằm trong \{"PENDING", "CONFIRMED"\}\newline -> trả về false.\newline 2. Gán this.status = "CANCELLED\_BY\_INCIDENT".\newline 3. Trả về true. \\
\hline
\end{longtable}
\vspace{-0.5\baselineskip}
\subsubsection{ChargingRegistration.cancelByIncident(incidentId)}
\begin{longtable}{|L{0.24\textwidth}|L{0.70\textwidth}|}
\hline
\textbf{Thuộc tính} & \textbf{Nội dung chi tiết} \\
\hline
\endfirsthead
\hline
\textbf{Thuộc tính} & \textbf{Nội dung chi tiết} \\
\hline
\endhead
\textbf{Tên phương thức} & ChargingRegistration.cancelByIncident(incidentId) \\
\hline
\textbf{Mục đích} & Hủy tự động lịch đăng ký sạc của Sinh viên khi cổng sạc gặp sự cố\newline phải tạm ngưngbảo trì(UC12 – Bước 6). \\
\hline
\textbf{Tham số đầu vào}\newline \textbf{(@param)} & incidentId: String – Mã sự cố làm cơ sở hủy lịch sạc. \\
\hline
\end{longtable}
\vspace{-0.5\baselineskip}
\begin{longtable}{|L{0.24\textwidth}|L{0.70\textwidth}|}
\hline
\textbf{Thuộc tính} & \textbf{Nội dung chi tiết} \\
\hline
\endfirsthead
\hline
\textbf{Thuộc tính} & \textbf{Nội dung chi tiết} \\
\hline
\endhead
\textbf{Giá trị trả về (@return)} & boolean – Trả về true nếu lịch sạc đang CONFIRMED được hủy\newline thành công, false nếu không. \\
\hline
\textbf{Ngoại lệ (@exception)} & Khôngcó. \\
\hline
\textbf{Tiền điều kiện (Pre-}\newline \textbf{conditions)} & Lịch đăng ký sạc thuộc về cổng sạc bị tạm ngưng hoạt động. \\
\hline
\textbf{Hậu điều kiện (Post-}\newline \textbf{conditions)} & Nếu đang CONFIRMED -> status chuyển sang\newline CANCELLED\_BY\_INCIDENT. \\
\hline
\textbf{Luồng thuật toán}\newline \textbf{(Pseudo-code)} & 1. Nếu this.status != "CONFIRMED" -> trả về false.\newline 2. Gán this.status = "CANCELLED\_BY\_INCIDENT".\newline 3. Trả về true. \\
\hline
\end{longtable}
\vspace{-0.5\baselineskip}
\subsection{Use Case UC13 - Lập lịch sạc thông minh}
\subsubsection{displaySchedule()}
\begin{longtable}{|L{0.24\textwidth}|L{0.70\textwidth}|}
\hline
\textbf{Thuộc tính} & \textbf{Nội dung chi tiết} \\
\hline
\endfirsthead
\hline
\textbf{Thuộc tính} & \textbf{Nội dung chi tiết} \\
\hline
\endhead
\textbf{Tên phương thức} & displaySchedule() \\
\hline
\textbf{Mục đích} & Hiển thị lịch sạc thông minh do hệ thống đề xuất, bao gồm phương\newline tiện, cổng sạc, thời điểm bắt đầu, thời điểm kết thúc và mức độ ưu\newline tiên. \\
\hline
\textbf{Tham số đầu vào}\newline \textbf{(@param)} & schedule: ChargingSchedule – Lịch sạc cần hiển thị. \\
\hline
\textbf{Giá trị trả về(@return)} & void – Khôngtrả vềgiá trị. \\
\hline
\textbf{Ngoại lệ (@exception)} & DisplayException – Xảy ra khi giao diện không thể hiển thị dữ liệu\newline lịch sạc. \\
\hline
\textbf{Tiền điều kiện (Pre-}\newline \textbf{conditions)} & Hệ thống đã tạo được lịch sạc và dữ liệu lịch hợp lệ. \\
\hline
\textbf{Hậu điều kiện (Post-}\newline \textbf{conditions)} & Lịch sạc được hiển thị đầy đủ trên giao diện để Đơn vị vận hành xem\newline xét. \\
\hline
\textbf{Luồng thuật toán}\newline \textbf{(Pseudo-code)} & 1. Tiếp nhận dữ liệu lịch sạc. 2. Trích xuất thông tin phương tiện,\newline cổng sạc, thời gian và mức độ ưu tiên. 3. Hiển thị lịch sạc trên giao\newline diện. \\
\hline
\end{longtable}
\vspace{-0.5\baselineskip}
\subsubsection{editSchedule()}
\begin{longtable}{|L{0.24\textwidth}|L{0.70\textwidth}|}
\hline
\textbf{Thuộc tính} & \textbf{Nội dung chi tiết} \\
\hline
\endfirsthead
\hline
\textbf{Thuộc tính} & \textbf{Nội dung chi tiết} \\
\hline
\endhead
\textbf{Tên phương thức} & editSchedule() \\
\hline
\textbf{Mục đích} & Cho phép Đơn vị vận hành chỉnh sửa vị trí Hub hoặc thời gian sạc\newline của một sốphươngtiện trước khiphê duyệt lịch. \\
\hline
\textbf{Tham số đầu vào}\newline \textbf{(@param)} & scheduleId: String – Mã lịch sạc cần chỉnh sửa. updatedData:\newline Map<String, Object> – Thông tin điều chỉnh về Hub hoặc thời gian\newline sạc. \\
\hline
\textbf{Giá trị trả về(@return)} & void – Khôngtrả vềgiá trị. \\
\hline
\textbf{Ngoại lệ (@exception)} & InvalidScheduleException – Xảy ra khi thông tin lịch sạc hoặc dữ\newline liệu điều chỉnh khônghợplệ. \\
\hline
\textbf{Tiền điều kiện (Pre-}\newline \textbf{conditions)} & Lịch sạc đã được tạo và hiển thị; lịch chưa được phê duyệt chính\newline thức. \\
\hline
\end{longtable}
\vspace{-0.5\baselineskip}
\begin{longtable}{|L{0.24\textwidth}|L{0.70\textwidth}|}
\hline
\textbf{Thuộc tính} & \textbf{Nội dung chi tiết} \\
\hline
\endfirsthead
\hline
\textbf{Thuộc tính} & \textbf{Nội dung chi tiết} \\
\hline
\endhead
\textbf{Hậu điều kiện (Post-}\newline \textbf{conditions)} & Thông tin điều chỉnh được tiếp nhận và chuyển đến bộ điều khiển để\newline kiểm tra tính hợplệ. \\
\hline
\textbf{Luồng thuật toán} & 1. Đơn vị vận hành chọn lịch cần chỉnh sửa. 2. Nhập thông tin Hub\newline hoặc thời gian sạc mới. 3. Gửi dữ liệu điều chỉnh đến \\
\hline
\textbf{(Pseudo-code)} & ChargingScheduleController. 4. Tiếp nhận kết quả kiểm tra và hiển\newline thị thôngbáo tươngứng. \\
\hline
\end{longtable}
\vspace{-0.5\baselineskip}
\subsubsection{approveSchedule()}
\begin{longtable}{|L{0.24\textwidth}|L{0.70\textwidth}|}
\hline
\textbf{Thuộc tính} & \textbf{Nội dung chi tiết} \\
\hline
\endfirsthead
\hline
\textbf{Thuộc tính} & \textbf{Nội dung chi tiết} \\
\hline
\endhead
\textbf{Tên phương thức} & approveSchedule() \\
\hline
\textbf{Mục đích} & Tiếp nhận thao tác phê duyệt lịch sạc của Đơn vị vận hành sau khi\newline xem xét lịch đề xuất. \\
\hline
\textbf{Tham số đầu vào}\newline \textbf{(@param)} & scheduleId: String – Mã lịch sạc cần phê duyệt. \\
\hline
\textbf{Giá trị trả về(@return)} & void – Khôngtrả vềgiá trị. \\
\hline
\textbf{Ngoại lệ (@exception)} & ApprovalException – Xảy ra khi yêu cầu phê duyệt không hợp lệ\newline hoặc khôngthể xử lý. \\
\hline
\textbf{Tiền điều kiện (Pre-}\newline \textbf{conditions)} & Lịch sạc đã được tạo, hiển thị và đáp ứng điều kiện phê duyệt. \\
\hline
\textbf{Hậu điều kiện (Post-}\newline \textbf{conditions)} & Yêu cầu phê duyệt được chuyển đến ChargingScheduleController để\newline xử lý. \\
\hline
\textbf{Luồng thuật toán}\newline \textbf{(Pseudo-code)} & 1. Đơn vị vận hành chọn chức năng phê duyệt. 2. Gửi mã lịch sạc\newline đến bộ điều khiển. 3. Chờ kết quả xử lý. 4. Hiển thị thông báo tương\newline ứngtrêngiao diện. \\
\hline
\end{longtable}
\vspace{-0.5\baselineskip}
\subsubsection{displayWarning()}
\begin{longtable}{|L{0.24\textwidth}|L{0.70\textwidth}|}
\hline
\textbf{Thuộc tính} & \textbf{Nội dung chi tiết} \\
\hline
\endfirsthead
\hline
\textbf{Thuộc tính} & \textbf{Nội dung chi tiết} \\
\hline
\endhead
\textbf{Tên phương thức} & displayWarning() \\
\hline
\textbf{Mục đích} & Hiển thị cảnh báo khi hệ thống phát hiện tình trạng quá tải tài\newline nguyên sạc hoặc lỗi truyxuất dữ liệu. \\
\hline
\textbf{Tham số đầu vào}\newline \textbf{(@param)} & warningMessage: String – Nội dung cảnh báo cần hiển thị. \\
\hline
\textbf{Giá trị trả về(@return)} & void – Khôngtrả vềgiá trị. \\
\hline
\textbf{Ngoại lệ (@exception)} & DisplayException – Xảy ra khi giao diện không thể hiển thị thông\newline báo cảnh báo. \\
\hline
\textbf{Tiền điều kiện (Pre-}\newline \textbf{conditions)} & Hệ thống đã phát hiện lỗi hoặc tình trạng bất thường và xác định\newline được nội dungcảnh báo. \\
\hline
\textbf{Hậu điều kiện (Post-}\newline \textbf{conditions)} & Đơn vị vận hành nhận được thông tin cảnh báo trên giao diện. \\
\hline
\textbf{Luồng thuật toán}\newline \textbf{(Pseudo-code)} & 1. Tiếp nhận nội dung cảnh báo. 2. Xác định loại cảnh báo. 3. Hiển\newline thị thông báo trên màn hình. 4. Hướng dẫn Đơn vị vận hành thực\newline hiện thao tácphù hợp. \\
\hline
\end{longtable}
\vspace{-0.5\baselineskip}
\subsubsection{generateSchedule()}
\begin{longtable}{|L{0.24\textwidth}|L{0.70\textwidth}|}
\hline
\textbf{Thuộc tính} & \textbf{Nội dung chi tiết} \\
\hline
\endfirsthead
\hline
\textbf{Thuộc tính} & \textbf{Nội dung chi tiết} \\
\hline
\endhead
\textbf{Tên phương thức} & generateSchedule() \\
\hline
\textbf{Mục đích} & Thu thập dữ liệu trạng thái tài nguyên và xây dựng lịch sạc ưu tiên\newline dựa trên mức pin, nhu cầu sử dụng dự kiến, thời gian chờ, trạng thái\newline cổngsạc và sức chứa Hub. \\
\hline
\textbf{Tham số đầu vào}\newline \textbf{(@param)} & Không có. Hệ thống tự truy xuất dữ liệu trạng thái đã được cập nhật\newline từ UC16. \\
\hline
\textbf{Giá trị trả về(@return)} & ChargingSchedule – Lịch sạc được hệ thốngđề xuất. \\
\hline
\textbf{Ngoại lệ (@exception)} & DataRetrievalException – Xảy ra khi không thể truy xuất dữ liệu từ\newline simulator hoặc cảm biến. ResourceOverloadException – Xảy ra khi\newline số lượng phương tiện cần sạc khẩn cấp vượt quá số cổng sạc khả\newline dụng. \\
\hline
\textbf{Tiền điều kiện (Pre-}\newline \textbf{conditions)} & Đơn vị vận hành đã đăng nhập thành công và hệ thống có dữ liệu\newline hợp lệ về mức pin, nhu cầu sử dụng, thời gian chờ, trạng thái cổng\newline sạc và sức chứa Hub. \\
\hline
\textbf{Hậu điều kiện (Post-}\newline \textbf{conditions)} & Lịch sạc đề xuất được tạo dựa trên dữ liệu tài nguyên hiện tại, chưa\newline được lưu chính thức vào hệ thống. \\
\hline
\textbf{Luồng thuật toán}\newline \textbf{(Pseudo-code)} & 1. Thu thập dữ liệu trạng thái phương tiện và tài nguyên sạc. 2. Kiểm\newline tra tính hợp lệ của dữ liệu. 3. Ưu tiên phương tiện có mức pin thấp\newline hoặc có nhu cầu sử dụng gần. 4. Phân bổ cổng sạc và khung thời\newline gian phù hợp. 5. Kiểm tra xung đột cổng sạc, khung giờ và sức chứa\newline Hub. 6. Tạo lịch sạc đề xuất. \\
\hline
\end{longtable}
\vspace{-0.5\baselineskip}
\subsubsection{validateSchedule()}
\begin{longtable}{|L{0.24\textwidth}|L{0.70\textwidth}|}
\hline
\textbf{Thuộc tính} & \textbf{Nội dung chi tiết} \\
\hline
\endfirsthead
\hline
\textbf{Thuộc tính} & \textbf{Nội dung chi tiết} \\
\hline
\endhead
\textbf{Tên phương thức} & validateSchedule() \\
\hline
\textbf{Mục đích} & Kiểm tra tính hợp lệ của lịch sạc, bảo đảm không xảy ra xung đột\newline cổngsạc, trùngkhung giờ hoặc vượtquá sức chứa Hub. \\
\hline
\textbf{Tham số đầu vào}\newline \textbf{(@param)} & schedule: ChargingSchedule – Lịch sạc cần kiểm tra. \\
\hline
\textbf{Giá trị trả về(@return)} & Boolean – Trả về true nếu lịch hợplệ, false nếuphát hiện viphạm. \\
\hline
\textbf{Ngoại lệ (@exception)} & InvalidScheduleException – Xảy ra khi dữ liệu lịch sạc không đầy\newline đủ hoặc khônghợplệ. \\
\hline
\textbf{Tiền điều kiện (Pre-}\newline \textbf{conditions)} & Lịch sạc đã được tạo hoặc có thay đổi từ thao tác chỉnh sửa của Đơn\newline vị vận hành. \\
\hline
\textbf{Hậu điều kiện (Post-}\newline \textbf{conditions)} & Kết quả kiểm tra tính hợp lệ của lịch sạc được xác định, làm cơ sở\newline cho bướcphê duyệt. \\
\hline
\textbf{Luồng thuật toán}\newline \textbf{(Pseudo-code)} & 1. Tiếp nhận lịch sạc cần kiểm tra. 2. Kiểm tra thời gian của từng\newline lượt sạc. 3. Kiểm tra tình trạng khả dụng của cổng sạc. 4. Kiểm tra\newline trùng cổng và khung giờ. 5. Kiểm tra sức chứa Hub. 6. Trả về kết\newline quả kiểm tra. \\
\hline
\end{longtable}
\vspace{-0.5\baselineskip}
\subsubsection{editSchedule()}
\begin{longtable}{|L{0.24\textwidth}|L{0.70\textwidth}|}
\hline
\textbf{Thuộc tính} & \textbf{Nội dung chi tiết} \\
\hline
\endfirsthead
\hline
\textbf{Thuộc tính} & \textbf{Nội dung chi tiết} \\
\hline
\endhead
\textbf{Tên phương thức} & editSchedule() \\
\hline
\textbf{Mục đích} & Xử lý yêu cầu điều chỉnh lịch sạc của Đơn vị vận hành và kiểm tra\newline tính hợplệ của các thayđổi trước khiphê duyệt. \\
\hline
\textbf{Tham số đầu vào}\newline \textbf{(@param)} & scheduleId: String – Mã lịch sạc cần chỉnh sửa. updatedData:\newline Map<String, Object> – Thông tin Hub hoặc thời gian sạc được điều\newline chỉnh. \\
\hline
\textbf{Giá trị trả về(@return)} & ChargingSchedule – Lịch sạc sau khi được cậpnhật hợplệ. \\
\hline
\textbf{Ngoại lệ (@exception)} & InvalidScheduleException – Xảy ra khi thông tin điều chỉnh không\newline hợp lệ. ResourceConflictException – Xảy ra khi cổng sạc hoặc\newline khung giờ được chọn khôngcòn khả dụng. \\
\hline
\textbf{Tiền điều kiện (Pre-}\newline \textbf{conditions)} & Lịch sạc đã được tạo và chưa được phê duyệt chính thức. \\
\hline
\textbf{Hậu điều kiện (Post-}\newline \textbf{conditions)} & Lịch sạc được cập nhật nếu các thay đổi đáp ứng điều kiện hợp lệ;\newline nếu không, lịchgiữ nguyên trạngthái trước đó. \\
\hline
\textbf{Luồng thuật toán}\newline \textbf{(Pseudo-code)} & 1. Tiếp nhận yêu cầu chỉnh sửa. 2. Xác định lịch cần cập nhật. 3. Áp\newline dụng thông tin điều chỉnh. 4. Gọi validateSchedule() để kiểm tra tính\newline hợp lệ. 5. Nếu hợp lệ, cập nhật lịch. 6. Nếu không hợp lệ, trả về\newline thôngbáo lỗi vàyêu cầu điều chỉnh lại. \\
\hline
\end{longtable}
\vspace{-0.5\baselineskip}
\subsubsection{approveSchedule()}
\begin{longtable}{|L{0.24\textwidth}|L{0.70\textwidth}|}
\hline
\textbf{Thuộc tính} & \textbf{Nội dung chi tiết} \\
\hline
\endfirsthead
\hline
\textbf{Thuộc tính} & \textbf{Nội dung chi tiết} \\
\hline
\endhead
\textbf{Tên phương thức} & approveSchedule() \\
\hline
\textbf{Mục đích} & Xác nhận lịch sạc sau khi được Đơn vị vận hành phê duyệt hoặc\newline được hệ thốngtự độngápdụngtheo chế độphê duyệt tự động. \\
\hline
\textbf{Tham số đầu vào}\newline \textbf{(@param)} & scheduleId: String – Mã lịch sạc cần phê duyệt. approvalMode:\newline String– Chế độphê duyệt,gồm MANUAL hoặc AUTO. \\
\hline
\textbf{Giá trị trả về (@return)} & Boolean – Trả về true nếu lịch được phê duyệt thành công, false nếu\newline khôngthểphê duyệt. \\
\hline
\textbf{Ngoại lệ (@exception)} & ApprovalException – Xảy ra khi lịch sạc không đáp ứng điều kiện\newline phê duyệt. ResourceConflictException – Xảy ra khi tài nguyên sạc\newline phát sinh xungđột tại thời điểmphê duyệt. \\
\hline
\textbf{Tiền điều kiện (Pre-}\newline \textbf{conditions)} & Lịch sạc đã được tạo và vượt qua bước kiểm tra tính hợp lệ. \\
\hline
\textbf{Hậu điều kiện (Post-}\newline \textbf{conditions)} & Lịch sạc được xác nhận phê duyệt và chuyển sang bước lưu chính\newline thức. \\
\hline
\textbf{Luồng thuật toán}\newline \textbf{(Pseudo-code)} & 1. Tiếp nhận yêu cầu phê duyệt và chế độ phê duyệt. 2. Kiểm tra\newline trạng thái lịch sạc. 3. Kiểm tra lại tính hợp lệ và khả dụng của tài\newline nguyên. 4. Xác nhận phê duyệt thủ công hoặc tự động theo cấu hình.\newline 5. Chuyển lịch sangbước lưu chính thức. \\
\hline
\end{longtable}
\vspace{-0.5\baselineskip}
\subsubsection{saveSchedule()}
\begin{longtable}{|L{0.24\textwidth}|L{0.70\textwidth}|}
\hline
\textbf{Thuộc tính} & \textbf{Nội dung chi tiết} \\
\hline
\endfirsthead
\hline
\textbf{Thuộc tính} & \textbf{Nội dung chi tiết} \\
\hline
\endhead
\textbf{Tên phương thức} & saveSchedule() \\
\hline
\textbf{Mục đích} & Lưu lịch sạc đã được phê duyệt, cập nhật trạng thái các cổng sạc\newline thành “đã giữ chỗ” và ghi nhận thao tác vào AuditLog trong cùng\newline mộtgiao dịch nguyên tử. \\
\hline
\textbf{Tham số đầu vào}\newline \textbf{(@param)} & schedule: ChargingSchedule – Lịch sạc đã được phê duyệt cần lưu. \\
\hline
\textbf{Giá trị trả về (@return)} & Boolean – Trả về true nếu toàn bộ giao dịch thành công, false nếu\newline giao dịch thất bại. \\
\hline
\textbf{Ngoại lệ (@exception)} & TransactionException – Xảy ra khi giao dịch lưu lịch, giữ cổng sạc\newline hoặcghi logkhôngthể hoàn tất. \\
\hline
\textbf{Tiền điều kiện (Pre-}\newline \textbf{conditions)} & Lịch sạc đã được phê duyệt và các cổng sạc được phân bổ vẫn khả\newline dụng. \\
\hline
\textbf{Hậu điều kiện (Post-}\newline \textbf{conditions)} & Nếu giao dịch thành công, lịch được lưu, cổng sạc được chuyển sang\newline trạng thái “đã giữ chỗ” và AuditLog được ghi nhận. Nếu giao dịch\newline thất bại, toàn bộ thayđổi trong giao dịch được hoàn tác. \\
\hline
\textbf{Luồng thuật toán}\newline \textbf{(Pseudo-code)} & 1. Bắt đầu giao dịch. 2. Kiểm tra lại khả năng giữ cổng sạc. 3. Lưu\newline lịch sạc. 4. Cập nhật trạng thái các cổng sạc thành “đã giữ chỗ”. 5.\newline Ghi nhận thao tác vào AuditLog. 6. Commit giao dịch nếu tất cả\newline bước thành công. 7. Rollback toàn bộ thayđổi nếu xảyra lỗi. \\
\hline
\end{longtable}
\vspace{-0.5\baselineskip}
\subsection{Use Case UC14 - Mô phỏng và phân tích kịch bản}
\subsubsection{displayInterface()}
\begin{longtable}{|L{0.24\textwidth}|L{0.70\textwidth}|}
\hline
\textbf{Thuộc tính} & \textbf{Nội dung chi tiết} \\
\hline
\endfirsthead
\hline
\textbf{Thuộc tính} & \textbf{Nội dung chi tiết} \\
\hline
\endhead
\textbf{Tên phương thức} & displayInterface() \\
\hline
\textbf{Mục đích} & Hiển thị giao diện mô phỏng What-if, cho phép Đơn vị vận hành\newline thiết lậpcác tham sốgiả định cho kịch bản môphỏng. \\
\hline
\textbf{Tham số đầu vào}\newline \textbf{(@param)} & Không có. \\
\hline
\textbf{Giá trị trả về(@return)} & void - Khôngtrả vềgiá trị. \\
\hline
\textbf{Ngoại lệ (@exception)} & DisplayException - Xảy ra khi hệ thống không thể hiển thị giao diện\newline môphỏng. \\
\hline
\textbf{Tiền điều kiện (Pre-}\newline \textbf{conditions)} & Đơn vị vận hành đã đăng nhập thành công và truy cập chức năng Mô\newline phỏngWhat-if. \\
\hline
\textbf{Hậu điều kiện (Post-}\newline \textbf{conditions)} & Giao diện mô phỏng được hiển thị, sẵn sàng tiếp nhận các tham số\newline giả định. \\
\hline
\textbf{Luồng thuật toán}\newline \textbf{(Pseudo-code)} & 1. Tiếp nhận yêu cầu truy cập chức năng mô phỏng.\newline 2. Khởi tạo giao diện What-if.\newline 3. Hiển thị các trường thiết lập tham số giả định.\newline 4. ChophépĐơn vị vận hành nhậpthôngtin kịch bản. \\
\hline
\end{longtable}
\vspace{-0.5\baselineskip}
\subsubsection{setupParameters()}
\begin{longtable}{|L{0.24\textwidth}|L{0.70\textwidth}|}
\hline
\textbf{Thuộc tính} & \textbf{Nội dung chi tiết} \\
\hline
\endfirsthead
\hline
\textbf{Thuộc tính} & \textbf{Nội dung chi tiết} \\
\hline
\endhead
\textbf{Tên phương thức} & setupParameters() \\
\hline
\textbf{Mục đích} & Tiếp nhận và thiết lập các tham số giả định cho kịch bản mô phỏng,\newline chẳng hạn tăng số lượng xe, giảm số cổng sạc hoặc đóng cửa một\newline Hub. \\
\hline
\textbf{Tham số đầu vào}\newline \textbf{(@param)} & parameters: Map<String, Object> - Tập hợp các tham số giả định do\newline Đơn vị vận hành thiết lập. \\
\hline
\textbf{Giá trị trả về (@return)} & Map<String, Object> - Tập hợp tham số mô phỏng đã được tiếp\newline nhận. \\
\hline
\textbf{Ngoại lệ (@exception)} & InvalidParameterException - Xảy ra khi dữ liệu đầu vào không đúng\newline định dạng. \\
\hline
\textbf{Tiền điều kiện (Pre-}\newline \textbf{conditions)} & Giao diện mô phỏng đã được hiển thị và Đơn vị vận hành đang thiết\newline lậpkịch bản. \\
\hline
\textbf{Hậu điều kiện (Post-}\newline \textbf{conditions)} & Các tham số giả định được tiếp nhận và chuyển đến bộ điều khiển để\newline kiểm tra tính hợplệ. \\
\hline
\textbf{Luồng thuật toán}\newline \textbf{(Pseudo-code)} & 1. Tiếp nhận các tham số giả định.\newline 2. Ghi nhận giá trị của từng tham số.\newline 3. Kiểm tra sơ bộ định dạng dữ liệu.\newline 4. Chuyển tậptham số đến SimulationController. \\
\hline
\end{longtable}
\vspace{-0.5\baselineskip}
\subsubsection{runSimulation()}
\begin{longtable}{|L{0.24\textwidth}|L{0.70\textwidth}|}
\hline
\textbf{Thuộc tính} & \textbf{Nội dung chi tiết} \\
\hline
\endfirsthead
\hline
\textbf{Thuộc tính} & \textbf{Nội dung chi tiết} \\
\hline
\endhead
\textbf{Tên phương thức} & runSimulation() \\
\hline
\textbf{Mục đích} & Tiếp nhận yêu cầu chạy mô phỏng và gửi các tham số kịch bản đến\newline bộ điều khiển để thực hiện môphỏngtrên môi trườngảo cách ly. \\
\hline
\textbf{Tham số đầu vào}\newline \textbf{(@param)} & parameters: Map<String, Object> - Tập hợp tham số giả định cần sử\newline dụngđể chạymôphỏng. \\
\hline
\textbf{Giá trị trả về(@return)} & void - Khôngtrả vềgiá trị. \\
\hline
\textbf{Ngoại lệ (@exception)} & InvalidParameterException - Tham số mô phỏng không hợp lệ.\newline SimulationInitializationException - Không thể khởi tạo môi trường\newline môphỏngcách ly. \\
\hline
\textbf{Tiền điều kiện (Pre-}\newline \textbf{conditions)} & Các tham số mô phỏng đã được thiết lập và Đơn vị vận hành chọn\newline chức năngChạymôphỏng. \\
\hline
\textbf{Hậu điều kiện (Post-}\newline \textbf{conditions)} & Yêu cầu mô phỏng được chuyển đến bộ điều khiển; kết quả được\newline chuyển đếngiao diện sau khi xử lýhoàn tất. \\
\hline
\textbf{Luồng thuật toán}\newline \textbf{(Pseudo-code)} & 1. Tiếp nhận thao tác Chạy mô phỏng.\newline 2. Lấy tập tham số đã thiết lập.\newline 3. Gửi yêu cầu đến SimulationController.\newline 4. Chờ kết quả xử lý.\newline 5. Chuyển kếtquả đếnphươngthức hiển thị báo cáo. \\
\hline
\end{longtable}
\vspace{-0.5\baselineskip}
\subsubsection{showReport()}
\begin{longtable}{|L{0.24\textwidth}|L{0.70\textwidth}|}
\hline
\textbf{Thuộc tính} & \textbf{Nội dung chi tiết} \\
\hline
\endfirsthead
\hline
\textbf{Thuộc tính} & \textbf{Nội dung chi tiết} \\
\hline
\endhead
\textbf{Tên phương thức} & showReport() \\
\hline
\textbf{Mục đích} & Hiển thị báo cáo đánh giá tác động và các đề xuất điều phối được tạo\newline ra sau khi hoàn tất môphỏng. \\
\hline
\textbf{Tham số đầu vào}\newline \textbf{(@param)} & report: String - Nội dung báo cáo đánh giá tác động và đề xuất điều\newline phối. \\
\hline
\textbf{Giá trị trả về(@return)} & void - Khôngtrả vềgiá trị. \\
\hline
\textbf{Ngoại lệ (@exception)} & DisplayException - Xảyra khi khôngthể hiển thị báo cáo môphỏng. \\
\hline
\textbf{Tiền điều kiện (Pre-}\newline \textbf{conditions)} & Quá trình mô phỏng đã hoàn tất và hệ thống đã tạo được báo cáo kết\newline quả. \\
\hline
\textbf{Hậu điều kiện (Post-}\newline \textbf{conditions)} & Báo cáo đánh giá tác động và các đề xuất điều phối được hiển thị để\newline Đơn vị vận hành xem xét. \\
\hline
\textbf{Luồng thuật toán}\newline \textbf{(Pseudo-code)} & 1. Tiếp nhận nội dung báo cáo.\newline 2. Hiển thị kết quả đánh giá tác động.\newline 3. Hiển thị các đề xuất điều phối.\newline 4. Cho phép Đơn vị vận hành xem xét, kết thúc hoặc chọn áp dụng\newline đề xuất. \\
\hline
\end{longtable}
\vspace{-0.5\baselineskip}
\subsubsection{displayParameterError()}
\begin{longtable}{|L{0.24\textwidth}|L{0.70\textwidth}|}
\hline
\textbf{Thuộc tính} & \textbf{Nội dung chi tiết} \\
\hline
\endfirsthead
\hline
\textbf{Thuộc tính} & \textbf{Nội dung chi tiết} \\
\hline
\endhead
\textbf{Tên phương thức} & displayParameterError() \\
\hline
\textbf{Mục đích} & Hiển thị cảnh báo và đánh dấu các trường dữ liệu có tham số không\newline hợplệ để Đơn vị vận hành chỉnh sửa. \\
\hline
\textbf{Tham số đầu vào}\newline \textbf{(@param)} & errorMessage: String - Nội dung lỗi cần hiển thị. parameterName:\newline String- Tên trườngtham số khônghợplệ. \\
\hline
\textbf{Giá trị trả về(@return)} & void - Khôngtrả vềgiá trị. \\
\hline
\textbf{Ngoại lệ (@exception)} & DisplayException - Xảyra khi khôngthể hiển thị thôngbáo lỗi. \\
\hline
\textbf{Tiền điều kiện (Pre-}\newline \textbf{conditions)} & Hệ thống đã phát hiện tham số mô phỏng không hợp lệ. \\
\hline
\textbf{Hậu điều kiện (Post-}\newline \textbf{conditions)} & Trường tham số không hợp lệ được đánh dấu và thông báo lỗi được\newline hiển thị; Đơn vị vận hành có thể nhậplại dữ liệu. \\
\hline
\textbf{Luồng thuật toán}\newline \textbf{(Pseudo-code)} & 1. Tiếp nhận thông tin lỗi từ bộ điều khiển.\newline 2. Xác định trường tham số có lỗi.\newline 3. Đánh dấu trường dữ liệu tương ứng.\newline 4. Hiển thị thông báo cảnh báo.\newline 5. Yêu cầu Đơn vị vận hành nhậplại tham số. \\
\hline
\end{longtable}
\vspace{-0.5\baselineskip}
\subsubsection{applyRecommendation()}
\begin{longtable}{|L{0.24\textwidth}|L{0.70\textwidth}|}
\hline
\textbf{Thuộc tính} & \textbf{Nội dung chi tiết} \\
\hline
\endfirsthead
\hline
\textbf{Thuộc tính} & \textbf{Nội dung chi tiết} \\
\hline
\endhead
\textbf{Tên phương thức} & applyRecommendation() \\
\hline
\textbf{Mục đích} & Tiếp nhận thao tác của Đơn vị vận hành khi lựa chọn áp dụng một đề\newline xuất điềuphối từ báo cáo môphỏngvàoquytrình vận hành thực tế. \\
\hline
\end{longtable}
\vspace{-0.5\baselineskip}
\begin{longtable}{|L{0.24\textwidth}|L{0.70\textwidth}|}
\hline
\textbf{Thuộc tính} & \textbf{Nội dung chi tiết} \\
\hline
\endfirsthead
\hline
\textbf{Thuộc tính} & \textbf{Nội dung chi tiết} \\
\hline
\endhead
\textbf{Tham số đầu vào}\newline \textbf{(@param)} & targetUseCase: String - Use Case đích được lựa chọn, gồm UC11\newline hoặc UC13. parameters: Map<String, Object> - Các tham số tương\newline ứngvới đề xuất được lựa chọn. \\
\hline
\textbf{Giá trị trả về(@return)} & void - Khôngtrả vềgiá trị. \\
\hline
\textbf{Ngoại lệ (@exception)} & InvalidRecommendationException - Đề xuất được chọn không hợp\newline lệ hoặc khôngthuộc UC11, UC13. \\
\hline
\textbf{Tiền điều kiện (Pre-}\newline \textbf{conditions)} & Báo cáo mô phỏng đã được hiển thị và có đề xuất điều phối khả thi. \\
\hline
\textbf{Hậu điều kiện (Post-}\newline \textbf{conditions)} & Yêu cầu áp dụng đề xuất được gửi đến SimulationController để\newline chuyển tiếpsangUse Case tươngứng. \\
\hline
\textbf{Luồng thuật toán}\newline \textbf{(Pseudo-code)} & 1. Đơn vị vận hành chọn Áp dụng vào thực tế.\newline 2. Xác định Use Case đích UC11 hoặc UC13.\newline 3. Lấy các tham số của đề xuất được chọn.\newline 4. Gửiyêu cầu chuyển tiếpđến SimulationController. \\
\hline
\end{longtable}
\vspace{-0.5\baselineskip}
\subsubsection{initVirtualEnvironment()}
\begin{longtable}{|L{0.24\textwidth}|L{0.70\textwidth}|}
\hline
\textbf{Thuộc tính} & \textbf{Nội dung chi tiết} \\
\hline
\endfirsthead
\hline
\textbf{Thuộc tính} & \textbf{Nội dung chi tiết} \\
\hline
\endhead
\textbf{Tên phương thức} & initVirtualEnvironment() \\
\hline
\textbf{Mục đích} & Sao chép hoặc khởi tạo trạng thái hiện tại của toàn mạng lưới vào\newline môi trường mô phỏng ảo cách ly, bảo đảm quá trình mô phỏng\newline khôngảnh hưởngđến hệ thốngvận hành thực tế. \\
\hline
\textbf{Tham số đầu vào}\newline \textbf{(@param)} & Không có. \\
\hline
\textbf{Giá trị trả về (@return)} & Boolean - Trả về true nếu môi trường mô phỏng được khởi tạo thành\newline công, false nếu thất bại. \\
\hline
\textbf{Ngoại lệ (@exception)} & SimulationInitializationException - Xảy ra khi hệ thống không thể\newline cấp phát tài nguyên hoặc khởi tạo môi trườngảo cách ly. \\
\hline
\textbf{Tiền điều kiện (Pre-}\newline \textbf{conditions)} & Hệ thống có thể truy xuất trạng thái hiện tại của mạng lưới và có khả\newline năngcấp phát môi trườngmôphỏng. \\
\hline
\textbf{Hậu điều kiện (Post-}\newline \textbf{conditions)} & Trạng thái hiện tại của mạng lưới được khởi tạo trong môi trường\newline môphỏngriêngbiệt; dữ liệu vận hành thực tế khôngbị thayđổi. \\
\hline
\textbf{Luồng thuật toán}\newline \textbf{(Pseudo-code)} & 1. Tiếp nhận yêu cầu khởi tạo môi trường mô phỏng.\newline 2. Truy xuất trạng thái hiện tại của mạng lưới.\newline 3. Sao chép trạng thái vào môi trường mô phỏng riêng biệt.\newline 4. Kiểm tra khả năng hoạt động của môi trường.\newline 5. Trả về kếtquả khởi tạo. \\
\hline
\end{longtable}
\vspace{-0.5\baselineskip}
\subsubsection{validateParameters()}
\begin{longtable}{|L{0.24\textwidth}|L{0.70\textwidth}|}
\hline
\textbf{Thuộc tính} & \textbf{Nội dung chi tiết} \\
\hline
\endfirsthead
\hline
\textbf{Thuộc tính} & \textbf{Nội dung chi tiết} \\
\hline
\endhead
\textbf{Tên phương thức} & validateParameters() \\
\hline
\textbf{Mục đích} & Kiểm tra tính hợp lệ và khả năng thực hiện của các tham số giả định\newline trước khi chạymôphỏng. \\
\hline
\textbf{Tham số đầu vào} & parameters: Map<String, Object> - Tập hợp tham số giả định do \\
\hline
\textbf{(@param)} & Đơn vị vận hành cungcấp. \\
\hline
\end{longtable}
\vspace{-0.5\baselineskip}
\begin{longtable}{|L{0.24\textwidth}|L{0.70\textwidth}|}
\hline
\textbf{Thuộc tính} & \textbf{Nội dung chi tiết} \\
\hline
\endfirsthead
\hline
\textbf{Thuộc tính} & \textbf{Nội dung chi tiết} \\
\hline
\endhead
\textbf{Giá trị trả về (@return)} & Boolean - Trả về true nếu toàn bộ tham số hợp lệ, false nếu phát hiện\newline tham số khônghợplệ hoặc mâu thuẫn. \\
\hline
\textbf{Ngoại lệ (@exception)} & InvalidParameterException - Xảy ra khi tham số sai định dạng hoặc\newline khôngthể kiểm tra. \\
\hline
\textbf{Tiền điều kiện (Pre-}\newline \textbf{conditions)} & Các tham số mô phỏng đã được tiếp nhận từ giao diện. \\
\hline
\textbf{Hậu điều kiện (Post-}\newline \textbf{conditions)} & Kết quả kiểm tra tính hợp lệ của tham số được xác định; chỉ tham số\newline hợplệ mới được chuyển sangbước môphỏng. \\
\hline
\textbf{Luồng thuật toán}\newline \textbf{(Pseudo-code)} & 1. Tiếp nhận tập tham số.\newline 2. Kiểm tra định dạng và kiểu dữ liệu.\newline 3. Kiểm tra giá trị không hợp lệ, chẳng hạn số lượng xe âm.\newline 4. Kiểm tra giới hạn khả thi của môi trường mô phỏng.\newline 5. Trả về kếtquả kiểm tra. \\
\hline
\end{longtable}
\vspace{-0.5\baselineskip}
\subsubsection{calculateImpact()}
\begin{longtable}{|L{0.24\textwidth}|L{0.70\textwidth}|}
\hline
\textbf{Thuộc tính} & \textbf{Nội dung chi tiết} \\
\hline
\endfirsthead
\hline
\textbf{Thuộc tính} & \textbf{Nội dung chi tiết} \\
\hline
\endhead
\textbf{Tên phương thức} & calculateImpact() \\
\hline
\textbf{Mục đích} & Tính toán sự thay đổi trạng thái của mạng lưới dựa trên các tham số\newline giả định trong môi trường mô phỏng ảo cách ly nhằm đánh giá tác\newline độngcủa kịch bản. \\
\hline
\textbf{Tham số đầu vào}\newline \textbf{(@param)} & parameters: Map<String, Object> - Tập hợp tham số giả định được\newline sử dụngtrong quá trình môphỏng. \\
\hline
\textbf{Giá trị trả về (@return)} & String - Thông tin tổng hợp về kết quả thay đổi trạng thái và tác\newline độngcủa kịch bản môphỏng. \\
\hline
\textbf{Ngoại lệ (@exception)} & SimulationCalculationException - Xảy ra khi hệ thống không thể\newline thực hiện tính toán môphỏng. \\
\hline
\textbf{Tiền điều kiện (Pre-}\newline \textbf{conditions)} & Môi trường mô phỏng ảo đã được khởi tạo thành công và các tham\newline sốgiả định đã vượtqua bước kiểm tra tính hợplệ. \\
\hline
\textbf{Hậu điều kiện (Post-}\newline \textbf{conditions)} & Kết quả tính toán tác động được tạo ra trong môi trường mô phỏng;\newline trạngthái vận hành thực tế khôngbị thayđổi. \\
\hline
\textbf{Luồng thuật toán}\newline \textbf{(Pseudo-code)} & 1. Tiếp nhận các tham số giả định hợp lệ.\newline 2. Truy xuất trạng thái mạng lưới trong môi trường mô phỏng.\newline 3. Áp dụng các thay đổi giả định lên dữ liệu mô phỏng.\newline 4. Tính toán sự thay đổi về phương tiện, cổng sạc và tình trạng Hub.\newline 5. Tổng hợp kết quả đánh giá tác động.\newline 6. Trả về kếtquả tính toán. \\
\hline
\end{longtable}
\vspace{-0.5\baselineskip}
\subsubsection{generateReport()}
\begin{longtable}{|L{0.24\textwidth}|L{0.70\textwidth}|}
\hline
\textbf{Thuộc tính} & \textbf{Nội dung chi tiết} \\
\hline
\endfirsthead
\hline
\textbf{Thuộc tính} & \textbf{Nội dung chi tiết} \\
\hline
\endhead
\textbf{Tên phương thức} & generateReport() \\
\hline
\textbf{Mục đích} & Tổng hợp kết quả mô phỏng để tạo báo cáo đánh giá tác động và đề\newline xuấtphươngán điềuphốiphù hợpcho Đơn vị vận hành. \\
\hline
\textbf{Tham số đầu vào} & impactResult: String - Kết quả tính toán tác động từ phương thức \\
\hline
\textbf{(@param)} & calculateImpact(). \\
\hline
\end{longtable}
\vspace{-0.5\baselineskip}
\begin{longtable}{|L{0.24\textwidth}|L{0.70\textwidth}|}
\hline
\textbf{Thuộc tính} & \textbf{Nội dung chi tiết} \\
\hline
\endfirsthead
\hline
\textbf{Thuộc tính} & \textbf{Nội dung chi tiết} \\
\hline
\endhead
\textbf{Giá trị trả về (@return)} & String - Nội dung báo cáo mô phỏng bao gồm kết quả đánh giá tác\newline độngvà các đề xuất điềuphối. \\
\hline
\textbf{Ngoại lệ (@exception)} & ReportGenerationException - Xảy ra khi hệ thống không thể tổng\newline hợphoặc tạo báo cáo môphỏng. \\
\hline
\textbf{Tiền điều kiện (Pre-}\newline \textbf{conditions)} & Quá trình tính toán mô phỏng đã hoàn tất và có kết quả tác động hợp\newline lệ. \\
\hline
\textbf{Hậu điều kiện (Post-}\newline \textbf{conditions)} & Báo cáo đánh giá tác động và các đề xuất điều phối được tạo thành\newline công, sẵn sànghiển thị cho Đơn vị vận hành. \\
\hline
\textbf{Luồng thuật toán}\newline \textbf{(Pseudo-code)} & 1. Tiếp nhận kết quả tính toán tác động.\newline 2. Phân tích các thay đổi trạng thái trong kịch bản.\newline 3. Xác định vấn đề phát sinh và mức độ ảnh hưởng.\newline 4. Xây dựng đề xuất điều phối như điều chuyển phương tiện hoặc\newline thay đổi lịch sạc.\newline 5. Tổng hợp nội dung báo cáo.\newline 6. Trả về báo cáo môphỏng. \\
\hline
\end{longtable}
\vspace{-0.5\baselineskip}
\subsubsection{saveReport()}
\begin{longtable}{|L{0.24\textwidth}|L{0.70\textwidth}|}
\hline
\textbf{Thuộc tính} & \textbf{Nội dung chi tiết} \\
\hline
\endfirsthead
\hline
\textbf{Thuộc tính} & \textbf{Nội dung chi tiết} \\
\hline
\endhead
\textbf{Tên phương thức} & saveReport() \\
\hline
\textbf{Mục đích} & Lưu báo cáo mô phỏng sau khi Đơn vị vận hành chọn chức năng Kết\newline thúc, đồngthời liên kết báo cáo với kịch bản môphỏngtươngứng. \\
\hline
\textbf{Tham số đầu vào}\newline \textbf{(@param)} & scenario: SimulationScenario - Kịch bản mô phỏng cần lưu, bao gồm\newline mã kịch bản và nội dungtóm tắt báo cáo. \\
\hline
\textbf{Giá trị trả về (@return)} & Boolean - Trả về true nếu báo cáo được lưu thành công, false nếu\newline lưu thất bại. \\
\hline
\textbf{Ngoại lệ (@exception)} & ReportSaveException - Xảy ra khi hệ thống không thể lưu báo cáo\newline môphỏng. \\
\hline
\textbf{Tiền điều kiện (Pre-}\newline \textbf{conditions)} & Quá trình mô phỏng đã hoàn tất, báo cáo đã được tạo và Đơn vị vận\newline hành chọn chức năngKết thúc. \\
\hline
\textbf{Hậu điều kiện (Post-}\newline \textbf{conditions)} & Báo cáo mô phỏng được lưu và gắn với SimulationScenario; trạng\newline thái vận hành thực tế của mạnglưới khôngbị thayđổi. \\
\hline
\textbf{Luồng thuật toán}\newline \textbf{(Pseudo-code)} & 1. Tiếp nhận yêu cầu lưu báo cáo.\newline 2. Tiếp nhận thông tin kịch bản và nội dung báo cáo.\newline 3. Cập nhật nội dung tóm tắt báo cáo vào SimulationScenario.\newline 4. Lưu thông tin kịch bản và báo cáo.\newline 5. Trả về kếtquả lưu. \\
\hline
\end{longtable}
\vspace{-0.5\baselineskip}
\subsubsection{redirectToUseCase()}
\begin{longtable}{|L{0.24\textwidth}|L{0.70\textwidth}|}
\hline
\textbf{Thuộc tính} & \textbf{Nội dung chi tiết} \\
\hline
\endfirsthead
\hline
\textbf{Thuộc tính} & \textbf{Nội dung chi tiết} \\
\hline
\endhead
\textbf{Tên phương thức} & redirectToUseCase() \\
\hline
\textbf{Mục đích} & Chuyển Đơn vị vận hành đến màn hình UC11 hoặc UC13 tương ứng\newline với đề xuất được chọn, đồng thời điền sẵn các tham số liên quan để\newline tiếptục xử lýnghiệpvụ. \\
\hline
\textbf{Tham số đầu vào}\newline \textbf{(@param)} & targetUseCase: String - Mã Use Case đích, gồm UC11 hoặc UC13.\newline parameters: Map<String, Object> - Các tham số tương ứng với đề\newline xuất điềuphối được lựa chọn. \\
\hline
\end{longtable}
\vspace{-0.5\baselineskip}
\begin{longtable}{|L{0.24\textwidth}|L{0.70\textwidth}|}
\hline
\textbf{Thuộc tính} & \textbf{Nội dung chi tiết} \\
\hline
\endfirsthead
\hline
\textbf{Thuộc tính} & \textbf{Nội dung chi tiết} \\
\hline
\endhead
\textbf{Giá trị trả về (@return)} & Boolean - Trả về true nếu chuyển tiếp thành công, false nếu không\newline thể chuyển tiếp. \\
\hline
\textbf{Ngoại lệ (@exception)} & InvalidTargetUseCaseException - Use Case đích không được hỗ trợ\newline hoặc tham số chuyển tiếpkhônghợplệ. \\
\hline
\textbf{Tiền điều kiện (Pre-}\newline \textbf{conditions)} & Báo cáo mô phỏng đã được hiển thị và Đơn vị vận hành đã lựa chọn\newline một đề xuất khả thi để ápdụng. \\
\hline
\textbf{Hậu điều kiện (Post-}\newline \textbf{conditions)} & Đơn vị vận hành được chuyển đến UC11 hoặc UC13 tương ứng và\newline các tham số đề xuất được điền sẵn; trạng thái vận hành thực tế chỉ\newline thay đổi khi nghiệp vụ tại Use Case đích được thực hiện và xác\newline nhận. \\
\hline
\textbf{Luồng thuật toán}\newline \textbf{(Pseudo-code)} & 1. Tiếp nhận Use Case đích và tham số đề xuất.\newline 2. Kiểm tra Use Case đích có phải UC11 hoặc UC13.\newline 3. Xác định màn hình tương ứng.\newline 4. Truyền các tham số đề xuất sang Use Case đích.\newline 5. Chuyển hướng giao diện đến màn hình tương ứng.\newline 6. Trả về kếtquả chuyển tiếp. \\
\hline
\end{longtable}
\vspace{-0.5\baselineskip}
\subsection{Use Case UC15 - Tùy chỉnh giao diện (Customize UI)}
\subsubsection{toggleTheme()}
\begin{longtable}{|L{0.24\textwidth}|L{0.70\textwidth}|}
\hline
\textbf{Thuộc tính} & \textbf{Nội dung chi tiết} \\
\hline
\endfirsthead
\hline
\textbf{Thuộc tính} & \textbf{Nội dung chi tiết} \\
\hline
\endhead
\textbf{Tên phương thức} & toggleTheme() \\
\hline
\textbf{Mục đích} & Hiển thị danh sách tùy chọn giao diện để người dùng lựa chọn chế\newline độ sáng, tối hoặc theo hệ thống. \\
\hline
\textbf{Tham số đầu vào}\newline \textbf{(@param)} & Không có. \\
\hline
\textbf{Giá trị trả về(@return)} & void – Khôngtrả vềgiá trị. \\
\hline
\textbf{Ngoại lệ (@exception)} & Khôngcó ngoại lệ nghiệpvụ được xác định. \\
\hline
\textbf{Tiền điều kiện (Pre-}\newline \textbf{conditions)} & Người dùng đang truy cập giao diện hệ thống. \\
\hline
\textbf{Hậu điều kiện (Post-}\newline \textbf{conditions)} & Danh sách tùy chọn giao diện được hiển thị. Chế độ giao diện hiện\newline tại chưa thayđổi. \\
\hline
\textbf{Luồng thuật toán}\newline \textbf{(Pseudo-code)} & 1. Người dùng nhấn biểu tượng giao diện.\newline 2. Hệ thống hiển thị danh sách Light, Dark và System.\newline 3. Chờ người dùnglựa chọn chế độ. \\
\hline
\end{longtable}
\vspace{-0.5\baselineskip}
\subsubsection{switchLanguage()}
\begin{longtable}{|L{0.24\textwidth}|L{0.70\textwidth}|}
\hline
\textbf{Thuộc tính} & \textbf{Nội dung chi tiết} \\
\hline
\endfirsthead
\hline
\textbf{Thuộc tính} & \textbf{Nội dung chi tiết} \\
\hline
\endhead
\textbf{Tên phương thức} & switchLanguage() \\
\hline
\textbf{Mục đích} & Hiển thị danh sách ngôn ngữ để người dùng lựa chọn ngôn ngữ hiển\newline thị của hệ thống. \\
\hline
\textbf{Tham số đầu vào}\newline \textbf{(@param)} & Không có. \\
\hline
\textbf{Giá trị trả về(@return)} & void – Khôngtrả vềgiá trị. \\
\hline
\textbf{Ngoại lệ (@exception)} & Khôngcó ngoại lệ nghiệpvụ được xác định. \\
\hline
\end{longtable}
\vspace{-0.5\baselineskip}
\begin{longtable}{|L{0.24\textwidth}|L{0.70\textwidth}|}
\hline
\textbf{Thuộc tính} & \textbf{Nội dung chi tiết} \\
\hline
\endfirsthead
\hline
\textbf{Thuộc tính} & \textbf{Nội dung chi tiết} \\
\hline
\endhead
\textbf{Tiền điều kiện (Pre-}\newline \textbf{conditions)} & Người dùng đang truy cập giao diện hệ thống. \\
\hline
\textbf{Hậu điều kiện (Post-}\newline \textbf{conditions)} & Danh sách ngôn ngữ được hiển thị, bao gồm Tiếng Việt và Tiếng\newline Anh. \\
\hline
\textbf{Luồng thuật toán}\newline \textbf{(Pseudo-code)} & 1. Người dùng nhấn biểu tượng ngôn ngữ.\newline 2. Hệ thống hiển thị danh sách ngôn ngữ hỗ trợ.\newline 3. Chờ người dùnglựa chọn ngôn ngữ. \\
\hline
\end{longtable}
\vspace{-0.5\baselineskip}
\subsubsection{displayError()}
\begin{longtable}{|L{0.24\textwidth}|L{0.70\textwidth}|}
\hline
\textbf{Thuộc tính} & \textbf{Nội dung chi tiết} \\
\hline
\endfirsthead
\hline
\textbf{Thuộc tính} & \textbf{Nội dung chi tiết} \\
\hline
\endhead
\textbf{Tên phương thức} & displayError() \\
\hline
\textbf{Mục đích} & Hiển thị thông báo khi hệ thống không thể lưu cấu hình giao diện\newline hoặc ngôn ngữ. \\
\hline
\textbf{Tham số đầu vào}\newline \textbf{(@param)} & errorMessage: String – Nội dung thông báo lỗi. \\
\hline
\textbf{Giá trị trả về(@return)} & void – Khôngtrả vềgiá trị. \\
\hline
\textbf{Ngoại lệ (@exception)} & Khôngcó ngoại lệ nghiệpvụ được xác định. \\
\hline
\textbf{Tiền điều kiện (Pre-}\newline \textbf{conditions)} & Hệ thống phát sinh lỗi trong quá trình lưu cấu hình. \\
\hline
\textbf{Hậu điều kiện (Post-}\newline \textbf{conditions)} & Thông báo lỗi được hiển thị cho người dùng; dữ liệu biểu mẫu đang\newline nhậpkhôngbị mất. \\
\hline
\textbf{Luồng thuật toán}\newline \textbf{(Pseudo-code)} & 1. Tiếp nhận nội dung lỗi.\newline 2. Hiển thị thông báo lỗi trên giao diện.\newline 3. Giữ nguyên dữ liệu biểu mẫu đangthực hiện. \\
\hline
\end{longtable}
\vspace{-0.5\baselineskip}
\subsubsection{applyTheme()}
\begin{longtable}{|L{0.24\textwidth}|L{0.70\textwidth}|}
\hline
\textbf{Thuộc tính} & \textbf{Nội dung chi tiết} \\
\hline
\endfirsthead
\hline
\textbf{Thuộc tính} & \textbf{Nội dung chi tiết} \\
\hline
\endhead
\textbf{Tên phương thức} & applyTheme() \\
\hline
\textbf{Mục đích} & Áp dụng chế độ giao diện được người dùng lựa chọn cho toàn bộ hệ\newline thống. \\
\hline
\textbf{Tham số đầu vào}\newline \textbf{(@param)} & theme: String – Chế độ giao diện được chọn (LIGHT, DARK,\newline SYSTEM). \\
\hline
\textbf{Giá trị trả về(@return)} & Boolean – Trả về true nếu áp dụngthành công, ngược lại trả về false. \\
\hline
\textbf{Ngoại lệ (@exception)} & ThemeApplicationException – Xảy ra khi không thể áp dụng chế độ\newline giao diện. \\
\hline
\textbf{Tiền điều kiện (Pre-}\newline \textbf{conditions)} & Người dùng đã lựa chọn chế độ giao diện hợp lệ. \\
\hline
\textbf{Hậu điều kiện (Post-}\newline \textbf{conditions)} & Giao diện được cập nhật theo chế độ đã chọn mà không tải lại trang\newline hoặc làm mất dữ liệu đangnhập. \\
\hline
\textbf{Luồng thuật toán}\newline \textbf{(Pseudo-code)} & 1. Tiếp nhận chế độ giao diện.\newline 2. Kiểm tra chế độ được lựa chọn.\newline 3. Áp dụng màu sắc và kiểu hiển thị tương ứng.\newline 4. Trả về kếtquả thực hiện. \\
\hline
\end{longtable}
\vspace{-0.5\baselineskip}
\subsubsection{applyLanguage()}
\begin{longtable}{|L{0.24\textwidth}|L{0.70\textwidth}|}
\hline
\textbf{Thuộc tính} & \textbf{Nội dung chi tiết} \\
\hline
\endfirsthead
\hline
\textbf{Thuộc tính} & \textbf{Nội dung chi tiết} \\
\hline
\endhead
\textbf{Tên phương thức} & applyLanguage() \\
\hline
\textbf{Mục đích} & Thay đổi ngôn ngữ hiển thị của hệ thống theo lựa chọn của người\newline dùng. \\
\hline
\textbf{Tham số đầu vào}\newline \textbf{(@param)} & language: String – Ngôn ngữ được chọn (VI, EN). \\
\hline
\textbf{Giá trị trả về (@return)} & Boolean – Trả về true nếu thay đổi thành công, ngược lại trả về\newline false. \\
\hline
\textbf{Ngoại lệ (@exception)} & LanguageApplicationException – Xảy ra khi không thể áp dụng\newline ngôn ngữ được chọn. \\
\hline
\textbf{Tiền điều kiện (Pre-}\newline \textbf{conditions)} & Người dùng đã lựa chọn ngôn ngữ được hệ thống hỗ trợ. \\
\hline
\textbf{Hậu điều kiện (Post-}\newline \textbf{conditions)} & Nội dung giao diện được chuyển sang ngôn ngữ đã chọn, không tải\newline lại trangvà khônglàm mất dữ liệu biểu mẫu. \\
\hline
\textbf{Luồng thuật toán}\newline \textbf{(Pseudo-code)} & 1. Tiếp nhận ngôn ngữ được chọn.\newline 2. Kiểm tra ngôn ngữ có được hỗ trợ hay không.\newline 3. Tải nội dung tương ứng với ngôn ngữ đã chọn.\newline 4. Cập nhật nội dung giao diện.\newline 5. Trả về kếtquả thực hiện. \\
\hline
\end{longtable}
\vspace{-0.5\baselineskip}
\subsubsection{preserveFormState()}
\begin{longtable}{|L{0.24\textwidth}|L{0.70\textwidth}|}
\hline
\textbf{Thuộc tính} & \textbf{Nội dung chi tiết} \\
\hline
\endfirsthead
\hline
\textbf{Thuộc tính} & \textbf{Nội dung chi tiết} \\
\hline
\endhead
\textbf{Tên phương thức} & preserveFormState() \\
\hline
\textbf{Mục đích} & Bảo toàn dữ liệu người dùng đang nhập trên biểu mẫu khi thay đổi\newline giao diện hoặc ngôn ngữ. \\
\hline
\textbf{Tham số đầu vào}\newline \textbf{(@param)} & formState: Map<String, Object> – Dữ liệu hiện tại của biểu mẫu. \\
\hline
\textbf{Giá trị trả về(@return)} & Map<String, Object> – Trả về dữ liệu biểu mẫu đã được bảo toàn. \\
\hline
\textbf{Ngoại lệ (@exception)} & FormStateException – Xảy ra khi không thể bảo toàn dữ liệu biểu\newline mẫu. \\
\hline
\textbf{Tiền điều kiện (Pre-}\newline \textbf{conditions)} & Người dùng đang thực hiện một biểu mẫu có dữ liệu chưa được gửi\newline đi. \\
\hline
\textbf{Hậu điều kiện (Post-}\newline \textbf{conditions)} & Dữ liệu biểu mẫu, bao gồm Hub đã chọn, được giữ nguyên sau khi\newline thay đổigiao diện hoặc ngôn ngữ. \\
\hline
\textbf{Luồng thuật toán}\newline \textbf{(Pseudo-code)} & 1. Tiếp nhận dữ liệu biểu mẫu hiện tại.\newline 2. Ghi nhận các giá trị người dùng đã nhập và lựa chọn.\newline 3. Bảo toàn dữ liệu trong quá trình cập nhật giao diện.\newline 4. Trả về trạngthái biểu mẫu đã được bảo toàn. \\
\hline
\end{longtable}
\vspace{-0.5\baselineskip}
\subsubsection{savePreference()}
\begin{longtable}{|L{0.24\textwidth}|L{0.70\textwidth}|}
\hline
\textbf{Thuộc tính} & \textbf{Nội dung chi tiết} \\
\hline
\endfirsthead
\hline
\textbf{Thuộc tính} & \textbf{Nội dung chi tiết} \\
\hline
\endhead
\textbf{Tên phương thức} & savePreference() \\
\hline
\textbf{Mục đích} & Lưu tùy chọn giao diện và ngôn ngữ của người dùng để duy trì thiết\newline lậpđã lựa chọn. \\
\hline
\end{longtable}
\vspace{-0.5\baselineskip}
\begin{longtable}{|L{0.24\textwidth}|L{0.70\textwidth}|}
\hline
\textbf{Thuộc tính} & \textbf{Nội dung chi tiết} \\
\hline
\endfirsthead
\hline
\textbf{Thuộc tính} & \textbf{Nội dung chi tiết} \\
\hline
\endhead
\textbf{Tham số đầu vào}\newline \textbf{(@param)} & theme: String – Chế độ giao diện; language: String – Ngôn ngữ được\newline lựa chọn. \\
\hline
\textbf{Giá trị trả về(@return)} & Boolean – Trả về true nếu lưu thành công, ngược lại trả về false. \\
\hline
\textbf{Ngoại lệ (@exception)} & PreferenceSaveException – Xảyra khi khôngthể lưu cấu hình. \\
\hline
\textbf{Tiền điều kiện (Pre-}\newline \textbf{conditions)} & Chế độ giao diện và ngôn ngữ đã được lựa chọn hợp lệ. \\
\hline
\textbf{Hậu điều kiện (Post-}\newline \textbf{conditions)} & Tùy chọn giao diện và ngôn ngữ được lưu thành công vào bộ nhớ\newline cấu hình cục bộ. \\
\hline
\textbf{Luồng thuật toán}\newline \textbf{(Pseudo-code)} & 1. Tiếp nhận chế độ giao diện và ngôn ngữ.\newline 2. Kiểm tra tính hợp lệ của các tùy chọn.\newline 3. Lưu cấu hình vào bộ nhớ cục bộ.\newline 4. Trả về kếtquả lưu cấu hình. \\
\hline
\end{longtable}
\vspace{-0.5\baselineskip}
\subsubsection{handleSaveError()}
\begin{longtable}{|L{0.24\textwidth}|L{0.70\textwidth}|}
\hline
\textbf{Thuộc tính} & \textbf{Nội dung chi tiết} \\
\hline
\endfirsthead
\hline
\textbf{Thuộc tính} & \textbf{Nội dung chi tiết} \\
\hline
\endhead
\textbf{Tên phương thức} & handleSaveError() \\
\hline
\textbf{Mục đích} & Xử lý lỗi phát sinh khi hệ thống không thể lưu tùy chọn giao diện\newline hoặc ngôn ngữ. \\
\hline
\textbf{Tham số đầu vào}\newline \textbf{(@param)} & error: Exception – Thông tin ngoại lệ phát sinh trong quá trình lưu\newline cấu hình. \\
\hline
\textbf{Giá trị trả về(@return)} & String– Nội dungthôngbáo lỗi cần hiển thị. \\
\hline
\textbf{Ngoại lệ (@exception)} & Khôngcó ngoại lệ nghiệpvụ bổ sungđược xác định. \\
\hline
\textbf{Tiền điều kiện (Pre-}\newline \textbf{conditions)} & Phương thức savePreference() thực hiện không thành công. \\
\hline
\textbf{Hậu điều kiện (Post-}\newline \textbf{conditions)} & Lỗi được xử lý và thông tin lỗi được chuyển đến giao diện; hệ thống\newline duytrì cấu hình hiện tại và khônglàm mất dữ liệu biểu mẫu. \\
\hline
\textbf{Luồng thuật toán}\newline \textbf{(Pseudo-code)} & 1. Tiếp nhận thông tin ngoại lệ.\newline 2. Xác định lỗi phát sinh trong quá trình lưu.\newline 3. Tạo nội dung thông báo lỗi.\newline 4. Trả về thôngbáo đểgiao diện hiển thị. \\
\hline
\end{longtable}
\vspace{-0.5\baselineskip}
\subsection{Use Case UC16 – Cập nhật trạng thái tài nguyên (Update Resource Status)}
\subsubsection{receiveTelemetryPacket(packet)}
\begin{longtable}{|L{0.24\textwidth}|L{0.70\textwidth}|}
\hline
\textbf{Thuộc tính} & \textbf{Nội dung chi tiết} \\
\hline
\endfirsthead
\hline
\textbf{Thuộc tính} & \textbf{Nội dung chi tiết} \\
\hline
\endhead
\textbf{Tên phương thức} & receiveTelemetryPacket(packet) \\
\hline
\textbf{Mục đích} & Tiếp nhận gói tin trạng thái thời gian thực từ Bộ giả lập dữ liệu (hoặc\newline cảm biến IoT), điều phối qua ResourceSyncController và trả mã\newline phản hồi HTTP chuẩn(200 OK hoặc 400 Bad Request). \\
\hline
\textbf{Tham số đầu vào}\newline \textbf{(@param)} & packet: ResourcePacket – Gói dữ liệu trạng thái gồm ID tài nguyên,\newline trạngthái, mứcpin, tọa độ, timestampvà cờ môphỏng. \\
\hline
\textbf{Giá trị trả về (@return)} & HttpResponse – Mã 200 OK nếu cập nhật thành công hoặc loại bỏ\newline gói trễ đúng quy định; mã 400 Bad Request nếu sai định dạng hoặc\newline khôngtìm thấyID(EX2). \\
\hline
\end{longtable}
\vspace{-0.5\baselineskip}
\begin{longtable}{|L{0.24\textwidth}|L{0.70\textwidth}|}
\hline
\textbf{Thuộc tính} & \textbf{Nội dung chi tiết} \\
\hline
\endfirsthead
\hline
\textbf{Thuộc tính} & \textbf{Nội dung chi tiết} \\
\hline
\endhead
\textbf{Ngoại lệ (@exception)} & Không ném ngoại lệ ra ngoài giao thức HTTP; các ngoại lệ nghiệp\newline vụ được ánh xạ thành mã trạngthái HTTP tươngứng. \\
\hline
\textbf{Tiền điều kiện (Pre-}\newline \textbf{conditions)} & Hệ thống đang hoạt động và lắng nghe kết nối từ Bộ giả lập dữ liệu. \\
\hline
\textbf{Hậu điều kiện (Post-}\newline \textbf{conditions)} & Gói tin được xử lý; nếu hợp lệ và mới nhất, giao diện người dùng\newline nhận cập nhật trong <= REALTIME\_SLA\_SECONDS (15 giây –\newline NFR02). \\
\hline
\textbf{Luồng thuật toán}\newline \textbf{(Pseudo-code)} & 1. Try:\newline a. Gọi result =\newline ResourceSyncController.ingestResourceStatus(packet).\newline b. Trả về HttpResponse(200, "OK", result).\newline 2. Catch (InvalidPacketException \textbar\{\} ResourceNotFoundException\newline ex):\newline a. Trả về HttpResponse(400, "Bad Request", ex.message). \\
\hline
\end{longtable}
\vspace{-0.5\baselineskip}
\subsubsection{ingestResourceStatus(packet)}
\begin{longtable}{|L{0.24\textwidth}|L{0.70\textwidth}|}
\hline
\textbf{Thuộc tính} & \textbf{Nội dung chi tiết} \\
\hline
\endfirsthead
\hline
\textbf{Thuộc tính} & \textbf{Nội dung chi tiết} \\
\hline
\endhead
\textbf{Tên phương thức} & ingestResourceStatus(packet) \\
\hline
\textbf{Mục đích} & Xử lý toàn bộ luồng đồng bộ trạng thái tài nguyên: kiểm tra hợp lệ\newline cấu trúc và ID (EX2), lọc/loại bỏ gói tin trễ theo Timestamp (EX1 –\newline NFR08), tách luồng What-If Sandbox (AF1 – NFR07) và phát tín\newline hiệu WebSocket/SSE(NFR02). \\
\hline
\textbf{Tham số đầu vào}\newline \textbf{(@param)} & packet: ResourcePacket – Gói tin telemetry nhận được từ Endpoint. \\
\hline
\textbf{Giá trị trả về (@return)} & SyncResult – Kết quả xử lý gồm trạng thái\newline \{UPDATED\_OPERATIONAL, UPDATED\_SANDBOX,\newline DROPPED\_OUTDATED\}. \\
\hline
\textbf{Ngoại lệ (@exception)} & InvalidPacketException (EX2) – Gói tin thiếu trường bắt buộc hoặc\newline giá trị không hợp lệ.\newline ResourceNotFoundException (EX2) – Không tìm thấy tài nguyên\newline tươngứngvớipacket.resourceId. \\
\hline
\textbf{Tiền điều kiện (Pre-}\newline \textbf{conditions)} & Gói tin đã được tiếp nhận tại ResourceIngestionEndpoint. \\
\hline
\textbf{Hậu điều kiện (Post-}\newline \textbf{conditions)} & Nếu gói mới hợp lệ (vận hành thực tế): Operational DB được cập\newline nhật và RealtimeGateway phát tín hiệu. Nếu isSimulation=true: chỉ\newline ghi vào SandboxStateStore. Nếu gói trễ: hủy gói, ghi log loại bỏ,\newline không ghi đè CSDL. \\
\hline
\end{longtable}
\vspace{-0.5\baselineskip}
\begin{longtable}{|L{0.24\textwidth}|L{0.70\textwidth}|}
\hline
\textbf{Thuộc tính} & \textbf{Nội dung chi tiết} \\
\hline
\endfirsthead
\hline
\textbf{Thuộc tính} & \textbf{Nội dung chi tiết} \\
\hline
\endhead
\textbf{Luồng thuật toán}\newline \textbf{(Pseudo-code)} & 1. Gọi validatePacket(packet); nếu lỗi -> ghi AuditLog("SYSTEM",\newline "INTEGRATION\_ERROR", packet.resourceId, err) và ném ngoại\newline lệ.\newline 2. Nếu packet.isSimulation == true (AF1 – NFR07):\newline a. Gọi SandboxStateStore.upsertSimulatedState(packet).\newline b. Trả về SyncResult(UPDATED\_SANDBOX) // Tuyệt đối không\newline ghi đè Operational DB.\newline 3. Tìm Resource resource =\newline resourceRepo.findById(packet.resourceId); nếu null -> ghi AuditLog\newline lỗi tích hợp và ném ResourceNotFoundException.\newline 4. Nếu isOutdatedPacket(packet.packetTimestamp,\newline resource.lastUpdatedAt) == true (EX1 – NFR08):\newline a. Gọi AuditLog.record("SIMULATOR",\newline "DROP\_OUTDATED\_PACKET", packet.resourceId, "Dropped\newline stale packet").\newline b. Trả về SyncResult(DROPPED\_OUTDATED) // Không ghi đè\newline dữ liệu mới hơn.\newline 5. Cập nhật Operational DB:\newline - Nếu resource là SharedVehicle -> gọi\newline resource.updateTelemetry(packet.status, packet.batteryLevel,\newline packet.packetTimestamp).\newline - Ngược lại -> gọi resource.updateStatus(packet.status,\newline packet.packetTimestamp).\newline 6. Gọi RealtimeGateway.broadcastResourceUpdate(resource.id,\newline resource.status, packet.packetTimestamp) (NFR02 <= 15s).\newline 7. Trả về SyncResult(UPDATED\_OPERATIONAL). \\
\hline
\end{longtable}
\vspace{-0.5\baselineskip}
\subsubsection{validatePacket(packet)}
\begin{longtable}{|L{0.24\textwidth}|L{0.70\textwidth}|}
\hline
\textbf{Thuộc tính} & \textbf{Nội dung chi tiết} \\
\hline
\endfirsthead
\hline
\textbf{Thuộc tính} & \textbf{Nội dung chi tiết} \\
\hline
\endhead
\textbf{Tên phương thức} & validatePacket(packet) \\
\hline
\textbf{Mục đích} & Kiểm tra tính đầy đủ và hợp lệ về cấu trúc dữ liệu của gói tin\newline telemetry trước khi truy xuất hoặc ghi CSDL (UC16 – Bước 2 \&\newline EX2). \\
\hline
\textbf{Tham số đầu vào}\newline \textbf{(@param)} & packet: ResourcePacket – Gói tin cần kiểm tra. \\
\hline
\textbf{Giá trị trả về(@return)} & void – Kết thúc bình thườngnếugói tin hợp lệ. \\
\hline
\textbf{Ngoại lệ (@exception)} & InvalidPacketException – Xảy ra khi packet null, thiếu\newline resourceId/status/packetTimestamp, hoặc batteryLevel nằm ngoài [0,\newline 100]. \\
\hline
\textbf{Tiền điều kiện (Pre-}\newline \textbf{conditions)} & Gói tin đã được deserialize thành đối tượng ResourcePacket. \\
\hline
\textbf{Hậu điều kiện (Post-}\newline \textbf{conditions)} & Đảm bảo gói tin đạt chuẩn cấu trúc trước khi so sánh timestamp\newline hoặc cậpnhật trạngthái. \\
\hline
\end{longtable}
\vspace{-0.5\baselineskip}
\begin{longtable}{|L{0.24\textwidth}|L{0.70\textwidth}|}
\hline
\textbf{Thuộc tính} & \textbf{Nội dung chi tiết} \\
\hline
\endfirsthead
\hline
\textbf{Thuộc tính} & \textbf{Nội dung chi tiết} \\
\hline
\endhead
\textbf{Luồng thuật toán} & 1. Nếu packet == null -> ném InvalidPacketException("Packet is\newline null").\newline 2. Nếu packet.resourceId rỗng hoặc packet.status == null hoặc\newline packet.packetTimestamp == null -> ném \\
\hline
\textbf{(Pseudo-code)} & InvalidPacketException("Missing required fields").\newline 3. Nếu packet.batteryLevel != null và (packet.batteryLevel < 0 \textbar\{\}\textbar\{\}\newline packet.batteryLevel > 100) -> ném InvalidPacketException("Invalid\newline batterylevel"). \\
\hline
\end{longtable}
\vspace{-0.5\baselineskip}
\subsubsection{isOutdatedPacket(packetTimestamp, currentTimestamp)}
\begin{longtable}{|L{0.24\textwidth}|L{0.70\textwidth}|}
\hline
\textbf{Thuộc tính} & \textbf{Nội dung chi tiết} \\
\hline
\endfirsthead
\hline
\textbf{Thuộc tính} & \textbf{Nội dung chi tiết} \\
\hline
\endhead
\textbf{Tên phương thức} & isOutdatedPacket(packetTimestamp, currentTimestamp) \\
\hline
\textbf{Mục đích} & So sánh dấu thời gian của gói tin đến với thời điểm cập nhật hiện tại\newline trongCSDL đểphát hiệngói tin bị trễ hoặc trùnglặp (NFR08). \\
\hline
\textbf{Tham số đầu vào}\newline \textbf{(@param)} & packetTimestamp: LocalDateTime – Thời điểm ghi nhận trên gói tin\newline đến.\newline currentTimestamp: LocalDateTime (nullable) – Thời điểm cập nhật\newline hiện tại của tài nguyên trongCSDL. \\
\hline
\textbf{Giá trị trả về (@return)} & boolean – Trả về true nếu gói tin bị lỗi thời (packetTimestamp <=\newline currentTimestamp), false nếu gói tin mới hơn (packetTimestamp >\newline currentTimestamp). \\
\hline
\textbf{Ngoại lệ (@exception)} & Khôngcó. \\
\hline
\textbf{Tiền điều kiện (Pre-}\newline \textbf{conditions)} & packetTimestamp đã được kiểm tra khác null trong validatePacket(). \\
\hline
\textbf{Hậu điều kiện (Post-}\newline \textbf{conditions)} & Không thay đổi dữ liệu; kết quả quyết định cho phép cập nhật hay\newline hủybỏgói tin. \\
\hline
\textbf{Luồng thuật toán}\newline \textbf{(Pseudo-code)} & 1. Nếu currentTimestamp == null -> trả về false (chấp nhận bản ghi\newline đầu tiên).\newline 2. Trả vềpacketTimestamp<= currentTimestamp. \\
\hline
\end{longtable}
\vspace{-0.5\baselineskip}
\subsubsection{upsertSimulatedState(packet)}
\begin{longtable}{|L{0.24\textwidth}|L{0.70\textwidth}|}
\hline
\textbf{Thuộc tính} & \textbf{Nội dung chi tiết} \\
\hline
\endfirsthead
\hline
\textbf{Thuộc tính} & \textbf{Nội dung chi tiết} \\
\hline
\endhead
\textbf{Tên phương thức} & upsertSimulatedState(packet) \\
\hline
\textbf{Mục đích} & Cập nhật trạng thái tài nguyên giả lập vào kho lưu trữ ảo cách ly\newline (What-If Sandbox) mà hoàn toàn không tác động tới Operational DB\newline (UC16 – AF1, NFR07). \\
\hline
\textbf{Tham số đầu vào}\newline \textbf{(@param)} & packet: ResourcePacket – Gói tin mô phỏng có cờ isSimulation =\newline true. \\
\hline
\textbf{Giá trị trả về(@return)} & void – Khôngtrả vềgiá trị. \\
\hline
\textbf{Ngoại lệ (@exception)} & InvalidStateException – Xảyra khipacket.isSimulation == false. \\
\hline
\textbf{Tiền điều kiện (Pre-}\newline \textbf{conditions)} & Môi trường What-If Sandbox đã được khởi tạo theo scenarioId. \\
\hline
\textbf{Hậu điều kiện (Post-}\newline \textbf{conditions)} & Bản đồ simulatedResources trong Sandbox được cập nhật trạng thái\newline mới nhất (với kiểm tra timestamp nội bộ Sandbox); Operational DB\newline khôngthay đổi. \\
\hline
\end{longtable}
\vspace{-0.5\baselineskip}
\begin{longtable}{|L{0.24\textwidth}|L{0.70\textwidth}|}
\hline
\textbf{Thuộc tính} & \textbf{Nội dung chi tiết} \\
\hline
\endfirsthead
\hline
\textbf{Thuộc tính} & \textbf{Nội dung chi tiết} \\
\hline
\endhead
\textbf{Luồng thuật toán}\newline \textbf{(Pseudo-code)} & 1. Nếu !packet.isSimulation -> ném\newline InvalidStateException("Operational packet cannot enter sandbox").\newline 2. Lấy simulated = this.simulatedResources.get(packet.resourceId).\newline 3. Nếu simulated == null hoặc packet.packetTimestamp >\newline simulated.lastUpdatedAt:\newline a. Cập nhật trạng thái và lastUpdatedAt của bản sao trong\newline simulatedResources.\newline b. Cậpnhật this.lastSyncAt = now(). \\
\hline
\end{longtable}
\vspace{-0.5\baselineskip}
\subsection{Các lớp Thực thể Lõi \& Giao diện Dùng chung (Core Domain \& Interfaces)}
\subsubsection{MobilityHub.hasAvailableCapacity(requiredCount)}
\begin{longtable}{|L{0.24\textwidth}|L{0.70\textwidth}|}
\hline
\textbf{Thuộc tính} & \textbf{Nội dung chi tiết} \\
\hline
\endfirsthead
\hline
\textbf{Thuộc tính} & \textbf{Nội dung chi tiết} \\
\hline
\endhead
\textbf{Tên phương thức} & MobilityHub.hasAvailableCapacity(requiredCount) \\
\hline
\textbf{Mục đích} & Kiểm tra xem Hub có đủ số chỗ đỗ trống để tiếp nhận thêm phương\newline tiện điềuphối haykhông (UC11 – EX2). \\
\hline
\textbf{Tham số đầu vào}\newline \textbf{(@param)} & requiredCount: int – Số lượng phương tiện dự kiến điều phối đến\newline Hub. \\
\hline
\textbf{Giá trị trả về (@return)} & boolean – Trả về true nếu Hub đang ACTIVE và availableSpots >=\newline requiredCount, ngược lại false. \\
\hline
\textbf{Ngoại lệ (@exception)} & IllegalArgumentException – Khi requiredCount <= 0. \\
\hline
\textbf{Tiền điều kiện (Pre-}\newline \textbf{conditions)} & Hub đã được cấu hình capacity và availableSpots hợp lệ. \\
\hline
\textbf{Hậu điều kiện (Post-}\newline \textbf{conditions)} & Không làm thay đổi trạng thái Hub. \\
\hline
\textbf{Luồng thuật toán}\newline \textbf{(Pseudo-code)} & 1. Nếu requiredCount <= 0 -> ném IllegalArgumentException.\newline 2. Trả về (this.status == HubStatus.ACTIVE) \&\&\newline (this.availableSpots >= requiredCount). \\
\hline
\end{longtable}
\vspace{-0.5\baselineskip}
\subsubsection{MobilityHub.adjustExpectedCapacity(delta)}
\begin{longtable}{|L{0.24\textwidth}|L{0.70\textwidth}|}
\hline
\textbf{Thuộc tính} & \textbf{Nội dung chi tiết} \\
\hline
\endfirsthead
\hline
\textbf{Thuộc tính} & \textbf{Nội dung chi tiết} \\
\hline
\endhead
\textbf{Tên phương thức} & MobilityHub.adjustExpectedCapacity(delta) \\
\hline
\textbf{Mục đích} & Cập nhật số chỗ đỗ trống dự kiến tại Hub nguồn (+count) và Hub\newline đích(-count)khi tạo hoặc hủylệnh điềuphối. \\
\hline
\textbf{Tham số đầu vào}\newline \textbf{(@param)} & delta: int – Độ biến thiên số chỗ trống (dương khi xe rời đi, âm khi\newline xe được điềuphối đến). \\
\hline
\textbf{Giá trị trả về(@return)} & void – Khôngtrả vềgiá trị. \\
\hline
\textbf{Ngoại lệ (@exception)} & HubCapacityExceededException – Khi availableSpots + delta < 0\newline hoặc > capacity. \\
\hline
\textbf{Tiền điều kiện (Pre-}\newline \textbf{conditions)} & Hub đang ở trạng thái ACTIVE. \\
\hline
\textbf{Hậu điều kiện (Post-}\newline \textbf{conditions)} & Thuộc tính availableSpots của Hub được cập nhật nằm trong đoạn\newline [0, capacity]. \\
\hline
\end{longtable}
\vspace{-0.5\baselineskip}
\begin{longtable}{|L{0.24\textwidth}|L{0.70\textwidth}|}
\hline
\textbf{Thuộc tính} & \textbf{Nội dung chi tiết} \\
\hline
\endfirsthead
\hline
\textbf{Thuộc tính} & \textbf{Nội dung chi tiết} \\
\hline
\endhead
 & 1. Tính nextSpots = this.availableSpots + delta. \\
\hline
\textbf{Luồng thuật toán} & 2. Nếu nextSpots < 0 hoặc nextSpots > this.capacity -> ném \\
\hline
\textbf{(Pseudo-code)} & HubCapacityExceededException.\newline 3. Gán this.availableSpots = nextSpots. \\
\hline
\end{longtable}
\vspace{-0.5\baselineskip}
\subsubsection{MobilityHub.isolateHub()}
\begin{longtable}{|L{0.24\textwidth}|L{0.70\textwidth}|}
\hline
\textbf{Thuộc tính} & \textbf{Nội dung chi tiết} \\
\hline
\endfirsthead
\hline
\textbf{Thuộc tính} & \textbf{Nội dung chi tiết} \\
\hline
\endhead
\textbf{Tên phương thức} & MobilityHub.isolateHub() \\
\hline
\textbf{Mục đích} & Chuyển trạng thái Hub sang ISOLATED để cô lập giao dịch khi Hub\newline mất kết nối hoàn toàn, bảo vệ toàn mạnglưới(NFR05). \\
\hline
\textbf{Tham số đầu vào}\newline \textbf{(@param)} & Không có. \\
\hline
\textbf{Giá trị trả về(@return)} & void – Khôngtrả vềgiá trị. \\
\hline
\textbf{Ngoại lệ (@exception)} & Khôngcó. \\
\hline
\textbf{Tiền điều kiện (Pre-}\newline \textbf{conditions)} & Hệ thống xác nhận Hub bị mất kết nối hoàn toàn (UC12 – EX1). \\
\hline
\textbf{Hậu điều kiện (Post-}\newline \textbf{conditions)} & this.status = HubStatus.ISOLATED; mọi giao dịch mới tại Hub này\newline bị tạm ngưng. \\
\hline
\textbf{Luồng thuật toán}\newline \textbf{(Pseudo-code)} & 1. Gán this.status = HubStatus.ISOLATED. \\
\hline
\end{longtable}
\vspace{-0.5\baselineskip}
\subsubsection{Resource.updateStatus(newStatus, timestamp)}
\begin{longtable}{|L{0.24\textwidth}|L{0.70\textwidth}|}
\hline
\textbf{Thuộc tính} & \textbf{Nội dung chi tiết} \\
\hline
\endfirsthead
\hline
\textbf{Thuộc tính} & \textbf{Nội dung chi tiết} \\
\hline
\endhead
\textbf{Tên phương thức} & Resource.updateStatus(newStatus, timestamp) \\
\hline
\textbf{Mục đích} & Cập nhật trạng thái hiện hành và dấu thời gian cập nhật gần nhất cho\newline bất kỳtài nguyên nào(SharedVehicle, ChargingPort, ParkingSpot). \\
\hline
\textbf{Tham số đầu vào}\newline \textbf{(@param)} & newStatus: ResourceStatus – Trạng thái mới cần áp dụng.\newline timestamp: LocalDateTime – Mốc thờigian hiệu lực của cậpnhật. \\
\hline
\textbf{Giá trị trả về(@return)} & void – Khôngtrả vềgiá trị. \\
\hline
\textbf{Ngoại lệ (@exception)} & InvalidArgumentException – Khi newStatus hoặc timestampnull. \\
\hline
\textbf{Tiền điều kiện (Pre-}\newline \textbf{conditions)} & Đối tượng Resource tồn tại trong hệ thống. \\
\hline
\textbf{Hậu điều kiện (Post-}\newline \textbf{conditions)} & status = newStatus; lastUpdatedAt = timestamp. \\
\hline
\textbf{Luồng thuật toán}\newline \textbf{(Pseudo-code)} & 1. Kiểm tra newStatus != null và timestamp != null.\newline 2. Gán this.status = newStatus; this.lastUpdatedAt = timestamp. \\
\hline
\end{longtable}
\vspace{-0.5\baselineskip}
\subsubsection{Resource.isAvailable()}
\begin{longtable}{|L{0.24\textwidth}|L{0.70\textwidth}|}
\hline
\textbf{Thuộc tính} & \textbf{Nội dung chi tiết} \\
\hline
\endfirsthead
\hline
\textbf{Thuộc tính} & \textbf{Nội dung chi tiết} \\
\hline
\endhead
\textbf{Tên phương thức} & Resource.isAvailable() \\
\hline
\textbf{Mục đích} & Kiểm tra nhanh tài nguyên hiện có đang sẵn sàng để phục vụ đặt chỗ\newline hoặc điềuphối haykhông. \\
\hline
\textbf{Tham số đầu vào}\newline \textbf{(@param)} & Không có. \\
\hline
\end{longtable}
\vspace{-0.5\baselineskip}
\begin{longtable}{|L{0.24\textwidth}|L{0.70\textwidth}|}
\hline
\textbf{Thuộc tính} & \textbf{Nội dung chi tiết} \\
\hline
\endfirsthead
\hline
\textbf{Thuộc tính} & \textbf{Nội dung chi tiết} \\
\hline
\endhead
\textbf{Giá trị trả về(@return)} & boolean – Trả về true nếu status == ResourceStatus.AVAILABLE. \\
\hline
\textbf{Ngoại lệ (@exception)} & Khôngcó. \\
\hline
\textbf{Tiền điều kiện (Pre-}\newline \textbf{conditions)} & Đối tượng Resource đã được khởi tạo. \\
\hline
\textbf{Hậu điều kiện (Post-}\newline \textbf{conditions)} & Không thay đổi trạng thái tài nguyên. \\
\hline
\textbf{Luồng thuật toán}\newline \textbf{(Pseudo-code)} & 1. Trả về (this.status == ResourceStatus.AVAILABLE). \\
\hline
\end{longtable}
\vspace{-0.5\baselineskip}
\subsubsection{SharedVehicle.updateTelemetry(status, batteryLevel, timestamp)}
\begin{longtable}{|L{0.24\textwidth}|L{0.70\textwidth}|}
\hline
\textbf{Thuộc tính} & \textbf{Nội dung chi tiết} \\
\hline
\endfirsthead
\hline
\textbf{Thuộc tính} & \textbf{Nội dung chi tiết} \\
\hline
\endhead
\textbf{Tên phương thức} & SharedVehicle.updateTelemetry(status, batteryLevel, timestamp) \\
\hline
\textbf{Mục đích} & Cập nhật đồng bộ trạng thái vận hành, mức pin (\%) và dấu thời gian\newline mới nhất chophươngtiện điện dùngchungtừgói tin UC16. \\
\hline
\textbf{Tham số đầu vào}\newline \textbf{(@param)} & status: ResourceStatus – Trạng thái mới của xe.\newline batteryLevel: double – Mức pin hiện tại (0..100).\newline timestamp: LocalDateTime – Thời điểmghi nhận telemetry. \\
\hline
\textbf{Giá trị trả về(@return)} & void – Khôngtrả vềgiá trị. \\
\hline
\textbf{Ngoại lệ (@exception)} & IllegalArgumentException – Khi batteryLevel < 0 hoặc > 100. \\
\hline
\textbf{Tiền điều kiện (Pre-}\newline \textbf{conditions)} & Gói tin đã vượt qua bước kiểm tra cấu trúc và chống gói trễ\newline (NFR08). \\
\hline
\textbf{Hậu điều kiện (Post-}\newline \textbf{conditions)} & status, batteryLevel và lastUpdatedAt của SharedVehicle được cập\newline nhật đồngbộ. \\
\hline
\textbf{Luồng thuật toán}\newline \textbf{(Pseudo-code)} & 1. Nếu batteryLevel < 0 \textbar\{\}\textbar\{\} batteryLevel > 100 -> ném\newline IllegalArgumentException.\newline 2. Gọi super.updateStatus(status, timestamp).\newline 3. Gán this.batteryLevel = batteryLevel. \\
\hline
\end{longtable}
\vspace{-0.5\baselineskip}
\subsubsection{AuditLog.record(actorId, action, targetId, details)}
\begin{longtable}{|L{0.24\textwidth}|L{0.70\textwidth}|}
\hline
\textbf{Thuộc tính} & \textbf{Nội dung chi tiết} \\
\hline
\endfirsthead
\hline
\textbf{Thuộc tính} & \textbf{Nội dung chi tiết} \\
\hline
\endhead
\textbf{Tên phương thức} & AuditLog.record(actorId, action, targetId, details) \\
\hline
\textbf{Mục đích} & Tạo và lưu trữ bản ghi nhật ký kiểm toán phục vụ truy vết các thao\newline tác quan trọng (điều phối, tạm ngưng/khôi phục sự cố, lọc gói trễ) tối\newline thiểu 90 ngày (NFR11). \\
\hline
\textbf{Tham số đầu vào}\newline \textbf{(@param)} & actorId: String – Mã tác nhân thực hiện (Operator ID hoặc\newline SYSTEM/SIMULATOR).\newline action: String – Tên hành động nghiệp vụ.\newline targetId: String – Mã đối tượng bị tác động.\newline details: String– Thôngtin chi tiết thao tác. \\
\hline
\textbf{Giá trị trả về(@return)} & AuditLog– Đối tượngbảnghi AuditLogđã được lưu bền vững. \\
\hline
\textbf{Ngoại lệ (@exception)} & PersistenceException – Xảy ra khi không thể ghi log vào cơ sở dữ\newline liệu. \\
\hline
\textbf{Tiền điều kiện (Pre-}\newline \textbf{conditions)} & Thông tin tác nhân và hành động hợp lệ. \\
\hline
\end{longtable}
\vspace{-0.5\baselineskip}
\begin{longtable}{|L{0.24\textwidth}|L{0.70\textwidth}|}
\hline
\textbf{Thuộc tính} & \textbf{Nội dung chi tiết} \\
\hline
\endfirsthead
\hline
\textbf{Thuộc tính} & \textbf{Nội dung chi tiết} \\
\hline
\endhead
\textbf{Hậu điều kiện (Post-}\newline \textbf{conditions)} & Một bản ghi AuditLog mới với timestamp = now() được lưu trữ tối\newline thiểu 90 ngày. \\
\hline
\textbf{Luồng thuật toán} & 1. Tạo đối tượng log = new AuditLog(UUID, actorId, action,\newline targetId, details, now()). \\
\hline
\textbf{(Pseudo-code)} & 2. Lưu log vào kho lưu trữ AuditLog (retention >= 90 days).\newline 3. Trả về log. \\
\hline
\end{longtable}
\vspace{-0.5\baselineskip}
\subsubsection{NotificationService.sendNotification(recipientId, content)}
\begin{longtable}{|L{0.24\textwidth}|L{0.70\textwidth}|}
\hline
\textbf{Thuộc tính} & \textbf{Nội dung chi tiết} \\
\hline
\endfirsthead
\hline
\textbf{Thuộc tính} & \textbf{Nội dung chi tiết} \\
\hline
\endhead
\textbf{Tên phương thức} & NotificationService.sendNotification(recipientId, content) \\
\hline
\textbf{Mục đích} & Gửi thông báo nghiệp vụ ra ngoài hệ thống tới đội ngũ vận chuyển\newline thực địa(UC11)hoặc Sinh viên bị hủyđơn do sự cố bảo trì(UC12). \\
\hline
\textbf{Tham số đầu vào}\newline \textbf{(@param)} & recipientId: String – Mã người nhận hoặc nhóm nhận thông báo.\newline content: String– Nội dungthôngbáo. \\
\hline
\textbf{Giá trị trả về (@return)} & boolean – Trả về true nếu dịch vụ thông báo tiếp nhận gửi thành\newline công, false nếugửi thất bại. \\
\hline
\textbf{Ngoại lệ (@exception)} & Không ném ngoại lệ chặn luồng nghiệp vụ chính; lỗi gửi được ghi\newline nhận để thử lại. \\
\hline
\textbf{Tiền điều kiện (Pre-}\newline \textbf{conditions)} & Giao dịch nghiệp vụ chính (tạo lệnh điều phối hoặc tạm ngưng tài\newline nguyên)đã commit thành công. \\
\hline
\textbf{Hậu điều kiện (Post-}\newline \textbf{conditions)} & Thông báo được chuyển qua kênh tương ứng tới người nhận; nếu lỗi,\newline hệ thống ghi logđể xử lýbù. \\
\hline
\textbf{Luồng thuật toán}\newline \textbf{(Pseudo-code)} & 1. Đóng gói payload thông báo gồm recipientId, content và\newline timestamp.\newline 2. Gọi API của hệ thống Notification Service bên ngoài thông qua\newline Notify Adapter.\newline 3. Nếu phản hồi HTTP 2xx -> trả về true; ngược lại ghi log lỗi tích\newline hợpvà trả về false. \\
\hline
\end{longtable}
\vspace{-0.5\baselineskip}
\subsubsection{RealtimeGateway.broadcastResourceUpdate(resourceId, status, updatedAt)}
\begin{longtable}{|L{0.24\textwidth}|L{0.70\textwidth}|}
\hline
\textbf{Thuộc tính} & \textbf{Nội dung chi tiết} \\
\hline
\endfirsthead
\hline
\textbf{Thuộc tính} & \textbf{Nội dung chi tiết} \\
\hline
\endhead
\textbf{Tên phương thức} & RealtimeGateway.broadcastResourceUpdate(resourceId, status,\newline updatedAt) \\
\hline
\textbf{Mục đích} & Đẩy tín hiệu cập nhật trạng thái tài nguyên theo thời gian thực qua\newline WebSocket/SSE tới các giao diện Mobile App và Web Dashboard\newline đangkết nối(NFR02 <= 15giây). \\
\hline
\textbf{Tham số đầu vào}\newline \textbf{(@param)} & resourceId: String – Mã tài nguyên vừa thay đổi trạng thái.\newline status: ResourceStatus – Trạng thái mới nhất.\newline updatedAt: LocalDateTime – Thời điểm cậpnhật hiệu lực. \\
\hline
\textbf{Giá trị trả về(@return)} & void – Khôngtrả vềgiá trị. \\
\hline
\textbf{Ngoại lệ (@exception)} & Không ném ngoại lệ ra luồng ingestion; lỗi kết nối client đơn lẻ bị bỏ\newline qua và tự độngdọn session. \\
\hline
\textbf{Tiền điều kiện (Pre-}\newline \textbf{conditions)} & Dữ liệu trạng thái mới đã được cập nhật thành công vào Operational\newline DB. \\
\hline
\end{longtable}
\vspace{-0.5\baselineskip}
\begin{longtable}{|L{0.24\textwidth}|L{0.70\textwidth}|}
\hline
\textbf{Thuộc tính} & \textbf{Nội dung chi tiết} \\
\hline
\endfirsthead
\hline
\textbf{Thuộc tính} & \textbf{Nội dung chi tiết} \\
\hline
\endhead
\textbf{Hậu điều kiện (Post-}\newline \textbf{conditions)} & Tất cả client đang đăng ký kênh cập nhật nhận được sự kiện trạng\newline thái mới trongvòng<= 15giây. \\
\hline
 & 1. Tạo thông điệp sự kiện JSON \{resourceId, status, updatedAt\}. \\
\hline
\textbf{Luồng thuật toán} & 2. Phát sự kiện tới các kênh/topic WebSocket và SSE đang hoạt \\
\hline
\textbf{(Pseudo-code)} & động.\newline 3. Ghi nhận thời điểmphát để theo dõi SLA <= 15giây (NFR02). \\
\hline
\end{longtable}
\vspace{-0.5\baselineskip}
\section{Kịch bản kiểm thử (Test Cases)}
\subsection{Test case cho UC1 – Tìm phương tiện và Hub phù hợp}
\subsubsection{TCUC0101 – Ngoại lệ (AF3) - Không có kết quả phù hợp với tiêu chí tìm kiếm.}
\begin{longtable}{|L{0.24\textwidth}|L{0.70\textwidth}|}
\hline
\textbf{Trường thông tin} & \textbf{Đặc tả kịch bản} \\
\hline
\endfirsthead
\hline
\textbf{Trường thông tin} & \textbf{Đặc tả kịch bản} \\
\hline
\endhead
Test Case ID & TC\_UC01\_01 \\
\hline
Use Case liênquan & UC1 - Tìmphươngtiện và Hubphù hợp \\
\hline
Mô tả kịch bản & Ngoại lệ(AF3)- Khôngcó kếtquảphù hợpvới tiêu chí tìm kiếm. \\
\hline
Tiền điều kiện\newline (Preconditions) & Vị trí xác định được; không có phương tiện khả dụng nào trong bán\newline kính tìm kiếm. \\
\hline
Dữ liệu đầu vào (Input\newline Data) & location=(10.762622, 106.660172); radiusKm=0.1;\newline searchType=VEHICLE; vehicleType="E-bike";\newline minBatteryLevel=20 \\
\hline
Kết quả mong đợi\newline (Expected Result) & Hệ thống trả về danh sách rỗng và thông báo không có kết quả phù\newline hợp; tiêu chí tìm kiếm được giữ nguyên trên giao diện để Sinh viên\newline điều chỉnh; khôngthay đổi dữ liệu. \\
\hline
Kết quả thực tế (Actual /\newline Status) & Pending \\
\hline
\end{longtable}
\vspace{-0.5\baselineskip}
\subsubsection{TCUC0102 – Ngoại lệ (EX1) - Không xác định được vị trí tìm kiếm.}
\begin{longtable}{|L{0.24\textwidth}|L{0.70\textwidth}|}
\hline
\textbf{Trường thông tin} & \textbf{Đặc tả kịch bản} \\
\hline
\endfirsthead
\hline
\textbf{Trường thông tin} & \textbf{Đặc tả kịch bản} \\
\hline
\endhead
Test Case ID & TC\_UC01\_02 \\
\hline
Use Case liênquan & UC1 - Tìmphươngtiện và Hubphù hợp \\
\hline
Mô tả kịch bản & Ngoại lệ(EX1)- Khôngxác định được vị trí tìm kiếm. \\
\hline
Tiền điều kiện\newline (Preconditions) & Sinh viên từ chối cấp quyền định vị và không chọn vị trí thủ công\newline trên bản đồ. \\
\hline
Dữ liệu đầu vào (Input\newline Data) & location=null; radiusKm=2; searchType=VEHICLE \\
\hline
Kết quả mong đợi\newline (Expected Result) & Hệ thống ném InvalidLocationException, hiển thị thông báo lỗi và\newline cho phép thử lại hoặc chuyển sang chọn vị trí thủ công (nhánh 3a);\newline khôngthực hiện truyvấn dữ liệu. \\
\hline
Kết quả thực tế (Actual /\newline Status) & Pending \\
\hline
\end{longtable}
\vspace{-0.5\baselineskip}
\subsubsection{TCUC0103 – Ngoại lệ (EX2) - Tiêu chí tìm kiếm không hợp lệ.}
\begin{longtable}{|L{0.24\textwidth}|L{0.70\textwidth}|}
\hline
\textbf{Trường thông tin} & \textbf{Đặc tả kịch bản} \\
\hline
\endfirsthead
\hline
\textbf{Trường thông tin} & \textbf{Đặc tả kịch bản} \\
\hline
\endhead
Test Case ID & TC\_UC01\_03 \\
\hline
Use Case liênquan & UC1 - Tìmphươngtiện và Hubphù hợp \\
\hline
Mô tả kịch bản & Ngoại lệ(EX2)- Tiêu chí tìm kiếm khônghợp lệ. \\
\hline
Tiền điều kiện\newline (Preconditions) & Vị trí xác định được; giao diện tìm kiếm đang hiển thị. \\
\hline
Dữ liệu đầu vào (Input & location=(10.762622, 106.660172); radiusKm=-5; \\
\hline
Data) & minBatteryLevel=150 \\
\hline
\end{longtable}
\vspace{-0.5\baselineskip}
\begin{longtable}{|L{0.24\textwidth}|L{0.70\textwidth}|}
\hline
\textbf{Trường thông tin} & \textbf{Đặc tả kịch bản} \\
\hline
\endfirsthead
\hline
\textbf{Trường thông tin} & \textbf{Đặc tả kịch bản} \\
\hline
\endhead
Kết quả mong đợi\newline (Expected Result) & Hệ thống ném InvalidSearchCriteriaException, chỉ rõ các trường cần\newline chỉnh sửa (bán kính phải lớn hơn 0, mức pin trong khoảng 0-100) và\newline yêu cầu nhậplại; khôngtruyvấn xuốngtầngdữ liệu. \\
\hline
Kết quả thực tế (Actual /\newline Status) & Pending \\
\hline
\end{longtable}
\vspace{-0.5\baselineskip}
\subsubsection{TCUC0104 – Ngoại lệ (EX3) - Map API không khả dụng.}
\begin{longtable}{|L{0.24\textwidth}|L{0.70\textwidth}|}
\hline
\textbf{Trường thông tin} & \textbf{Đặc tả kịch bản} \\
\hline
\endfirsthead
\hline
\textbf{Trường thông tin} & \textbf{Đặc tả kịch bản} \\
\hline
\endhead
Test Case ID & TC\_UC01\_04 \\
\hline
Use Case liênquan & UC1 - Tìmphươngtiện và Hubphù hợp \\
\hline
Mô tả kịch bản & Ngoại lệ(EX3)- MapAPI khôngkhả dụng. \\
\hline
Tiền điều kiện\newline (Preconditions) & Dịch vụ Map API bị timeout hoặc trả về lỗi kết nối; dữ liệu Hub vẫn\newline truy xuất được. \\
\hline
Dữ liệu đầu vào (Input\newline Data) & location=(10.762622, 106.660172); radiusKm=2; searchType=HUB \\
\hline
Kết quả mong đợi\newline (Expected Result) & Hệ thống thông báo lỗi bản đồ và chuyển sang hiển thị danh sách\newline Hub theo khu vực hoặc toạ độ đã lưu; chức năng tìm kiếm vẫn trả\newline kếtquả dạngdanh sách. \\
\hline
Kết quả thực tế (Actual /\newline Status) & Pending \\
\hline
\end{longtable}
\vspace{-0.5\baselineskip}
\subsubsection{TCUC0105 – Ngoại lệ (EX4) - Không truy xuất được dữ liệu phương tiện/Hub.}
\begin{longtable}{|L{0.24\textwidth}|L{0.70\textwidth}|}
\hline
\textbf{Trường thông tin} & \textbf{Đặc tả kịch bản} \\
\hline
\endfirsthead
\hline
\textbf{Trường thông tin} & \textbf{Đặc tả kịch bản} \\
\hline
\endhead
Test Case ID & TC\_UC01\_05 \\
\hline
Use Case liênquan & UC1 - Tìmphươngtiện và Hubphù hợp \\
\hline
Mô tả kịch bản & Ngoại lệ(EX4)- Khôngtruy xuất được dữ liệuphươngtiện/Hub. \\
\hline
Tiền điều kiện\newline (Preconditions) & Kết nối tới cơ sở dữ liệu bị gián đoạn trong lúc tìm kiếm. \\
\hline
Dữ liệu đầu vào (Input\newline Data) & location=(10.762622, 106.660172); radiusKm=2;\newline searchType=VEHICLE; vehicleType="E-bike";\newline minBatteryLevel=20 \\
\hline
Kết quả mong đợi\newline (Expected Result) & Hệ thống ném DataSourceUnavailableException, thông báo lỗi, giữ\newline lại tiêu chí đã nhập và cho phép Sinh viên thử lại; không hiển thị kết\newline quả sai lệch. \\
\hline
Kết quả thực tế (Actual /\newline Status) & Pending \\
\hline
\end{longtable}
\vspace{-0.5\baselineskip}
\subsection{Test case cho UC2 – Xem trạng thái phương tiện và Hub}
\subsubsection{TCUC0201 – Ngoại lệ (EX2) - Tài nguyên không tồn tại hoặc đã ngừng hoạt động.}
\begin{longtable}{|L{0.24\textwidth}|L{0.70\textwidth}|}
\hline
\textbf{Trường thông tin} & \textbf{Đặc tả kịch bản} \\
\hline
\endfirsthead
\hline
\textbf{Trường thông tin} & \textbf{Đặc tả kịch bản} \\
\hline
\endhead
Test Case ID & TC\_UC02\_01 \\
\hline
Use Case liênquan & UC2 - Xem trạngtháiphươngtiện và Hub \\
\hline
Mô tả kịch bản & Ngoại lệ(EX2)- Tài nguyên khôngtồn tại hoặc đã ngừnghoạt động. \\
\hline
\end{longtable}
\vspace{-0.5\baselineskip}
\begin{longtable}{|L{0.24\textwidth}|L{0.70\textwidth}|}
\hline
\textbf{Trường thông tin} & \textbf{Đặc tả kịch bản} \\
\hline
\endfirsthead
\hline
\textbf{Trường thông tin} & \textbf{Đặc tả kịch bản} \\
\hline
\endhead
Tiền điều kiện\newline (Preconditions) & resourceId không tồn tại trong hệ thống. \\
\hline
Dữ liệu đầu vào (Input\newline Data) & resourceType=VEHICLE; resourceId="V-9999-INVALID" \\
\hline
Kết quả mong đợi\newline (Expected Result) & Hệ thống ném ResourceNotFoundException, thông báo và đưa Sinh\newline viên về danh sách tìm kiếm; khôngtrả dữ liệu trạngthái nào. \\
\hline
Kết quả thực tế (Actual /\newline Status) & Pending \\
\hline
\end{longtable}
\vspace{-0.5\baselineskip}
\subsubsection{TCUC0202 – Ngoại lệ (EX1) - Chưa chọn tài nguyên hợp lệ.}
\begin{longtable}{|L{0.24\textwidth}|L{0.70\textwidth}|}
\hline
\textbf{Trường thông tin} & \textbf{Đặc tả kịch bản} \\
\hline
\endfirsthead
\hline
\textbf{Trường thông tin} & \textbf{Đặc tả kịch bản} \\
\hline
\endhead
Test Case ID & TC\_UC02\_02 \\
\hline
Use Case liênquan & UC2 - Xem trạngtháiphươngtiện và Hub \\
\hline
Mô tả kịch bản & Ngoại lệ(EX1)- Chưa chọn tài nguyên hợp lệ. \\
\hline
Tiền điều kiện\newline (Preconditions) & Sinh viên mở chức năng xem trạng thái nhưng chưa chọn phương\newline tiện hoặc Hub. \\
\hline
Dữ liệu đầu vào (Input\newline Data) & resourceType=VEHICLE; resourceId="" \\
\hline
Kết quả mong đợi\newline (Expected Result) & Hệ thống ném InvalidArgumentException và yêu cầu Sinh viên chọn\newline phươngtiện hoặc Hub trước khi tiếptục; khôngtruyvấn dữ liệu. \\
\hline
Kết quả thực tế (Actual /\newline Status) & Pending \\
\hline
\end{longtable}
\vspace{-0.5\baselineskip}
\subsubsection{TCUC0203 – Ngoại lệ (EX3) - Không truy xuất được dữ liệu trạng thái.}
\begin{longtable}{|L{0.24\textwidth}|L{0.70\textwidth}|}
\hline
\textbf{Trường thông tin} & \textbf{Đặc tả kịch bản} \\
\hline
\endfirsthead
\hline
\textbf{Trường thông tin} & \textbf{Đặc tả kịch bản} \\
\hline
\endhead
Test Case ID & TC\_UC02\_03 \\
\hline
Use Case liênquan & UC2 - Xem trạngtháiphươngtiện và Hub \\
\hline
Mô tả kịch bản & Ngoại lệ(EX3)- Khôngtruy xuất được dữ liệu trạngthái. \\
\hline
Tiền điều kiện\newline (Preconditions) & resourceId hợp lệ; kết nối tới cơ sở dữ liệu bị gián đoạn. \\
\hline
Dữ liệu đầu vào (Input\newline Data) & resourceType=HUB; resourceId="HUB-001" \\
\hline
Kết quả mong đợi\newline (Expected Result) & Hệ thống ném DataSourceUnavailableException, thông báo lỗi, giữ\newline lại tài nguyên đangchọn và cho phépSinh viên thử lại. \\
\hline
Kết quả thực tế (Actual /\newline Status) & Pending \\
\hline
\end{longtable}
\vspace{-0.5\baselineskip}
\begin{longtable}{|L{0.24\textwidth}|L{0.70\textwidth}|}
\hline
\textbf{Trường thông tin} & \textbf{Đặc tả kịch bản} \\
\hline
\endfirsthead
\hline
\textbf{Trường thông tin} & \textbf{Đặc tả kịch bản} \\
\hline
\endhead
Test Case ID & TC\_UC02\_04 \\
\hline
Use Case liênquan & UC2 - Xem trạngtháiphươngtiện và Hub \\
\hline
Mô tả kịch bản & Biên(EX4)- Dữ liệu trạngtháiquá cũ. \\
\hline
\end{longtable}
\vspace{-0.5\baselineskip}
\begin{longtable}{|L{0.24\textwidth}|L{0.70\textwidth}|}
\hline
\textbf{Trường thông tin} & \textbf{Đặc tả kịch bản} \\
\hline
\endfirsthead
\hline
\textbf{Trường thông tin} & \textbf{Đặc tả kịch bản} \\
\hline
\endhead
Tiền điều kiện\newline (Preconditions) & Tài nguyên tồn tại; thời điểm cập nhật gần nhất (lastUpdatedAt) cũ\newline hơn ngưỡngSTALE\_THRESHOLD. \\
\hline
Dữ liệu đầu vào (Input\newline Data) & resourceType=VEHICLE; resourceId="V-0012"; lastUpdatedAt =\newline thời điểm hiện tại - 30phút \\
\hline
Kết quả mong đợi\newline (Expected Result) & Hệ thống vẫn hiển thị thời điểm cập nhật gần nhất nhưng đánh dấu\newline dữ liệu chưa đáng tin cậy (stale) và không khẳng định phương tiện\newline còn khả dụng. \\
\hline
Kết quả thực tế (Actual /\newline Status) & Pending \\
\hline
\end{longtable}
\vspace{-0.5\baselineskip}
\subsubsection{TCUC0205 – Biên (AF4) - Tài nguyên đang không khả dụng.}
\begin{longtable}{|L{0.24\textwidth}|L{0.70\textwidth}|}
\hline
\textbf{Trường thông tin} & \textbf{Đặc tả kịch bản} \\
\hline
\endfirsthead
\hline
\textbf{Trường thông tin} & \textbf{Đặc tả kịch bản} \\
\hline
\endhead
Test Case ID & TC\_UC02\_05 \\
\hline
Use Case liênquan & UC2 - Xem trạngtháiphươngtiện và Hub \\
\hline
Mô tả kịch bản & Biên(AF4)- Tài nguyên đangkhôngkhả dụng. \\
\hline
Tiền điều kiện\newline (Preconditions) & Phương tiện tồn tại và đang ở trạng thái "Đã đặt" bởi Sinh viên khác. \\
\hline
Dữ liệu đầu vào (Input\newline Data) & resourceType=VEHICLE; resourceId="V-0020" \\
\hline
Kết quả mong đợi\newline (Expected Result) & Hệ thống vẫn hiển thị trạng thái hiện tại ("Đã đặt") kèm thời điểm\newline cập nhật, không báo lỗi, và cho phép Sinh viên chọn tài nguyên\newline khác; khôngthay đổi dữ liệu đặt xe. \\
\hline
Kết quả thực tế (Actual /\newline Status) & Pending\newline 2.4. Test case cho UC4 - Đặt chỗ đỗ xe \\
\hline
\end{longtable}
\vspace{-0.5\baselineskip}
\subsection{Test case cho UC3 – Đặt xe}
\subsubsection{TCUC0301 – Đặt xe thành công}
\begin{longtable}{|L{0.24\textwidth}|L{0.70\textwidth}|}
\hline
\textbf{Trường thông tin} & \textbf{Đặc tả kịch bản} \\
\hline
\endfirsthead
\hline
\textbf{Trường thông tin} & \textbf{Đặc tả kịch bản} \\
\hline
\endhead
Test Case ID & TC\_UC03\_01 \\
\hline
Use Case liênquan & UC3 – Đặt xe \\
\hline
Mô tả kịch bản & Đặt xe thành công \\
\hline
Tiền điều kiện\newline (Preconditions) & Sinh viên đã đăng nhập; xe được chọn đang khả dụng tại Hub hoạt\newline động. \\
\hline
Dữ liệu đầu vào (Input\newline Data) & Người dùng: sinh viên An.\newline Phương tiện: xe điện số 1 tại Hub A.\newline Thao tác: chọn xe và xác nhận đặt. \\
\hline
Kết quả mong đợi\newline (Expected Result) & 1. Tạo một đơn thuộc sinh viên An.\newline 2. Xe chuyển sang trạng thái đã giữ.\newline 3. Lưu lịch sử và hiển thị xác nhận đặt xe. \\
\hline
Kếtquả thực tế / Status & Pending– Chưa thực hiện kiểm thử. \\
\hline
\end{longtable}
\vspace{-0.5\baselineskip}
\subsubsection{TCUC0302 – Hủy trước khi xác nhận}
\begin{longtable}{|L{0.24\textwidth}|L{0.70\textwidth}|}
\hline
\textbf{Trường thông tin} & \textbf{Đặc tả kịch bản} \\
\hline
\endfirsthead
\hline
\textbf{Trường thông tin} & \textbf{Đặc tả kịch bản} \\
\hline
\endhead
Test Case ID & TC\_UC03\_02 \\
\hline
Use Case liênquan & UC3 – Đặt xe \\
\hline
Mô tả kịch bản & Hủytrước khi xác nhận \\
\hline
Tiền điều kiện\newline (Preconditions) & Sinh viên đang xem hộp thoại đặt xe; chưa gửi yêu cầu lưu. \\
\hline
Dữ liệu đầu vào (Input\newline Data) & Xe điện số 1 tại Hub A.\newline Thao tác: chọn Hủy. \\
\hline
Kết quả mong đợi\newline (Expected Result) & 1. Đóng hộp thoại và xóa lựa chọn tạm.\newline 2. Khôngtạo đơn; trạngthái xe khôngđổi. \\
\hline
Kếtquả thực tế / Status & Pending– Chưa thực hiện kiểm thử. \\
\hline
\end{longtable}
\vspace{-0.5\baselineskip}
\subsubsection{TCUC0303 – Nhiều người cùng đặt một xe}
\begin{longtable}{|L{0.24\textwidth}|L{0.70\textwidth}|}
\hline
\textbf{Trường thông tin} & \textbf{Đặc tả kịch bản} \\
\hline
\endfirsthead
\hline
\textbf{Trường thông tin} & \textbf{Đặc tả kịch bản} \\
\hline
\endhead
Test Case ID & TC\_UC03\_03 \\
\hline
Use Case liênquan & UC3 – Đặt xe \\
\hline
Mô tả kịch bản & Nhiều người cùngđặt một xe \\
\hline
Tiền điều kiện\newline (Preconditions) & Hai sinh viên An và Bình đã đăng nhập; chỉ có xe điện số 1 đang khả\newline dụng. \\
\hline
Dữ liệu đầu vào (Input\newline Data) & Tài nguyên: cùng xe điện số 1.\newline Thao tác: gửi hai xác nhận đồng thời; lặp lại 10 lần, khôi phục xe\newline trước mỗi lần. \\
\hline
Kết quả mong đợi\newline (Expected Result) & 1. Mỗi lần chỉ có một đơn đặt xe thành công.\newline 2. Người còn lại nhận thông báo xe không khả dụng.\newline 3. Không có hai đơn cùng giữ xe; lịch sử phản ánh đúng đơn thành\newline công. \\
\hline
Kếtquả thực tế / Status & Pending– Chưa thực hiện kiểm thử. \\
\hline
\end{longtable}
\vspace{-0.5\baselineskip}
\subsubsection{TCUC0304 – Lỗi khi lưu đơn}
\begin{longtable}{|L{0.24\textwidth}|L{0.70\textwidth}|}
\hline
\textbf{Trường thông tin} & \textbf{Đặc tả kịch bản} \\
\hline
\endfirsthead
\hline
\textbf{Trường thông tin} & \textbf{Đặc tả kịch bản} \\
\hline
\endhead
Test Case ID & TC\_UC03\_04 \\
\hline
Use Case liênquan & UC3 – Đặt xe \\
\hline
Mô tả kịch bản & Lỗi khi lưu đơn \\
\hline
Tiền điều kiện\newline (Preconditions) & Xe điện số 1 khả dụng; mô phỏng lỗi lưu trước khi giao dịch hoàn\newline tất. \\
\hline
Dữ liệu đầu vào (Input\newline Data) & Sinh viên An chọn xe điện số 1 và xác nhận đặt. \\
\hline
Kết quả mong đợi\newline (Expected Result) & 1. Thông báo chưa đặt xe thành công.\newline 2. Không lưu đơn dở dang; xe vẫn khả dụng.\newline 3. Chỉ cho thử lại sau khi xác địnhyêu cầu trước chưa được lưu. \\
\hline
Kếtquả thực tế / Status & Pending– Chưa thực hiện kiểm thử. \\
\hline
\end{longtable}
\vspace{-0.5\baselineskip}
\subsection{Test case cho UC4 – Đặt chỗ đỗ xe}
\subsubsection{TCUC0401 – Đặt chỗ đỗ thành công}
\begin{longtable}{|L{0.24\textwidth}|L{0.70\textwidth}|}
\hline
\textbf{Trường thông tin} & \textbf{Đặc tả kịch bản} \\
\hline
\endfirsthead
\hline
\textbf{Trường thông tin} & \textbf{Đặc tả kịch bản} \\
\hline
\endhead
Test Case ID & TC\_UC04\_01 \\
\hline
Use Case liênquan & UC4 – Đặt chỗ đỗ xe \\
\hline
Mô tả kịch bản & Đặt chỗ đỗ thành công \\
\hline
Tiền điều kiện\newline (Preconditions) & Sinh viên An đã đăng nhập; chỗ đỗ số 1 tại Hub A trống và chưa\newline đượcgiữ trongkhung giờ chọn. \\
\hline
Dữ liệu đầu vào (Input\newline Data) & Chỗ đỗ: số 1 tại Hub A.\newline Thời gian: 09:00–10:00 ngày 01/12/2026.\newline Thao tác: chọn chỗ và xác nhận. \\
\hline
Kết quả mong đợi\newline (Expected Result) & 1. Lưu đơn thuộc sinh viên An và giữ đúng khung giờ.\newline 2. Hiển thị thôngtin đặt chỗ;ghi lịch sử. \\
\hline
Kếtquả thực tế / Status & Pending– Chưa thực hiện kiểm thử. \\
\hline
\end{longtable}
\vspace{-0.5\baselineskip}
\subsubsection{TCUC0402 – Khoảng thời gian không hợp lệ}
\begin{longtable}{|L{0.24\textwidth}|L{0.70\textwidth}|}
\hline
\textbf{Trường thông tin} & \textbf{Đặc tả kịch bản} \\
\hline
\endfirsthead
\hline
\textbf{Trường thông tin} & \textbf{Đặc tả kịch bản} \\
\hline
\endhead
Test Case ID & TC\_UC04\_02 \\
\hline
Use Case liênquan & UC4 – Đặt chỗ đỗ xe \\
\hline
Mô tả kịch bản & Khoảngthờigian khônghợp lệ \\
\hline
Tiền điều kiện\newline (Preconditions) & Sinh viên đã đăng nhập và chọn chỗ đỗ. \\
\hline
Dữ liệu đầu vào (Input\newline Data) & Lần 1: bắt đầu 10:00, kết thúc 09:00.\newline Lần 2: bắt đầu và kết thúc đều 10:00.\newline Thao tác: xác nhận từnglần độc lập. \\
\hline
Kết quả mong đợi\newline (Expected Result) & 1. Từ chối cả hai khoảng thời gian.\newline 2. Yêu cầu sửa thờigian; khôngtạo đơn hoặcgiữ chỗ. \\
\hline
Kếtquả thực tế / Status & Pending– Chưa thực hiện kiểm thử. \\
\hline
\end{longtable}
\vspace{-0.5\baselineskip}
\subsubsection{TCUC0403 – Đặt chỗ đỗ đồng thời bị trùng giờ}
\begin{longtable}{|L{0.24\textwidth}|L{0.70\textwidth}|}
\hline
\textbf{Trường thông tin} & \textbf{Đặc tả kịch bản} \\
\hline
\endfirsthead
\hline
\textbf{Trường thông tin} & \textbf{Đặc tả kịch bản} \\
\hline
\endhead
Test Case ID & TC\_UC04\_03 \\
\hline
Use Case liênquan & UC4 – Đặt chỗ đỗ xe \\
\hline
Mô tả kịch bản & Đặt chỗ đỗ đồngthời bị trùng giờ \\
\hline
Tiền điều kiện\newline (Preconditions) & Hai sinh viên đã đăng nhập; chỗ đỗ số 1 tại Hub A chưa có đơn. \\
\hline
Dữ liệu đầu vào (Input\newline Data) & An đặt 09:00–10:00; Bình đặt 09:30–10:30 cùng ngày, cùng chỗ.\newline Thao tác:gửi đồngthời; lặplại 10 lần với dữ liệu được khôiphục. \\
\hline
Kết quả mong đợi\newline (Expected Result) & 1. Mỗi lần chỉ một yêu cầu thành công.\newline 2. Yêu cầu còn lại nhận thông báo trùng khung giờ.\newline 3. Khônglưu hai đơn có khoảng giữ chỗgiao nhau. \\
\hline
Kếtquả thực tế / Status & Pending– Chưa thực hiện kiểm thử. \\
\hline
\end{longtable}
\vspace{-0.5\baselineskip}
\subsubsection{TCUC0404 – Chưa đăng nhập hoặc không có quyền}
\begin{longtable}{|L{0.24\textwidth}|L{0.70\textwidth}|}
\hline
\textbf{Trường thông tin} & \textbf{Đặc tả kịch bản} \\
\hline
\endfirsthead
\hline
\textbf{Trường thông tin} & \textbf{Đặc tả kịch bản} \\
\hline
\endhead
Test Case ID & TC\_UC04\_04 \\
\hline
Use Case liênquan & UC4 – Đặt chỗ đỗ xe \\
\hline
Mô tả kịch bản & Chưa đăngnhậphoặc khôngcóquyền \\
\hline
Tiền điều kiện\newline (Preconditions) & Người dùng chưa có phiên; kiểm tra thêm một tài khoản vận hành\newline khôngcóquyền sinh viên. \\
\hline
Dữ liệu đầu vào (Input\newline Data) & Chỗ đỗ số 1; thời gian 09:00–10:00.\newline Thao tác:gửiyêu cầu đặt bằngtừngtài khoản. \\
\hline
Kết quả mong đợi\newline (Expected Result) & 1. Người chưa đăng nhập được yêu cầu xác thực.\newline 2. Tài khoản thiếu quyền bị từ chối.\newline 3. Khônglưu đơn trongcả hai trườnghợp. \\
\hline
Kếtquả thực tế / Status & Pending– Chưa thực hiện kiểm thử. \\
\hline
\end{longtable}
\vspace{-0.5\baselineskip}
\subsection{Test case cho UC5 – Nhận phương tiện}
\subsubsection{TCUC0501 – Nhận xe thành công}
\begin{longtable}{|L{0.24\textwidth}|L{0.70\textwidth}|}
\hline
\textbf{Trường thông tin} & \textbf{Đặc tả kịch bản} \\
\hline
\endfirsthead
\hline
\textbf{Trường thông tin} & \textbf{Đặc tả kịch bản} \\
\hline
\endhead
Test Case ID & TC\_UC05\_01 \\
\hline
Use Case liênquan & UC5 – Nhậnphươngtiện \\
\hline
Mô tả kịch bản & Nhận xe thành công \\
\hline
Tiền điều kiện\newline (Preconditions) & An có đơn đặt xe còn hiệu lực tại Hub A; xe đang được giữ và không\newline có sự cố. \\
\hline
Dữ liệu đầu vào (Input\newline Data) & Đơn: đơn đặt xe điện số 1 của An.\newline Hub nhận: Hub A.\newline Thao tác: xác nhận nhận xe môphỏng. \\
\hline
Kết quả mong đợi\newline (Expected Result) & 1. Xe chuyển sang đang sử dụng; đơn chuyển sang đã nhận.\newline 2. Lưu người nhận, thời gian nhận và lịch sử.\newline 3. Hiển thị thành công; khôngđiều khiển khóa vật lý. \\
\hline
Kếtquả thực tế / Status & Pending– Chưa thực hiện kiểm thử. \\
\hline
\end{longtable}
\vspace{-0.5\baselineskip}
\subsubsection{TCUC0502 – Đơn không hợp lệ hoặc thuộc người khác}
\begin{longtable}{|L{0.24\textwidth}|L{0.70\textwidth}|}
\hline
\textbf{Trường thông tin} & \textbf{Đặc tả kịch bản} \\
\hline
\endfirsthead
\hline
\textbf{Trường thông tin} & \textbf{Đặc tả kịch bản} \\
\hline
\endhead
Test Case ID & TC\_UC05\_02 \\
\hline
Use Case liênquan & UC5 – Nhậnphươngtiện \\
\hline
Mô tả kịch bản & Đơn khônghợp lệ hoặc thuộc người khác \\
\hline
Tiền điều kiện\newline (Preconditions) & An đã đăng nhập; có một đơn thuộc sinh viên Bình. \\
\hline
Dữ liệu đầu vào (Input\newline Data) & Lần 1: nhập mã đơn không tồn tại.\newline Lần 2: dùng mã đơn của Bình.\newline Thao tác: xác nhận nhận xe. \\
\hline
Kết quả mong đợi\newline (Expected Result) & 1. Từ chối cả hai yêu cầu.\newline 2. Không tiết lộ chi tiết đơn của Bình.\newline 3. Khôngđổi trạngthái xe hoặc đơn. \\
\hline
Kếtquả thực tế / Status & Pending– Chưa thực hiện kiểm thử. \\
\hline
\end{longtable}
\vspace{-0.5\baselineskip}
\subsubsection{TCUC0503 – Kiểm tra thời hạn nhận xe}
\begin{longtable}{|L{0.24\textwidth}|L{0.70\textwidth}|}
\hline
\textbf{Trường thông tin} & \textbf{Đặc tả kịch bản} \\
\hline
\endfirsthead
\hline
\textbf{Trường thông tin} & \textbf{Đặc tả kịch bản} \\
\hline
\endhead
Test Case ID & TC\_UC05\_03 \\
\hline
Use Case liênquan & UC5 – Nhậnphươngtiện \\
\hline
Mô tả kịch bản & Kiểm tra thời hạn nhận xe \\
\hline
Tiền điều kiện\newline (Preconditions) & Đơn của An có hiệu lực từ 08:00 đến 08:15; xe không có sự cố. \\
\hline
Dữ liệu đầu vào (Input\newline Data) & Thử độc lập tại 08:14:59, 08:15:00 và 08:15:01.\newline Thao tác: xác nhận nhận xe tại mỗi thời điểm. \\
\hline
Kết quả mong đợi\newline (Expected Result) & 1. Trước 08:15 được nhận xe.\newline 2. Từ đúng 08:15 trở đi, thông báo đơn hết hiệu lực và không ghi\newline nhận nhận xe.Quyước biên thờigian cần thốngnhất trongnhóm. \\
\hline
Kếtquả thực tế / Status & Pending– Chưa thực hiện kiểm thử. \\
\hline
\end{longtable}
\vspace{-0.5\baselineskip}
\subsubsection{TCUC0504 – Lỗi cập nhật sau thao tác nhận}
\begin{longtable}{|L{0.24\textwidth}|L{0.70\textwidth}|}
\hline
\textbf{Trường thông tin} & \textbf{Đặc tả kịch bản} \\
\hline
\endfirsthead
\hline
\textbf{Trường thông tin} & \textbf{Đặc tả kịch bản} \\
\hline
\endhead
Test Case ID & TC\_UC05\_04 \\
\hline
Use Case liênquan & UC5 – Nhậnphươngtiện \\
\hline
Mô tả kịch bản & Lỗi cậpnhật sau thao tác nhận \\
\hline
Tiền điều kiện\newline (Preconditions) & Đã lưu thông tin nhận xe mô phỏng; gây lỗi ở bước cập nhật cuối. \\
\hline
Dữ liệu đầu vào (Input\newline Data) & An xác nhận nhận xe điện số 1.\newline Sau lỗi, Bình thử đặt lại chính xe đó. \\
\hline
Kết quả mong đợi\newline (Expected Result) & 1. Không báo nhận xe hoàn tất.\newline 2. Giữ thông tin nhận và chặn Bình đặt lại xe.\newline 3. Chuyển vận hành kiểm tra, khôngtựgiảiphóngxe. \\
\hline
Kếtquả thực tế / Status & Pending– Chưa thực hiện kiểm thử. \\
\hline
\end{longtable}
\vspace{-0.5\baselineskip}
\subsection{Test case cho UC6 – Trả phương tiện}
\subsubsection{TCUC0601 – Thành công: trả xe pin thấp, xe chuyển sang "Chờ sạc".}
\begin{longtable}{|L{0.24\textwidth}|L{0.70\textwidth}|}
\hline
\textbf{Trường thông tin} & \textbf{Đặc tả kịch bản} \\
\hline
\endfirsthead
\hline
\textbf{Trường thông tin} & \textbf{Đặc tả kịch bản} \\
\hline
\endhead
Test Case ID & TC\_UC06\_01 \\
\hline
Use Case liênquan & UC6: Trảphươngtiện \\
\hline
Mô tả kịch bản & Thành công: trả xepin thấp, xe chuyển sang"Chờ sạc". \\
\hline
Tiền điều kiện\newline (Preconditions) & SV01 đã đăng nhập. Xe V001 có lượt mượn IN\_USE của SV01, pin\newline 15\%, khônglỗi kỹthuật. Hub A còn chỗ tiếpnhận. \\
\hline
Dữ liệu đầu vào (Input\newline Data) & studentId = SV01, vehicleId = V001, hubId = HubA \\
\hline
Kết quả mong đợi\newline (Expected Result) & ReturnRecord.resultingStatus = WAITING\_CHARGE.\newline RentalRecord = ENDED. AuditLog ghi thời điểm, người trả, xe và\newline Hub nhận. Hiển thị xác nhận trả xe thành công. \\
\hline
Kết quả thực tế (Actual / & Pending \\
\hline
Status) & 2.6. Test case cho UC6 - Trảphươngtiện \\
\hline
\end{longtable}
\vspace{-0.5\baselineskip}
\subsubsection{TCUC0602 – Ngoại lệ: xe không có lượt mượn hợp lệ.}
\begin{longtable}{|L{0.24\textwidth}|L{0.70\textwidth}|}
\hline
\textbf{Trường thông tin} & \textbf{Đặc tả kịch bản} \\
\hline
\endfirsthead
\hline
\textbf{Trường thông tin} & \textbf{Đặc tả kịch bản} \\
\hline
\endhead
Test Case ID & TC\_UC06\_02 \\
\hline
Use Case liênquan & UC6: Trảphươngtiện \\
\hline
Mô tả kịch bản & Ngoại lệ: xe khôngcó lượt mượn hợp lệ. \\
\hline
Tiền điều kiện\newline (Preconditions) & SV01 đã đăng nhập. Xe V002 ở trạng thái AVAILABLE, không có\newline RentalRecord IN\_USE. \\
\hline
Dữ liệu đầu vào (Input\newline Data) & studentId = SV01, vehicleId = V002, hubId = HubA \\
\hline
Kết quả mong đợi\newline (Expected Result) & Ném InvalidReturnException. Không có ReturnRecord, AuditLog\newline hay thay đổi trạng thái xe. Hiển thị thông báo xe/lượt mượn không\newline hợp lệ và hướngdẫn liên hệ đơn vị vận hành. \\
\hline
Kết quả thực tế (Actual /\newline Status) & Pending \\
\hline
\end{longtable}
\vspace{-0.5\baselineskip}
\subsubsection{TCUC0603 – Ngoại lệ: gửi thông báo thất bại nhưng việc trả xe vẫn thành công.}
\begin{longtable}{|L{0.24\textwidth}|L{0.70\textwidth}|}
\hline
\textbf{Trường thông tin} & \textbf{Đặc tả kịch bản} \\
\hline
\endfirsthead
\hline
\textbf{Trường thông tin} & \textbf{Đặc tả kịch bản} \\
\hline
\endhead
Test Case ID & TC\_UC06\_03 \\
\hline
Use Case liênquan & UC6: Trảphươngtiện \\
\hline
Mô tả kịch bản & Ngoại lệ:gửi thôngbáo thất bại nhưngviệc trả xe vẫn thành công. \\
\hline
Tiền điều kiện\newline (Preconditions) & Như TC\_UC06\_01. NotificationService được giả lập trả lỗi. \\
\hline
Dữ liệu đầu vào (Input\newline Data) & studentId = SV01, vehicleId = V001, hubId = HubA \\
\hline
Kết quả mong đợi\newline (Expected Result) & returnVehicle không ném ngoại lệ. RentalRecord = ENDED, trạng\newline thái xe đã cập nhật, ReturnRecord đã lưu. Lỗi thông báo được ghi\newline log để xử lý lại. Sinh viên vẫn thấy xác nhận trả xe thành công.\newline 2.7. Test case cho UC7 - Đăngkýsạc \\
\hline
\end{longtable}
\vspace{-0.5\baselineskip}
\subsection{Test case cho UC7 – Đăng ký sạc}
\subsubsection{TCUC0701 – Ngoại lệ: khung giờ đã được người khác đặt.}
\begin{longtable}{|L{0.24\textwidth}|L{0.70\textwidth}|}
\hline
\textbf{Trường thông tin} & \textbf{Đặc tả kịch bản} \\
\hline
\endfirsthead
\hline
\textbf{Trường thông tin} & \textbf{Đặc tả kịch bản} \\
\hline
\endhead
Test Case ID & TC\_UC07\_01 \\
\hline
Use Case liênquan & UC7: Đăngkýsạc \\
\hline
Mô tả kịch bản & Ngoại lệ: khung giờ đã được người khác đặt. \\
\hline
Tiền điều kiện\newline (Preconditions) & Cổng P1 AVAILABLE. SV02 đã có đăng ký CONFIRMED cho P1,\newline 14:00 đến 15:00. \\
\hline
Dữ liệu đầu vào (Input\newline Data) & studentId = SV01, vehicleId = V001, portId = P1, start = 14:00, end\newline = 15:00 \\
\hline
Kết quả mong đợi\newline (Expected Result) & Ném SlotConflictException. Không tạo ChargingRegistration mới,\newline không ghi AuditLog thành công. Hiển thị xung đột và gợi ý cổng\newline hoặc khung giờ khác. P1 chỉ có đúng1 đăngkýtrongkhung giờ này. \\
\hline
Kết quả thực tế (Actual /\newline Status) & Pending \\
\hline
\end{longtable}
\vspace{-0.5\baselineskip}
\subsubsection{TCUC0702 – Race condition: hai sinh viên đăng ký cùng cổng, cùng khung giờ đồng thời.}
\begin{longtable}{|L{0.24\textwidth}|L{0.70\textwidth}|}
\hline
\textbf{Trường thông tin} & \textbf{Đặc tả kịch bản} \\
\hline
\endfirsthead
\hline
\textbf{Trường thông tin} & \textbf{Đặc tả kịch bản} \\
\hline
\endhead
Test Case ID & TC\_UC07\_02 \\
\hline
Use Case liênquan & UC7: Đăngkýsạc \\
\hline
Mô tả kịch bản & Race condition: hai sinh viên đăng ký cùng cổng, cùng khung giờ\newline đồngthời. \\
\hline
Tiền điều kiện\newline (Preconditions) & P1 AVAILABLE, chưa có đăng ký nào. Hai luồng (thread) bắt đầu\newline cùnglúc, mỗi luồngđãqua bước kiểm tra khả dụng. \\
\hline
Dữ liệu đầu vào (Input\newline Data) & Luồng A: SV01, V001, P1, 14:00, 15:00. Luồng B: SV02, V002, P1,\newline 14:00, 15:00. \\
\hline
Kết quả mong đợi\newline (Expected Result) & Đúng một luồng nhận ChargingRegistration CONFIRMED, luồng\newline còn lại nhận SlotConflictException. CSDL chỉ có 1 đăng ký cho P1\newline trong khung giờ này, P1 chỉ bị giữ một lần. Lặp lại 100 lần vẫn\newline khôngcó lần nào tạo 2 đăngký. \\
\hline
Kết quả thực tế (Actual /\newline Status) & Pending \\
\hline
\end{longtable}
\vspace{-0.5\baselineskip}
\subsubsection{TCUC0703 – Thành công: đăng ký cổng còn trống.}
\begin{longtable}{|L{0.24\textwidth}|L{0.70\textwidth}|}
\hline
\textbf{Trường thông tin} & \textbf{Đặc tả kịch bản} \\
\hline
\endfirsthead
\hline
\textbf{Trường thông tin} & \textbf{Đặc tả kịch bản} \\
\hline
\endhead
Test Case ID & TC\_UC07\_03 \\
\hline
Use Case liênquan & UC7: Đăngkýsạc \\
\hline
Mô tả kịch bản & Thành công: đăngkýcổngcòn trống. \\
\hline
Tiền điều kiện\newline (Preconditions) & P1 AVAILABLE, không có đăng ký chồng lấn. V001 thuộc lượt\newline dùnghợp lệ của SV01. \\
\hline
Dữ liệu đầu vào (Input\newline Data) & studentId = SV01, vehicleId = V001, portId = P1, start = 16:00, end\newline = 17:00 \\
\hline
Kết quả mong đợi\newline (Expected Result) & Trả về ChargingRegistration (CONFIRMED). P1 được giữ cho\newline 16:00 đến 17:00. AuditLog được ghi. Yêu cầu thông báo được\newline chuyển đến Hệ thống thông báo. Hiển thị xác nhận đăng ký.\newline 2.8. Test case cho UC8 - Giám sát mạnglưới MobilityHub \\
\hline
Kết quả thực tế (Actual /\newline Status) & Pending \\
\hline
\end{longtable}
\vspace{-0.5\baselineskip}
\subsection{Test case cho UC8 – Giám sát mạng lưới Mobility Hub}
\subsubsection{TCUC0801 – Biên (EX4) - Một Hub có dữ liệu cấu hình không hợp lệ.}
\begin{longtable}{|L{0.24\textwidth}|L{0.70\textwidth}|}
\hline
\textbf{Trường thông tin} & \textbf{Đặc tả kịch bản} \\
\hline
\endfirsthead
\hline
\textbf{Trường thông tin} & \textbf{Đặc tả kịch bản} \\
\hline
\endhead
Test Case ID & TC\_UC08\_01 \\
\hline
Use Case liênquan & UC8 - Giám sát mạnglưới MobilityHub \\
\hline
Mô tả kịch bản & Biên(EX4)- Một Hub có dữ liệu cấu hình khônghợp lệ. \\
\hline
Tiền điều kiện\newline (Preconditions) & Đã đăng nhập với vai trò Đơn vị vận hành; tồn tại 1 Hub thiếu thuộc\newline tính capacitytrongcấu hình. \\
\hline
Dữ liệu đầu vào (Input\newline Data) & filter=null (lấy toàn bộ mạng lưới) \\
\hline
\end{longtable}
\vspace{-0.5\baselineskip}
\begin{longtable}{|L{0.24\textwidth}|L{0.70\textwidth}|}
\hline
\textbf{Trường thông tin} & \textbf{Đặc tả kịch bản} \\
\hline
\endfirsthead
\hline
\textbf{Trường thông tin} & \textbf{Đặc tả kịch bản} \\
\hline
\endhead
Kết quả mong đợi & Dashboard vẫn hiển thị đúng các Hub hợp lệ; Hub lỗi cấu hình bị \\
\hline
(Expected Result) & đánh dấu "cần kiểm tra" và bị loại khỏiphéptính mức độ sử dụng. \\
\hline
Kết quả thực tế (Actual / & Pending \\
\hline
Status) & 2.9. Test case cho UC9 - Theo dõi tài nguyên \\
\hline
\end{longtable}
\vspace{-0.5\baselineskip}
\subsubsection{TCUC0802 – Ngoại lệ (AF5) - Không có Hub phù hợp với bộ lọc.}
\begin{longtable}{|L{0.24\textwidth}|L{0.70\textwidth}|}
\hline
\textbf{Trường thông tin} & \textbf{Đặc tả kịch bản} \\
\hline
\endfirsthead
\hline
\textbf{Trường thông tin} & \textbf{Đặc tả kịch bản} \\
\hline
\endhead
Test Case ID & TC\_UC08\_02 \\
\hline
Use Case liênquan & UC8 - Giám sát mạnglưới MobilityHub \\
\hline
Mô tả kịch bản & Ngoại lệ(AF5)- Khôngcó Hubphù hợpvới bộ lọc. \\
\hline
Tiền điều kiện\newline (Preconditions) & Đã đăng nhập với vai trò Đơn vị vận hành; không có Hub nào thuộc\newline khu vực được lọc. \\
\hline
Dữ liệu đầu vào (Input\newline Data) & filter = \{ khuVuc: "Khu vực không tồn tại" \} \\
\hline
Kết quả mong đợi\newline (Expected Result) & Hệ thống thông báo không có Hub đáp ứng tiêu chí, giữ nguyên bộ\newline lọc để Đơn vị vận hành điều chỉnh hoặc xoá bộ lọc. \\
\hline
Kết quả thực tế (Actual /\newline Status) & Pending \\
\hline
\end{longtable}
\vspace{-0.5\baselineskip}
\subsubsection{TCUC0803 – Ngoại lệ (EX1) - Người dùng không đủ quyền truy cập.}
\begin{longtable}{|L{0.24\textwidth}|L{0.70\textwidth}|}
\hline
\textbf{Trường thông tin} & \textbf{Đặc tả kịch bản} \\
\hline
\endfirsthead
\hline
\textbf{Trường thông tin} & \textbf{Đặc tả kịch bản} \\
\hline
\endhead
Test Case ID & TC\_UC08\_03 \\
\hline
Use Case liênquan & UC8 - Giám sát mạnglưới MobilityHub \\
\hline
Mô tả kịch bản & Ngoại lệ(EX1)- Người dùngkhôngđủquyền truycập. \\
\hline
Tiền điều kiện\newline (Preconditions) & Người dùng đăng nhập bằng tài khoản có vai trò Sinh viên. \\
\hline
Dữ liệu đầu vào (Input\newline Data) & session.role=STUDENT; filter=null \\
\hline
Kết quả mong đợi\newline (Expected Result) & Hệ thống ném UnauthorizedException, từ chối truy cập dashboard\newline vàyêu cầu xác thực lại bằngtài khoản cóquyền vận hành(NFR10). \\
\hline
Kết quả thực tế (Actual /\newline Status) & Pending \\
\hline
\end{longtable}
\vspace{-0.5\baselineskip}
\subsubsection{TCUC0804 – Ngoại lệ (EX2) - Không truy xuất được dữ liệu giám sát.}
\begin{longtable}{|L{0.24\textwidth}|L{0.70\textwidth}|}
\hline
\textbf{Trường thông tin} & \textbf{Đặc tả kịch bản} \\
\hline
\endfirsthead
\hline
\textbf{Trường thông tin} & \textbf{Đặc tả kịch bản} \\
\hline
\endhead
Test Case ID & TC\_UC08\_04 \\
\hline
Use Case liênquan & UC8 - Giám sát mạnglưới MobilityHub \\
\hline
Mô tả kịch bản & Ngoại lệ(EX2)- Khôngtruy xuất được dữ liệugiám sát. \\
\hline
Tiền điều kiện\newline (Preconditions) & Đã đăng nhập với vai trò Đơn vị vận hành; kết nối tới cơ sở dữ liệu\newline bịgián đoạn. \\
\hline
Dữ liệu đầu vào (Input\newline Data) & filter=null \\
\hline
\end{longtable}
\vspace{-0.5\baselineskip}
\begin{longtable}{|L{0.24\textwidth}|L{0.70\textwidth}|}
\hline
\textbf{Trường thông tin} & \textbf{Đặc tả kịch bản} \\
\hline
\endfirsthead
\hline
\textbf{Trường thông tin} & \textbf{Đặc tả kịch bản} \\
\hline
\endhead
Kết quả mong đợi & Hệ thống ném DataSourceUnavailableException, thông báo không \\
\hline
(Expected Result) & thể tải dashboard và cho phépĐơn vị vận hành thử lại. \\
\hline
Kết quả thực tế (Actual /\newline Status) & Pending \\
\hline
\end{longtable}
\vspace{-0.5\baselineskip}
\subsubsection{TCUC0805 – Ngoại lệ (EX3, NFR05) - Một Hub mất kết nối, các Hub khác vẫn hoạt động.}
\begin{longtable}{|L{0.24\textwidth}|L{0.70\textwidth}|}
\hline
\textbf{Trường thông tin} & \textbf{Đặc tả kịch bản} \\
\hline
\endfirsthead
\hline
\textbf{Trường thông tin} & \textbf{Đặc tả kịch bản} \\
\hline
\endhead
Test Case ID & TC\_UC08\_05 \\
\hline
Use Case liênquan & UC8 - Giám sát mạnglưới MobilityHub \\
\hline
Mô tả kịch bản & Ngoại lệ (EX3, NFR05) - Một Hub mất kết nối, các Hub khác vẫn\newline hoạt động. \\
\hline
Tiền điều kiện\newline (Preconditions) & Đã đăng nhập với vai trò Đơn vị vận hành; 1 Hub mất kết nối hoàn\newline toàn, các Hub còn lại bình thường. \\
\hline
Dữ liệu đầu vào (Input\newline Data) & filter=null; Hub "HUB-003" không gửi dữ liệu trong 10 phút \\
\hline
Kết quả mong đợi\newline (Expected Result) & Hub mất kết nối được hiển thị với dữ liệu gần nhất và đánh dấu chưa\newline xác định, không khẳng định tài nguyên còn khả dụng; dashboard của\newline các Hub còn lại vẫn hiển thị bình thường, khôngbịgián đoạn. \\
\hline
Kết quả thực tế (Actual /\newline Status) & Pending\newline 2.5. Test case cho UC5 - Nhậnphươngtiện \\
\hline
\end{longtable}
\vspace{-0.5\baselineskip}
\subsection{Test case cho UC9 – Theo dõi tài nguyên}
\subsubsection{TCUC0901 – Biên: một dòng tài nguyên có timestamp quá cũ.}
\begin{longtable}{|L{0.24\textwidth}|L{0.70\textwidth}|}
\hline
\textbf{Trường thông tin} & \textbf{Đặc tả kịch bản} \\
\hline
\endfirsthead
\hline
\textbf{Trường thông tin} & \textbf{Đặc tả kịch bản} \\
\hline
\endhead
Test Case ID & TC\_UC09\_01 \\
\hline
Use Case liênquan & UC9: Theo dõi tài nguyên \\
\hline
Mô tả kịch bản & Biên: một dòngtài nguyên có timestamp quá cũ. \\
\hline
Tiền điều kiện\newline (Preconditions) & Vận hành đã đăng nhập từ UC8. Hub A: 20 xe, 24 chỗ đỗ, 5 cổng\newline sạc khả dụng. Port 2 FAULTY, cậpnhật 07:47; Hub cậpnhật 10:39. \\
\hline
Dữ liệu đầu vào (Input\newline Data) & hubId = HubA, giờ hiện tại 10:39 \\
\hline
Kết quả mong đợi\newline (Expected Result) & HubResource với 20 / 24 / 5 và lastUpdated = 10:39. Dòng Port 2 có\newline isStale = true và được đánh dấu trên giao diện. Không dữ liệu vận\newline hành nào bị thay đổi. \\
\hline
Kết quả thực tế (Actual /\newline Status) & Pending \\
\hline
\end{longtable}
\vspace{-0.5\baselineskip}
\subsubsection{TCUC0902 – Ngoại lệ: không truy xuất được dữ liệu (EX1).}
\begin{longtable}{|L{0.24\textwidth}|L{0.70\textwidth}|}
\hline
\textbf{Trường thông tin} & \textbf{Đặc tả kịch bản} \\
\hline
\endfirsthead
\hline
\textbf{Trường thông tin} & \textbf{Đặc tả kịch bản} \\
\hline
\endhead
Test Case ID & TC\_UC09\_02 \\
\hline
Use Case liênquan & UC9: Theo dõi tài nguyên \\
\hline
Mô tả kịch bản & Ngoại lệ: khôngtruy xuất được dữ liệu(EX1). \\
\hline
\end{longtable}
\vspace{-0.5\baselineskip}
\begin{longtable}{|L{0.24\textwidth}|L{0.70\textwidth}|}
\hline
\textbf{Trường thông tin} & \textbf{Đặc tả kịch bản} \\
\hline
\endfirsthead
\hline
\textbf{Trường thông tin} & \textbf{Đặc tả kịch bản} \\
\hline
\endhead
Tiền điều kiện\newline (Preconditions) & Vận hành đã đăng nhập. Kho dữ liệu tài nguyên được giả lập lỗi kết\newline nối. \\
\hline
Dữ liệu đầu vào (Input\newline Data) & hubId = HubA \\
\hline
Kết quả mong đợi\newline (Expected Result) & Ném DataUnavailableException. Giao diện hiển thị thông báo lỗi\newline truy cập dữ liệu và cho phép thử lại ("Refresh information"). Không\newline hiển thị số liệu cũ như thể còn hiệu lực. \\
\hline
Kết quả thực tế (Actual / & Pending \\
\hline
Status) & 2.10. Test case cho UC10 - Xem cảnh báoquá tải \\
\hline
\end{longtable}
\vspace{-0.5\baselineskip}
\subsection{Test case cho UC10 – Xem cảnh báo quá tải}
\subsubsection{TCUC1001 – Thành công: chỗ đỗ 95\% vượt ngưỡng 75\%, gọi hai lần không tạo cảnh báo trùng.}
\begin{longtable}{|L{0.24\textwidth}|L{0.70\textwidth}|}
\hline
\textbf{Trường thông tin} & \textbf{Đặc tả kịch bản} \\
\hline
\endfirsthead
\hline
\textbf{Trường thông tin} & \textbf{Đặc tả kịch bản} \\
\hline
\endhead
Test Case ID & TC\_UC10\_01 \\
\hline
Use Case liênquan & UC10: Xem cảnh báoquá tải \\
\hline
Mô tả kịch bản & Thành công: chỗ đỗ 95\% vượt ngưỡng 75\%, gọi hai lần không tạo\newline cảnh báo trùng. \\
\hline
Tiền điều kiện\newline (Preconditions) & Hub A có DataState = FRESH. Ngưỡng PARKING\_SPOT = 75\%.\newline Chưa có cảnh báo ACTIVE cho Hub A, loại PARKING\_SPOT. \\
\hline
Dữ liệu đầu vào (Input\newline Data) & hubId = HubA, mức sử dụng chỗ đỗ 95\%. Gọi evaluateHubLoad hai\newline lần liên tiếp. \\
\hline
Kết quả mong đợi\newline (Expected Result) & Lần 1: tạo 1 OverloadAlert (utilization = 0.95, threshold = 0.75,\newline severity = HIGH, ACTIVE). Lần 2: vẫn 1 cảnh báo ACTIVE, được\newline cập nhật chứ không tạo mới. Chi tiết hiển thị Hub, loại, mức sử\newline dụng, ngưỡng, vị trí, nguyên nhân dự kiến. \\
\hline
Kết quả thực tế (Actual /\newline Status) & Pending \\
\hline
\end{longtable}
\vspace{-0.5\baselineskip}
\subsubsection{TCUC1002 – Biên: mức sử dụng thấp hơn ngưỡng, không có cảnh báo.}
\begin{longtable}{|L{0.24\textwidth}|L{0.70\textwidth}|}
\hline
\textbf{Trường thông tin} & \textbf{Đặc tả kịch bản} \\
\hline
\endfirsthead
\hline
\textbf{Trường thông tin} & \textbf{Đặc tả kịch bản} \\
\hline
\endhead
Test Case ID & TC\_UC10\_02 \\
\hline
Use Case liênquan & UC10: Xem cảnh báoquá tải \\
\hline
Mô tả kịch bản & Biên: mức sử dụngthấphơn ngưỡng, khôngcó cảnh báo. \\
\hline
Tiền điều kiện\newline (Preconditions) & Hub C có DataState = FRESH. Ngưỡng 75\% cho chỗ đỗ và cổng\newline sạc. Chưa có cảnh báo ACTIVE. \\
\hline
Dữ liệu đầu vào (Input\newline Data) & hubId = HubC, chỗ đỗ 70\%, cổng sạc 60\% \\
\hline
Kết quả mong đợi\newline (Expected Result) & Trả về danh sách rỗng, không tạo OverloadAlert nào. Giao diện gọi\newline showNormalStatus() (trạngthái mạnglưới bình thường). \\
\hline
Kết quả thực tế (Actual /\newline Status) & Pending \\
\hline
\end{longtable}
\vspace{-0.5\baselineskip}
\subsubsection{TCUC1003 – Ngoại lệ: dữ liệu quá cũ, không tạo cảnh báo.}
\begin{longtable}{|L{0.24\textwidth}|L{0.70\textwidth}|}
\hline
\textbf{Trường thông tin} & \textbf{Đặc tả kịch bản} \\
\hline
\endfirsthead
\hline
\textbf{Trường thông tin} & \textbf{Đặc tả kịch bản} \\
\hline
\endhead
Test Case ID & TC\_UC10\_03 \\
\hline
Use Case liênquan & UC10: Xem cảnh báoquá tải \\
\hline
Mô tả kịch bản & Ngoại lệ: dữ liệuquá cũ, khôngtạo cảnh báo. \\
\hline
Tiền điều kiện\newline (Preconditions) & Hub A có DataState = STALE (dữ liệu cũ hơn 15 giây). Ngưỡng đã\newline cấu hình. \\
\hline
Dữ liệu đầu vào (Input\newline Data) & hubId = HubA \\
\hline
Kết quả mong đợi\newline (Expected Result) & Ném DataUnavailableException. Không tạo hay cập nhật cảnh báo\newline nào, tránh cảnh báo sai do dữ liệu lỗi thời. Giao diện hiển thị thông\newline báo lỗi vàgiữ nguyên dashboard để thử lại. \\
\hline
Kết quả thực tế (Actual /\newline Status) & Pending\newline 2.11. Test case cho UC11 - Điềuphốiphươngtiện \\
\hline
\end{longtable}
\vspace{-0.5\baselineskip}
\subsection{Test case cho UC11 – Điều phối phương tiện}
\subsubsection{TCUC1101 – Thành công (Normal Flow): Điều phối 2 xe từ Hub KTX B sang Hub Ga Metro còn đủ chỗ đỗ trống.}
\begin{longtable}{|L{0.24\textwidth}|L{0.70\textwidth}|}
\hline
\textbf{Trường thông tin} & \textbf{Đặc tả kịch bản} \\
\hline
\endfirsthead
\hline
\textbf{Trường thông tin} & \textbf{Đặc tả kịch bản} \\
\hline
\endhead
Test Case ID & TC\_UC11\_01 \\
\hline
Use Case liênquan & UC11 – Điềuphốiphươngtiện \\
\hline
Mô tả kịch bản & Thành công (Normal Flow): Điều phối 2 xe từ Hub KTX B sang\newline Hub Ga Metro còn đủ chỗ đỗ trống. \\
\hline
Tiền điều kiện\newline (Preconditions) & 1. Operator OP01 đã đăng nhập hợp lệ qua SSO (NFR10).\newline 2. Hub KTX B (HUB\_KTXB) có xe EB\_01 (AVAILABLE, pin\newline 90\%) và EB\_02 (AVAILABLE, pin 85\%).\newline 3. Hub Ga Metro (HUB\_METRO) đang ACTIVE và còn\newline availableSpots = 10. \\
\hline
Dữ liệu đầu vào (Input\newline Data) & operatorId = "OP01", sourceHubId = "HUB\_KTXB", destHubId =\newline "HUB\_METRO", vehicleIds =["EB\_01", "EB\_02"] \\
\hline
Kết quả mong đợi\newline (Expected Result) & 1. Tạo thành công DispatchOrder (status = PENDING) với 2 xe\newline EB\_01, EB\_02.\newline 2. EB\_01 và EB\_02 chuyển sang trạng thái DISPATCHING (Sinh\newline viên không thể đặt – NFR09).\newline 3. availableSpots tại HUB\_METRO giảm từ 10 xuống 8; tại\newline HUB\_KTXB tăng tương ứng.\newline 4. Ghi 1 bản ghi AuditLog (NFR11) và gửi thông báo công việc tới\newline đội vận chuyểnqua NotificationService. \\
\hline
Kết quả thực tế (Actual /\newline Status) & Pending \\
\hline
\end{longtable}
\vspace{-0.5\baselineskip}
\subsubsection{TCUC1102 – Ngoại lệ (EX1 – Race Condition / NFR09): Xe được chọn vừa bị Sinh viên đặt trước ngay trước khi Operator xác nhận lệnh.}
\begin{longtable}{|L{0.24\textwidth}|L{0.70\textwidth}|}
\hline
\textbf{Trường thông tin} & \textbf{Đặc tả kịch bản} \\
\hline
\endfirsthead
\hline
\textbf{Trường thông tin} & \textbf{Đặc tả kịch bản} \\
\hline
\endhead
Test Case ID & TC\_UC11\_02 \\
\hline
Use Case liênquan & UC11 – Điềuphốiphươngtiện \\
\hline
\end{longtable}
\vspace{-0.5\baselineskip}
\begin{longtable}{|L{0.24\textwidth}|L{0.70\textwidth}|}
\hline
\textbf{Trường thông tin} & \textbf{Đặc tả kịch bản} \\
\hline
\endfirsthead
\hline
\textbf{Trường thông tin} & \textbf{Đặc tả kịch bản} \\
\hline
\endhead
Mô tả kịch bản & Ngoại lệ (EX1 – Race Condition / NFR09): Xe được chọn vừa bị\newline Sinh viên đặt trước ngaytrước khi Operator xác nhận lệnh. \\
\hline
Tiền điều kiện\newline (Preconditions) & 1. Operator OP01 mở màn hình điều phối và thấy xe EB\_01 khả\newline dụng.\newline 2. Ngay trước thời điểm Operator bấm xác nhận, Sinh viên SV05 đặt\newline thành côngEB\_01(status chuyển sangRESERVED). \\
\hline
Dữ liệu đầu vào (Input\newline Data) & operatorId = "OP01", sourceHubId = "HUB\_KTXB", destHubId =\newline "HUB\_METRO", vehicleIds =["EB\_01", "EB\_02"] \\
\hline
Kết quả mong đợi\newline (Expected Result) & 1. Cơ chế khóa giao dịch phát hiện EB\_01 không còn ở trạng thái\newline AVAILABLE và ném VehicleConflictException.\newline 2. Toàn bộ giao dịch bị rollback: không tạo DispatchOrder nào,\newline EB\_02 vẫn giữ nguyên AVAILABLE, sức chứa các Hub không đổi.\newline 3. Giao diện thông báo xung đột ở EB\_01 và gợi ý Operator loại bỏ\newline xe xungđột hoặc chọn xe khác. \\
\hline
Kết quả thực tế (Actual /\newline Status) & Pending \\
\hline
\end{longtable}
\vspace{-0.5\baselineskip}
\subsubsection{TCUC1103 – Ngoại lệ (EX2): Hub đích không đủ chỗ đỗ trống cho số lượng phương tiện điều phối.}
\begin{longtable}{|L{0.24\textwidth}|L{0.70\textwidth}|}
\hline
\textbf{Trường thông tin} & \textbf{Đặc tả kịch bản} \\
\hline
\endfirsthead
\hline
\textbf{Trường thông tin} & \textbf{Đặc tả kịch bản} \\
\hline
\endhead
Test Case ID & TC\_UC11\_03 \\
\hline
Use Case liênquan & UC11 – Điềuphốiphươngtiện \\
\hline
Mô tả kịch bản & Ngoại lệ (EX2): Hub đích không đủ chỗ đỗ trống cho số lượng\newline phươngtiện điềuphối. \\
\hline
Tiền điều kiện\newline (Preconditions) & 1. Operator OP01 đã đăng nhập.\newline 2. Hub nguồn HUB\_KTXB có 3 xe AVAILABLE.\newline 3. Hub đích HUB\_THUVIEN đang ACTIVE nhưng chỉ còn\newline availableSpots = 1. \\
\hline
Dữ liệu đầu vào (Input\newline Data) & operatorId = "OP01", sourceHubId = "HUB\_KTXB", destHubId =\newline "HUB\_THUVIEN", vehicleIds =["EB\_01", "EB\_02", "EB\_03"] \\
\hline
Kết quả mong đợi\newline (Expected Result) & 1. hasAvailableCapacity(3) trả về false -> ném\newline HubCapacityExceededException.\newline 2. Không tạo DispatchOrder, trạng thái cả 3 xe vẫn là AVAILABLE.\newline 3. Giao diện hiển thị cảnh báo quá tải sức chứa Hub đích và yêu cầu\newline giảm số xe(<= 1)hoặc chọn Hub đích khác. \\
\hline
Kết quả thực tế (Actual /\newline Status) & Pending \\
\hline
\end{longtable}
\vspace{-0.5\baselineskip}
\subsubsection{TCUC1104 – Luồng thay thế (AF2): Hủy lệnh điều phối đang chờ xử lý trước khi vận chuyển.}
\begin{longtable}{|L{0.24\textwidth}|L{0.70\textwidth}|}
\hline
\textbf{Trường thông tin} & \textbf{Đặc tả kịch bản} \\
\hline
\endfirsthead
\hline
\textbf{Trường thông tin} & \textbf{Đặc tả kịch bản} \\
\hline
\endhead
Test Case ID & TC\_UC11\_04 \\
\hline
Use Case liênquan & UC11 – Điềuphốiphươngtiện \\
\hline
Mô tả kịch bản & Luồng thay thế (AF2): Hủy lệnh điều phối đang chờ xử lý trước khi\newline vận chuyển. \\
\hline
\end{longtable}
\vspace{-0.5\baselineskip}
\begin{longtable}{|L{0.24\textwidth}|L{0.70\textwidth}|}
\hline
\textbf{Trường thông tin} & \textbf{Đặc tả kịch bản} \\
\hline
\endfirsthead
\hline
\textbf{Trường thông tin} & \textbf{Đặc tả kịch bản} \\
\hline
\endhead
Tiền điều kiện\newline (Preconditions) & 1. Lệnh điều phối ORD\_101 đang ở trạng thái PENDING gồm 2 xe\newline EB\_01, EB\_02 (status = DISPATCHING).\newline 2. Đội vận chuyển chưa bắt đầu di chuyển xe. \\
\hline
Dữ liệu đầu vào (Input\newline Data) & operatorId = "OP01", orderId = "ORD\_101" \\
\hline
Kết quả mong đợi\newline (Expected Result) & 1. cancelDispatchOrder() trả về true.\newline 2. ORD\_101 chuyển sang status = CANCELLED.\newline 3. EB\_01 và EB\_02 được mở khóa và khôi phục về AVAILABLE\newline tại HUB\_KTXB.\newline 4. Sức chứa dự kiến của HUB\_KTXB và HUB\_METRO được hoàn\newline tác vàghi nhận vào AuditLog. \\
\hline
Kết quả thực tế (Actual /\newline Status) & Pending \\
\hline
\end{longtable}
\vspace{-0.5\baselineskip}
\subsection{Test case cho UC12 – Xử lý các sự cố}
\subsubsection{TCUC1201 – Thành công (Normal Flow): Tạm ngưng cổng sạc bị lỗi đang có lịch đặt trước của Sinh viên, hủy lịch, thông báo và ghi AuditLog.}
\begin{longtable}{|L{0.24\textwidth}|L{0.70\textwidth}|}
\hline
\textbf{Trường thông tin} & \textbf{Đặc tả kịch bản} \\
\hline
\endfirsthead
\hline
\textbf{Trường thông tin} & \textbf{Đặc tả kịch bản} \\
\hline
\endhead
Test Case ID & TC\_UC12\_01 \\
\hline
Use Case liênquan & UC12 – Xử lýcác sự cố \\
\hline
Mô tả kịch bản & Thành công (Normal Flow): Tạm ngưng cổng sạc bị lỗi đang có lịch\newline đặt trước của Sinh viên, hủylịch, thôngbáo vàghi AuditLog. \\
\hline
Tiền điều kiện\newline (Preconditions) & 1. Operator OP01 đã đăng nhập hợp lệ.\newline 2. Sự cố INC\_201 (CRITICAL/HIGH) được ghi nhận cho cổng sạc\newline CP\_03 (status = AVAILABLE).\newline 3. Sinh viên SV12 đang có lịch ChargingRegistration CR\_55\newline (CONFIRMED)cho CP\_03. \\
\hline
Dữ liệu đầu vào (Input\newline Data) & operatorId = "OP01", incidentId = "INC\_201", reason = "Cổng sạc\newline mất kết nối bo mạch nguồn", fixingMethod = "Cử kỹ thuật viên thay\newline rơ-le bảo vệ" \\
\hline
Kết quả mong đợi\newline (Expected Result) & 1. Incident INC\_201 chuyển sang IN\_PROGRESS với đầy đủ\newline nguyên nhân và hướng khắc phục.\newline 2. Cổng CP\_03 chuyển sang MAINTENANCE (bị loại khỏi danh\newline sách tìm kiếm/đặt chỗ).\newline 3. ChargingRegistration CR\_55 tự động chuyển sang\newline CANCELLED\_BY\_INCIDENT.\newline 4. NotificationService gửi thông báo xin lỗi tới SV12 và ghi\newline AuditLog (lưu >= 90 ngày– NFR11). \\
\hline
Kết quả thực tế (Actual /\newline Status) & Pending \\
\hline
\end{longtable}
\vspace{-0.5\baselineskip}
\subsubsection{TCUC1202 – Luồng thay thế (AF1): Hoàn tất sửa chữa và khôi phục tài nguyên về trạng thái Khả dụng.}
\begin{longtable}{|L{0.24\textwidth}|L{0.70\textwidth}|}
\hline
\textbf{Trường thông tin} & \textbf{Đặc tả kịch bản} \\
\hline
\endfirsthead
\hline
\textbf{Trường thông tin} & \textbf{Đặc tả kịch bản} \\
\hline
\endhead
Test Case ID & TC\_UC12\_02 \\
\hline
Use Case liênquan & UC12 – Xử lýcác sự cố \\
\hline
\end{longtable}
\vspace{-0.5\baselineskip}
\begin{longtable}{|L{0.24\textwidth}|L{0.70\textwidth}|}
\hline
\textbf{Trường thông tin} & \textbf{Đặc tả kịch bản} \\
\hline
\endfirsthead
\hline
\textbf{Trường thông tin} & \textbf{Đặc tả kịch bản} \\
\hline
\endhead
Mô tả kịch bản & Luồng thay thế (AF1): Hoàn tất sửa chữa và khôi phục tài nguyên về\newline trạngthái Khả dụng. \\
\hline
Tiền điều kiện\newline (Preconditions) & 1. Sự cố INC\_201 đang ở trạng thái IN\_PROGRESS; cổng sạc\newline CP\_03 đangở trạngthái MAINTENANCE. \\
\hline
Dữ liệu đầu vào (Input\newline Data) & operatorId = "OP01", incidentId = "INC\_201" \\
\hline
Kết quả mong đợi\newline (Expected Result) & 1. Incident INC\_201 chuyển sang RESOLVED và ghi nhận\newline resolvedAt = thời điểm hiện tại.\newline 2. Cổng sạc CP\_03 khôi phục từ MAINTENANCE về\newline AVAILABLE.\newline 3. Ghi nhận AuditLog thao tác hoàn tất xử lý và cập nhật trạng thái\newline mạnglưới. \\
\hline
Kết quả thực tế (Actual /\newline Status) & Pending \\
\hline
\end{longtable}
\vspace{-0.5\baselineskip}
\subsubsection{TCUC1203 – Ngoại lệ (EX1 – NFR05 Fault Tolerance): Mất kết nối hoàn toàn tại một Hub đơn lẻ kích hoạt cô lập lỗi mà không làm gián đoạn các Hub khác.}
\begin{longtable}{|L{0.24\textwidth}|L{0.70\textwidth}|}
\hline
\textbf{Trường thông tin} & \textbf{Đặc tả kịch bản} \\
\hline
\endfirsthead
\hline
\textbf{Trường thông tin} & \textbf{Đặc tả kịch bản} \\
\hline
\endhead
Test Case ID & TC\_UC12\_03 \\
\hline
Use Case liênquan & UC12 – Xử lýcác sự cố \\
\hline
Mô tả kịch bản & Ngoại lệ (EX1 – NFR05 Fault Tolerance): Mất kết nối hoàn toàn tại\newline một Hub đơn lẻ kích hoạt cô lập lỗi mà không làm gián đoạn các\newline Hub khác. \\
\hline
Tiền điều kiện\newline (Preconditions) & 1. Mạng lưới có HUB\_A, HUB\_B, HUB\_C đang hoạt động bình\newline thường (ACTIVE).\newline 2. HUB\_B mất kết nối hoàn toàn khỏi hệ thống giám sát. \\
\hline
Dữ liệu đầu vào (Input\newline Data) & hubId = "HUB\_B" \\
\hline
Kết quả mong đợi\newline (Expected Result) & 1. HUB\_B chuyển sang trạng thái ISOLATED; mọi giao dịch mới\newline tại HUB\_B bị tạm ngưng.\newline 2. Tự động tạo sự cố loại HUB\_DISCONNECTED mức CRITICAL\newline và gửi cảnh báo cho Operator.\newline 3. HUB\_A và HUB\_C vẫn ở trạng thái ACTIVE, phục vụ tìm\newline kiếm/đặt xe bình thườngkhôngbị treo hay gián đoạn(NFR05). \\
\hline
Kết quả thực tế (Actual /\newline Status) & Pending \\
\hline
\end{longtable}
\vspace{-0.5\baselineskip}
\begin{longtable}{|L{0.24\textwidth}|L{0.70\textwidth}|}
\hline
\textbf{Trường thông tin} & \textbf{Đặc tả kịch bản} \\
\hline
\endfirsthead
\hline
\textbf{Trường thông tin} & \textbf{Đặc tả kịch bản} \\
\hline
\endhead
Test Case ID & TC\_UC12\_04 \\
\hline
Use Case liênquan & UC12 – Xử lýcác sự cố \\
\hline
Mô tả kịch bản & Ngoại lệ (Bảo mật RBAC – NFR10): Tài khoản Sinh viên cố gắng\newline gọi chức năngxử lýsự cố bị từ chối truycập. \\
\hline
Tiền điều kiện & 1. Người dùng đang đăng nhập với tài khoản vai trò STUDENT \\
\hline
(Preconditions) & (khôngcóquyền OPERATOR). \\
\hline
\end{longtable}
\vspace{-0.5\baselineskip}
\begin{longtable}{|L{0.24\textwidth}|L{0.70\textwidth}|}
\hline
\textbf{Trường thông tin} & \textbf{Đặc tả kịch bản} \\
\hline
\endfirsthead
\hline
\textbf{Trường thông tin} & \textbf{Đặc tả kịch bản} \\
\hline
\endhead
Dữ liệu đầu vào (Input\newline Data) & actorId = "SV01" (role = STUDENT), incidentId = "INC\_201" \\
\hline
Kết quả mong đợi\newline (Expected Result) & 1. Hệ thống ném UnauthorizedException ngay bước kiểm tra quyền.\newline 2. Không có bất kỳ thay đổi nào lên Incident, Resource hay\newline AuditLognghiệpvụ. \\
\hline
Kết quả thực tế (Actual / & Pending \\
\hline
Status) & 2.13. Test case cho UC13 - Lậplịch sạc thôngminh \\
\hline
\end{longtable}
\vspace{-0.5\baselineskip}
\subsection{Test case cho UC13 – Lập lịch sạc thông minh}
\subsubsection{TCUC1301 – Lập lịch sạc và phê duyệt thủ công thành công}
\begin{longtable}{|L{0.24\textwidth}|L{0.70\textwidth}|}
\hline
\textbf{Trường thông tin} & \textbf{Đặc tả kịch bản} \\
\hline
\endfirsthead
\hline
\textbf{Trường thông tin} & \textbf{Đặc tả kịch bản} \\
\hline
\endhead
Test Case ID & TC\_UC13\_01 \\
\hline
Use Case liênquan & UC13 – Smart ChargingSchedule \\
\hline
Mô tả kịch bản & Kiểm tra hệ thống tạo lịch sạc phù hợp và cho phép Operator phê\newline duyệt, lưu lịch thành côngkhi các tài nguyên sạc khả dụng. \\
\hline
Tiền điều kiện\newline (Preconditions) & 1. Operator đã đăng nhập và có quyền lập lịch sạc.\newline 2. Dữ liệu phương tiện, mức pin, nhu cầu sử dụng, thời gian chờ và\newline trạng thái trụ sạc đã được cập nhật.\newline 3. Hub còn khả năng tiếp nhận lịch sạc.\newline 4. Trụ sạc được chọn khôngcó lịch trùngthờigian. \\
\hline
Dữ liệu đầu vào (Input\newline Data) & Operator: OP001; Hub: HUB001;Vehicle: VEH001; Battery Level:\newline 15\%; Charging Port: CP001; Charging Time: 14:00–15:00;\newline Approval Mode: MANUAL. \\
\hline
Kết quả mong đợi\newline (Expected Result) & 1. Hệ thống tạo lịch sạc phù hợp với dữ liệu đầu vào.\newline 2. Lịch sạc hiển thị đầy đủ phương tiện, trụ sạc, thời gian và mức độ\newline ưu tiên.\newline 3. Operator phê duyệt lịch thành công.\newline 4. Lịch được lưu vào hệ thống.\newline 5. Trụ sạc được chuyển sang trạng thái đã giữ chỗ.\newline 6. Thao tác đượcghi nhận trongAuditLog. \\
\hline
Kếtquả thực tế / Status & Pending– Chưa thực hiện kiểm thử. \\
\hline
\end{longtable}
\vspace{-0.5\baselineskip}
\subsubsection{TCUC1302 – Operator chỉnh sửa lịch sạc trước khi phê duyệt}
\begin{longtable}{|L{0.24\textwidth}|L{0.70\textwidth}|}
\hline
\textbf{Trường thông tin} & \textbf{Đặc tả kịch bản} \\
\hline
\endfirsthead
\hline
\textbf{Trường thông tin} & \textbf{Đặc tả kịch bản} \\
\hline
\endhead
Test Case ID & TC\_UC13\_02 \\
\hline
Use Case liênquan & UC13 – Smart ChargingSchedule \\
\hline
Mô tả kịch bản & Kiểm tra khả năng Operator chỉnh sửa lịch sạc do hệ thống đề xuất\newline trước khiphê duyệt. \\
\hline
Tiền điều kiện\newline (Preconditions) & 1. Operator đã đăng nhập và có quyền lập lịch sạc.\newline 2. Hệ thống đã tạo lịch sạc ở trạng thái chờ phê duyệt.\newline 3. Hub và trụ sạc có đủ tài nguyên để đápứnglịch sau khi chỉnh sửa. \\
\hline
\end{longtable}
\vspace{-0.5\baselineskip}
\begin{longtable}{|L{0.24\textwidth}|L{0.70\textwidth}|}
\hline
\textbf{Trường thông tin} & \textbf{Đặc tả kịch bản} \\
\hline
\endfirsthead
\hline
\textbf{Trường thông tin} & \textbf{Đặc tả kịch bản} \\
\hline
\endhead
Dữ liệu đầu vào (Input\newline Data) & Operator: OP001\newline Schedule ID: SCH001\newline Vehicle: VEH001\newline Charging Port: CP001\newline Thời gian ban đầu: 14:00–15:00\newline Thờigian sau chỉnh sửa: 15:00–16:00 \\
\hline
Kết quả mong đợi\newline (Expected Result) & 1. Hệ thống hiển thị lịch sạc hiện tại để Operator chỉnh sửa.\newline 2. Operator cập nhật thời gian sạc thành công.\newline 3. Hệ thống kiểm tra và xác nhận khung giờ mới không bị xung đột\newline tài nguyên.\newline 4. Lịch sạc được cập nhật theo thông tin mới nhưng vẫn ở trạng thái\newline chờ phê duyệt.\newline 5. Operator có thể tiếptụcphê duyệt lịch sau khi chỉnh sửa. \\
\hline
Kếtquả thực tế / Status & Pending– Chưa thực hiện kiểm thử. \\
\hline
\end{longtable}
\vspace{-0.5\baselineskip}
\subsubsection{TCUC1303 – Hệ thống tự động phê duyệt lịch sạc}
\begin{longtable}{|L{0.24\textwidth}|L{0.70\textwidth}|}
\hline
\textbf{Trường thông tin} & \textbf{Đặc tả kịch bản} \\
\hline
\endfirsthead
\hline
\textbf{Trường thông tin} & \textbf{Đặc tả kịch bản} \\
\hline
\endhead
Test Case ID & TC\_UC13\_03 \\
\hline
Use Case liênquan & UC13 – Smart ChargingSchedule \\
\hline
Mô tả kịch bản & Kiểm tra khả năng hệ thống tự động phê duyệt lịch sạc khi chế độ\newline phê duyệt tự động được kích hoạt và lịch đáp ứng các điều kiện hợp\newline lệ. \\
\hline
Tiền điều kiện\newline (Preconditions) & 1. Operator đã đăng nhập và có quyền lập lịch sạc.\newline 2. Chế độ tự động phê duyệt đã được kích hoạt.\newline 3. Dữ liệu phương tiện, mức pin, nhu cầu sử dụng và tài nguyên sạc\newline hợp lệ.\newline 4. Trụ sạc khôngcó lịch trùngthờigian và Hub khôngbịquá tải. \\
\hline
Dữ liệu đầu vào (Input\newline Data) & Operator: OP001\newline Hub: HUB001\newline Vehicle: VEH002\newline Battery Level: 18\%\newline Charging Port: CP002\newline Charging Time: 16:00–17:00\newline Approval Mode: AUTO \\
\hline
Kết quả mong đợi\newline (Expected Result) & 1. Hệ thống tiếp nhận dữ liệu và tạo lịch sạc hợp lệ.\newline 2. Hệ thống tự động phê duyệt lịch mà không yêu cầu Operator xác\newline nhận thủ công.\newline 3. Lịch sạc được lưu thành công với trạng thái đã phê duyệt.\newline 4. Trụ sạc tương ứng được chuyển sang trạng thái đã giữ chỗ.\newline 5. Thao tácphê duyệt đượcghi nhận trongAuditLog. \\
\hline
Kếtquả thực tế / Status & Pending– Chưa thực hiện kiểm thử. \\
\hline
\end{longtable}
\vspace{-0.5\baselineskip}
\subsubsection{TCUC1304 – Xử lý xung đột tài nguyên và Hub quá tải}
\begin{longtable}{|L{0.24\textwidth}|L{0.70\textwidth}|}
\hline
\textbf{Trường thông tin} & \textbf{Đặc tả kịch bản} \\
\hline
\endfirsthead
\hline
\textbf{Trường thông tin} & \textbf{Đặc tả kịch bản} \\
\hline
\endhead
Test Case ID & TC\_UC13\_04 \\
\hline
Use Case liênquan & UC13 – Smart ChargingSchedule \\
\hline
\end{longtable}
\vspace{-0.5\baselineskip}
\begin{longtable}{|L{0.24\textwidth}|L{0.70\textwidth}|}
\hline
\textbf{Trường thông tin} & \textbf{Đặc tả kịch bản} \\
\hline
\endfirsthead
\hline
\textbf{Trường thông tin} & \textbf{Đặc tả kịch bản} \\
\hline
\endhead
Mô tả kịch bản & Kiểm tra khả năng hệ thống phát hiện xung đột trụ sạc, khung giờ\newline hoặc tình trạngHubquá tải trong quá trình lậplịch sạc. \\
\hline
Tiền điều kiện\newline (Preconditions) & 1. Operator đã đăng nhập và có quyền lập lịch sạc.\newline 2. Hub HUB001 đang hoạt động.\newline 3. Trụ sạc CP001 đã được giữ chỗ trong khoảng thời gian 14:00–\newline 15:00.\newline 4. Dữ liệu tải của Hub được cập nhật và có thể xác định tình trạng\newline quá tải. \\
\hline
Dữ liệu đầu vào (Input\newline Data) & Operator: OP001\newline Hub: HUB001\newline Vehicle: VEH003\newline Battery Level: 12\%\newline Charging Port: CP001\newline Charging Time: 14:30–15:30\newline Tình trạngHub:Quá tải trongtrườnghợpkiểm thử bổ sung. \\
\hline
Kết quả mong đợi\newline (Expected Result) & 1. Hệ thống phát hiện trụ sạc CP001 bị trùng lịch trong khoảng\newline 14:30–15:00.\newline 2. Hệ thống không cho phép lưu lịch sạc bị xung đột.\newline 3. Hệ thống hiển thị cảnh báo và đề xuất khung giờ hoặc tài nguyên\newline thay thế phù hợp.\newline 4. Khi Hub quá tải, hệ thống thông báo tình trạng quá tải và đề xuất\newline chuyển yêu cầu đến chức năng điều phối tại UC11.\newline 5. Không phát sinh lịch sạc mới hoặc giữ chỗ tài nguyên không hợp\newline lệ. \\
\hline
Kếtquả thực tế / Status & Pending– Chưa thực hiện kiểm thử. \\
\hline
\end{longtable}
\vspace{-0.5\baselineskip}
\subsubsection{TCUC1305 – Xử lý lỗi truy xuất dữ liệu và bảo đảm tính toàn vẹn khi lưu lịch}
\begin{longtable}{|L{0.24\textwidth}|L{0.70\textwidth}|}
\hline
\textbf{Trường thông tin} & \textbf{Đặc tả kịch bản} \\
\hline
\endfirsthead
\hline
\textbf{Trường thông tin} & \textbf{Đặc tả kịch bản} \\
\hline
\endhead
Test Case ID & TC\_UC13\_05 \\
\hline
Use Case liênquan & UC13 – Smart ChargingSchedule \\
\hline
Mô tả kịch bản & Kiểm tra khả năng hệ thống xử lý lỗi khi không truy xuất được dữ\newline liệu từ hệ thống mô phỏng hoặc cảm biến, đồng thời bảo đảm các dữ\newline liệu liên quan được lưu đồng bộ khi phát sinh lỗi trong quá trình lưu\newline lịch. \\
\hline
Tiền điều kiện\newline (Preconditions) & 1. Operator đã đăng nhập và có quyền lập lịch sạc.\newline 2. Hệ thống đang sử dụng dữ liệu từ hệ thống mô phỏng hoặc cảm\newline biến.\newline 3. Có thể mô phỏng tình huống lỗi truy xuất dữ liệu và lỗi trong quá\newline trình lưu lịch.\newline 4. Có sẵn lịch sạc và trụ sạc để kiểm tra trạng thái trước và sau giao\newline dịch. \\
\hline
\end{longtable}
\vspace{-0.5\baselineskip}
\begin{longtable}{|L{0.24\textwidth}|L{0.70\textwidth}|}
\hline
\textbf{Trường thông tin} & \textbf{Đặc tả kịch bản} \\
\hline
\endfirsthead
\hline
\textbf{Trường thông tin} & \textbf{Đặc tả kịch bản} \\
\hline
\endhead
Dữ liệu đầu vào (Input\newline Data) & Operator: OP001\newline Hub: HUB001\newline Vehicle: VEH004\newline Charging Port: CP003\newline Charging Time: 17:00–18:00\newline Tình huống 1: Hệ thống mô phỏng/cảm biến không phản hồi dữ liệu.\newline Tình huống 2: Phát sinh lỗi khi lưu ChargingSchedule, giữ chỗ\newline ChargingPort hoặcghi AuditLog. \\
\hline
Kết quả mong đợi\newline (Expected Result) & 1. Khi không truy xuất được dữ liệu, hệ thống thông báo lỗi và\newline không tạo lịch dựa trên dữ liệu không đầy đủ.\newline 2. Khi phát sinh lỗi trong quá trình lưu, giao dịch phải được rollback.\newline 3. Không tồn tại lịch sạc mới được lưu một phần.\newline 4. Trạng thái giữ chỗ của trụ sạc được khôi phục về trạng thái trước\newline giao dịch nếu giao dịch thất bại.\newline 5. Hệ thống thông báo cho Operator biết thao tác lưu lịch không\newline thành công. \\
\hline
Kết quả thực tế / Status\newline \textbf{.14. Test case cho UC14} & Pending – Chưa thực hiện kiểm thử.\newline \textbf{– Mô phỏng và phân tích kịch bản} \\
\hline
\textbf{\_.14.1. TC\_UC14\_01 – Chạ\_}\newline \textbf{Trường thông tin} & \textbf{\_y mô phỏng thành công và lưu báo cáo\_}\newline \textbf{Nội dung} \\
\hline
1. Test Case ID & TC\_UC14\_01 \\
\hline
2. Use Case liênquan & UC14 – Môphỏngvàphân tích kịch bản \\
\hline
3. Mô tả kịch bản & Đơn vị vận hành thiết lập kịch bản giả định về tình trạng quá tải tại\newline một MobilityHub, chạymôphỏngvà lưu báo cáo đánhgiá. \\
\hline
4. Tiền điều kiện\newline (Preconditions) & 1. Đơn vị vận hành đã đăng nhập thành công và có quyền truy cập\newline UC14.\newline 2. Trạng thái hiện tại của toàn mạng lưới đã được sao chép vào môi\newline trường mô phỏng ảo cách ly.\newline 3. Dữ liệu Hub và tài nguyên hợp lệ. \\
\hline
5. Dữ liệu đầu vào (Input\newline Data) & - Operator ID: OP001\newline - Scenario ID: SCN001\newline - Hub ID: HUB001\newline - Kịch bản: Lượng sinh viên tăng 30\%.\newline - Trạngthái ban đầu: 20 xe, 10 chỗ đỗ, 5 cổngsạc khả dụng. \\
\hline
6. Kết quả mong đợi\newline (Expected Result) & 1. Hệ thống tiếp nhận tham số và thực hiện mô phỏng thành công\newline trên môi trường ảo cách ly.\newline 2. Báo cáo hiển thị tác động của kịch bản và đề xuất điều phối phù\newline hợp.\newline 3. Báo cáo được lưu khi Operator chọn “Kết thúc”.\newline 4. Dữ liệu vận hành thực tế không bị thay đổi sau quá trình mô\newline phỏng. \\
\hline
7. Kếtquả thực tế / Status & Pending \\
\hline
\end{longtable}
\vspace{-0.5\baselineskip}
\subsection{Test case cho UC14 – Mô phỏng và phân tích kịch bản}
\subsubsection{TCUC1401 – Chạy mô phỏng thành công và lưu báo cáo}
\subsubsection{TCUC1402 – Chuyển tiếp áp dụng đề xuất mô phỏng (AF1)}
\begin{longtable}{|L{0.24\textwidth}|L{0.70\textwidth}|}
\hline
\textbf{Trường thông tin} & \textbf{Nội dung} \\
\hline
\endfirsthead
\hline
\textbf{Trường thông tin} & \textbf{Nội dung} \\
\hline
\endhead
1. Test Case ID & TC\_UC14\_02 \\
\hline
2. Use Case liênquan & UC14 – Môphỏngvàphân tích kịch bản \\
\hline
3. Mô tả kịch bản & Sau khi mô phỏng, Đơn vị vận hành đánh giá đề xuất điều phối là\newline khả thi và chọn “Ápdụngvào thực tế”. \\
\hline
4. Tiền điều kiện\newline (Preconditions) & 1. Đơn vị vận hành đã đăng nhập thành công và có quyền truy cập\newline UC14.\newline 2. Mô phỏng đã hoàn tất và báo cáo đề xuất đã được tạo.\newline 3. Đề xuất điềuphối có thể chuyển tiếpsangUC11 hoặc UC13. \\
\hline
5. Dữ liệu đầu vào (Input\newline Data) & - Operator ID: OP001\newline - Scenario ID: SCN002\newline - Recommendation: Điều chuyển 5 phương tiện từ HUB001 sang\newline HUB002.\newline - Use Case đích: UC11. \\
\hline
6. Kết quả mong đợi\newline (Expected Result) & 1. Hệ thống ghi nhận lựa chọn áp dụng đề xuất của Operator.\newline 2. Hệ thống chuyển hướng đến giao diện UC11 với các tham số điều\newline phối tương ứng được điền sẵn.\newline 3. Operator có thể kiểm tra và tiếp tục thực hiện quy trình UC11.\newline 4. Việc chuyển tiếp không tự động thay đổi trạng thái vận hành thực\newline tế khi chưa thực hiệnquytrình của UC11. \\
\hline
7. Kếtquả thực tế / Status & Pending \\
\hline
\end{longtable}
\vspace{-0.5\baselineskip}
\subsubsection{TCUC1403 – Nhập tham số mô phỏng không hợp lệ (EX1)}
\begin{longtable}{|L{0.24\textwidth}|L{0.70\textwidth}|}
\hline
\textbf{Trường thông tin} & \textbf{Nội dung} \\
\hline
\endfirsthead
\hline
\textbf{Trường thông tin} & \textbf{Nội dung} \\
\hline
\endhead
1. Test Case ID & TC\_UC14\_03 \\
\hline
2. Use Case liênquan & UC14 – Môphỏngvàphân tích kịch bản \\
\hline
3. Mô tả kịch bản & Đơn vị vận hành nhập tham số mô phỏng không hợp lệ, chẳng hạn\newline số lượng phươngtiệngiả định cógiá trị âm. \\
\hline
4. Tiền điều kiện\newline (Preconditions) & 1. Đơn vị vận hành đã đăng nhập thành công và có quyền truy cập\newline UC14.\newline 2. Giao diện thiết lậptham số môphỏngđangđược hiển thị. \\
\hline
5. Dữ liệu đầu vào (Input\newline Data) & - Operator ID: OP001\newline - Scenario ID: SCN003\newline - Hub ID: HUB001\newline - Số lượng phươngtiệngiả định: -5. \\
\hline
6. Kết quả mong đợi\newline (Expected Result) & 1. Hệ thống phát hiện tham số không hợp lệ.\newline 2. Hệ thống đánh dấu trường dữ liệu lỗi và hiển thị thông báo yêu\newline cầu nhập lại.\newline 3. Mô phỏng không được thực hiện khi tham số chưa hợp lệ.\newline 4. Operator có thể chỉnh sửa và nhậplại tham số. \\
\hline
7. Kếtquả thực tế / Status & Pending \\
\hline
\end{longtable}
\vspace{-0.5\baselineskip}
\subsubsection{TCUC1404 – Không thể khởi tạo môi trường mô phỏng (EX2)}
\begin{longtable}{|L{0.24\textwidth}|L{0.70\textwidth}|}
\hline
\textbf{Trường thông tin} & \textbf{Nội dung} \\
\hline
\endfirsthead
\hline
\textbf{Trường thông tin} & \textbf{Nội dung} \\
\hline
\endhead
1. Test Case ID & TC\_UC14\_04 \\
\hline
2. Use Case liênquan & UC14 – Môphỏngvàphân tích kịch bản \\
\hline
3. Mô tả kịch bản & Hệ thống không thể cấp phát tài nguyên để khởi tạo môi trường mô\newline phỏngcách lyvới cơ sở dữ liệu vận hành thực tế. \\
\hline
4. Tiền điều kiện\newline (Preconditions) & 1. Đơn vị vận hành đã đăng nhập thành công và có quyền truy cập\newline UC14.\newline 2. Operator đã thiết lập tham số kịch bản hợp lệ.\newline 3. Hệ thống gặp lỗi cấp phát tài nguyên khi khởi tạo môi trường mô\newline phỏng. \\
\hline
5. Dữ liệu đầu vào (Input\newline Data) & - Operator ID: OP001\newline - Scenario ID: SCN004\newline - Hub ID: HUB001\newline - Kịch bản: Đóng cửa HUB001 để đánh giá khả năng điều phối\newline phương tiện.\newline - Trạngthái môi trườngmôphỏng: Khôngthể khởi tạo. \\
\hline
6. Kết quả mong đợi\newline (Expected Result) & 1. Hệ thống phát hiện không thể khởi tạo môi trường mô phỏng cách\newline ly.\newline 2. Hệ thống hiển thị thông báo: “Không thể chạy mô phỏng do lỗi\newline cấp phát môi trường”.\newline 3. Quá trình mô phỏng kết thúc sớm và không tạo báo cáo kết quả.\newline 4. Trạngthái dữ liệu vận hành thực tế đượcgiữ nguyên. \\
\hline
7. Kết quả thực tế / Status & Pending \\
\hline
\end{longtable}
\vspace{-0.5\baselineskip}
\subsection{Test case cho UC15 – Tùy chỉnh giao diện (Customize UI)}
\subsubsection{TCUC1501 – Thay đổi giao diện thành công}
\begin{longtable}{|L{0.24\textwidth}|L{0.70\textwidth}|}
\hline
\textbf{Trường thông tin} & \textbf{Nội dung} \\
\hline
\endfirsthead
\hline
\textbf{Trường thông tin} & \textbf{Nội dung} \\
\hline
\endhead
Test Case ID & TC\_UC15\_01 \\
\hline
Use Case liênquan & UC15 – Customize UI \\
\hline
Mô tả kịch bản & Người dùngchuyển đổigiao diện từ sángsangtối hoặc ngược lại. \\
\hline
Tiền điều kiện & Người dùng đang truy cập hệ thống; giao diện hiện tại đã được khởi\newline tạo. \\
\hline
Dữ liệu đầu vào & Theme hiện tại: Light; Theme được chọn: Dark. \\
\hline
Kết quả mong đợi & Hệ thống áp dụng giao diện tối thành công; các thành phần giao diện\newline được cập nhật đồng bộ; tùy chọn giao diện được lưu để sử dụng cho\newline lần truycậptiếptheo. \\
\hline
Kếtquả thực tế / Status & Pending \\
\hline
\end{longtable}
\vspace{-0.5\baselineskip}
\subsubsection{TCUC1502 – Thay đổi ngôn ngữ và giữ nguyên trạng thái biểu mẫu}
\begin{longtable}{|L{0.24\textwidth}|L{0.70\textwidth}|}
\hline
\textbf{Trường thông tin} & \textbf{Nội dung} \\
\hline
\endfirsthead
\hline
\textbf{Trường thông tin} & \textbf{Nội dung} \\
\hline
\endhead
Test Case ID & TC\_UC15\_02 \\
\hline
Use Case liênquan & UC15 – Customize UI \\
\hline
\end{longtable}
\vspace{-0.5\baselineskip}
\begin{longtable}{|L{0.24\textwidth}|L{0.70\textwidth}|}
\hline
\textbf{Trường thông tin} & \textbf{Nội dung} \\
\hline
\endfirsthead
\hline
\textbf{Trường thông tin} & \textbf{Nội dung} \\
\hline
\endhead
Mô tả kịch bản & Người dùng chuyển đổi ngôn ngữ hệ thống khi đang thao tác trên\newline một biểu mẫu. \\
\hline
Tiền điều kiện & Người dùng đang truy cập hệ thống; biểu mẫu đang có dữ liệu được\newline nhậpnhưngchưagửi. \\
\hline
Dữ liệu đầu vào & Ngôn ngữ hiện tại: Tiếng Việt (VN); ngôn ngữ được chọn: Tiếng\newline Anh(EN); dữ liệu biểu mẫu: đã nhậpmột số trườngthôngtin. \\
\hline
Kết quả mong đợi & Hệ thống chuyển đổi ngôn ngữ thành công; nội dung giao diện được\newline cập nhật theo ngôn ngữ mới; dữ liệu đã nhập trong biểu mẫu được\newline giữ nguyên, khôngbị mất hoặc đặt lại. \\
\hline
Kếtquả thực tế / Status & Pending \\
\hline
\end{longtable}
\vspace{-0.5\baselineskip}
\subsubsection{TCUC1503 – Lỗi khi lưu tùy chọn giao diện}
\begin{longtable}{|L{0.24\textwidth}|L{0.70\textwidth}|}
\hline
\textbf{Trường thông tin} & \textbf{Nội dung} \\
\hline
\endfirsthead
\hline
\textbf{Trường thông tin} & \textbf{Nội dung} \\
\hline
\endhead
Test Case ID & TC\_UC15\_03 \\
\hline
Use Case liênquan & UC15 – Customize UI \\
\hline
Mô tả kịch bản & Người dùng thay đổi giao diện hoặc ngôn ngữ nhưng hệ thống gặp\newline lỗi khi lưu tùychọn. \\
\hline
Tiền điều kiện & Người dùng đang truy cập hệ thống; chức năng thay đổi giao diện và\newline ngôn ngữ hoạt động bình thường; phát sinh lỗi tại quá trình lưu tùy\newline chọn. \\
\hline
Dữ liệu đầu vào & Theme mới: Dark; Language mới: EN; trạngthái lưu: thất bại. \\
\hline
Kết quả mong đợi & Hệ thống hiển thị thông báo lỗi lưu tùy chọn; không ghi nhận tùy\newline chọn mới là đã lưu thành công; giữ trạng thái biểu mẫu và cho phép\newline người dùngthử lưu lại. \\
\hline
Kết quả thực tế / Status & Pending\newline 2.16. Test case cho UC16 - Cậpnhật trạngthái tài nguyên \\
\hline
\end{longtable}
\vspace{-0.5\baselineskip}
\subsection{Test case cho UC16 – Cập nhật trạng thái tài nguyên}
\begin{longtable}{|L{0.24\textwidth}|L{0.70\textwidth}|}
\hline
\textbf{Trường thông tin} & \textbf{Đặc tả kịch bản} \\
\hline
\endfirsthead
\hline
\textbf{Trường thông tin} & \textbf{Đặc tả kịch bản} \\
\hline
\endhead
Test Case ID & TC\_UC16\_01 \\
\hline
Use Case liênquan & UC16 – Cậpnhật trạngthái tài nguyên \\
\hline
Mô tả kịch bản & Thành công (Normal Flow – NFR02): Tiếp nhận gói dữ liệu mới\newline hợp lệ, cập nhật Operational DB và đẩy cập nhật thời gian thực <=\newline 15giây. \\
\hline
Tiền điều kiện\newline (Preconditions) & 1. Xe EB\_01 tồn tại trong Operational DB với status =\newline AVAILABLE, batteryLevel = 80\%, lastUpdatedAt = 2026-11-\newline 13T10:00:00.\newline 2. Giao diện Sinh viên và Operator đang kết nối RealtimeGateway\newline (WebSocket/SSE). \\
\hline
Dữ liệu đầu vào (Input\newline Data) & packet = \{ packetId: "PKT\_901", resourceId: "EB\_01",\newline resourceType: "SharedVehicle", status: IN\_USE, batteryLevel: 78.5,\newline latitude: 10.8701, longitude: 106.8031, packetTimestamp: 2026-11-\newline 13T10:00:10, isSimulation: false\} \\
\hline
\end{longtable}
\vspace{-0.5\baselineskip}
\begin{longtable}{|L{0.24\textwidth}|L{0.70\textwidth}|}
\hline
\textbf{Trường thông tin} & \textbf{Đặc tả kịch bản} \\
\hline
\endfirsthead
\hline
\textbf{Trường thông tin} & \textbf{Đặc tả kịch bản} \\
\hline
\endhead
Kết quả mong đợi & 1. packetTimestamp (10:00:10) > lastUpdatedAt (10:00:00) -> hợp\newline lệ.\newline 2. Operational DB cập nhật EB\_01: status = IN\_USE, batteryLevel =\newline 78.5, lastUpdatedAt = 10:00:10. \\
\hline
(Expected Result) & 3. RealtimeGateway phát sự kiện cập nhật tới giao diện người dùng\newline trong <= 15 giây (NFR02).\newline 4. Endpoint trả về HTTP 200 OK (SyncResult =\newline UPDATED\_OPERATIONAL). \\
\hline
Kết quả thực tế (Actual /\newline Status) & Pending \\
\hline
\end{longtable}
\vspace{-0.5\baselineskip}
\subsubsection{TCUC1602 – Ngoại lệ BẮT BUỘC (EX1 – NFR08): Lọc/loại bỏ gói dữ liệu trễ (Out-ofdate Timestamp) để không ghi đè dữ liệu mới hơn.}
\begin{longtable}{|L{0.24\textwidth}|L{0.70\textwidth}|}
\hline
\textbf{Trường thông tin} & \textbf{Đặc tả kịch bản} \\
\hline
\endfirsthead
\hline
\textbf{Trường thông tin} & \textbf{Đặc tả kịch bản} \\
\hline
\endhead
Test Case ID & TC\_UC16\_02 \\
\hline
Use Case liênquan & UC16 – Cậpnhật trạngthái tài nguyên \\
\hline
Mô tả kịch bản & Ngoại lệ BẮT BUỘC (EX1 – NFR08): Lọc/loại bỏ gói dữ liệu trễ\newline (Out-of-date Timestamp)để không ghi đè dữ liệu mới hơn. \\
\hline
Tiền điều kiện\newline (Preconditions) & 1. Xe EB\_01 trong Operational DB đang có trạng thái mới nhất:\newline status = IN\_USE, batteryLevel = 75.0\%, lastUpdatedAt = 2026-11-\newline 13T10:05:00.\newline 2. Một gói tin cũ bị trễ mạng (packetTimestamp = 10:04:30) hoặc\newline trùngmốc thờigian(10:05:00)mới đến nơi. \\
\hline
Dữ liệu đầu vào (Input\newline Data) & packet = \{ packetId: "PKT\_899", resourceId: "EB\_01",\newline resourceType: "SharedVehicle", status: AVAILABLE, batteryLevel:\newline 82.0,packetTimestamp: 2026-11-13T10:04:30, isSimulation: false\} \\
\hline
Kết quả mong đợi\newline (Expected Result) & 1. isOutdatedPacket(10:04:30, 10:05:00) trả về true.\newline 2. Hệ thống tự động hủy gói tin PKT\_899, ghi AuditLog với action\newline = DROP\_OUTDATED\_PACKET.\newline 3. Dữ liệu Operational DB của EB\_01 giữ nguyên tuyệt đối (status =\newline IN\_USE, batteryLevel = 75.0\%, lastUpdatedAt = 10:05:00), không\newline bị ghi đè.\newline 4. Không phát tín hiệu đẩy trạng thái cũ lên RealtimeGateway;\newline Endpoint trả về 200 OK kèm SyncResult =\newline DROPPED\_OUTDATED. \\
\hline
Kết quả thực tế (Actual /\newline Status) & Pending \\
\hline
\end{longtable}
\vspace{-0.5\baselineskip}
\subsubsection{TCUC1603 – Ngoại lệ (EX2): Gói tin sai định dạng (mức pin âm) hoặc resourceId không tồn tại trong danh mục.}
\begin{longtable}{|L{0.24\textwidth}|L{0.70\textwidth}|}
\hline
\textbf{Trường thông tin} & \textbf{Đặc tả kịch bản} \\
\hline
\endfirsthead
\hline
\textbf{Trường thông tin} & \textbf{Đặc tả kịch bản} \\
\hline
\endhead
Test Case ID & TC\_UC16\_03 \\
\hline
Use Case liênquan & UC16 – Cậpnhật trạngthái tài nguyên \\
\hline
Mô tả kịch bản & Ngoại lệ (EX2): Gói tin sai định dạng (mức pin âm) hoặc resourceId\newline khôngtồn tại trongdanh mục. \\
\hline
Tiền điều kiện & 1. Hệ thống đang sẵn sàng nhận gói tin từ Bộ giả lập dữ liệu. \\
\hline
(Preconditions) & 2. Mã tài nguyên "EB\_NOT\_FOUND" khôngtồn tại tronghệ thống. \\
\hline
\end{longtable}
\vspace{-0.5\baselineskip}
\begin{longtable}{|L{0.24\textwidth}|L{0.70\textwidth}|}
\hline
\textbf{Trường thông tin} & \textbf{Đặc tả kịch bản} \\
\hline
\endfirsthead
\hline
\textbf{Trường thông tin} & \textbf{Đặc tả kịch bản} \\
\hline
\endhead
Dữ liệu đầu vào (Input\newline Data) & Trường hợp A: packetId = "PKT\_ERR1", resourceId = "EB\_01",\newline batteryLevel = -10, packetTimestamp = 10:06:00.\newline Trường hợp B: packetId = "PKT\_ERR2", resourceId =\newline "EB\_NOT\_FOUND", status = AVAILABLE, batteryLevel = 90,\newline packetTimestamp= 10:06:00. \\
\hline
Kết quả mong đợi\newline (Expected Result) & 1. Trường hợp A ném InvalidPacketException; Trường hợp B ném\newline ResourceNotFoundException.\newline 2. Hệ thống ghi nhận lỗi tích hợp (INTEGRATION\_ERROR) vào\newline AuditLog và bỏ qua gói tin.\newline 3. Endpoint trả về HTTP 400 Bad Request cho Bộgiả lậpdữ liệu. \\
\hline
Kết quả thực tế (Actual / & Pending \\
\hline
Status) & 2.17. Test case cho UC17 - Xác thựcqua SSO \\
\hline
\end{longtable}
\vspace{-0.5\baselineskip}
\subsubsection{TCUC1604 – Luồng thay thế (AF1 – NFR07): Gói tin gắn cờ mô phỏng What-If chỉ ghi vào môi trường ảo cách ly, không ảnh hưởng Operational DB.}
\begin{longtable}{|L{0.24\textwidth}|L{0.70\textwidth}|}
\hline
\textbf{Trường thông tin} & \textbf{Đặc tả kịch bản} \\
\hline
\endfirsthead
\hline
\textbf{Trường thông tin} & \textbf{Đặc tả kịch bản} \\
\hline
\endhead
Test Case ID & TC\_UC16\_04 \\
\hline
Use Case liênquan & UC16 – Cậpnhật trạngthái tài nguyên \\
\hline
Mô tả kịch bản & Luồng thay thế (AF1 – NFR07): Gói tin gắn cờ mô phỏng What-If\newline chỉ ghi vào môi trường ảo cách ly, không ảnh hưởng Operational\newline DB. \\
\hline
Tiền điều kiện\newline (Preconditions) & 1. Xe EB\_01 trong Operational DB có status = AVAILABLE,\newline batteryLevel = 90.0\%, lastUpdatedAt = 10:00:00.\newline 2. SandboxStateStore cho kịch bản SCN\_01 đã được khởi tạo. \\
\hline
Dữ liệu đầu vào (Input\newline Data) & packet = \{ packetId: "PKT\_SIM\_01", resourceId: "EB\_01", status:\newline MAINTENANCE, batteryLevel: 12.0, packetTimestamp: 2026-11-\newline 13T10:10:00, isSimulation: true\} \\
\hline
Kết quả mong đợi\newline (Expected Result) & 1. Hệ thống điều hướng gói tin tới\newline SandboxStateStore.upsertSimulatedState(packet).\newline 2. Bản sao EB\_01 trong SandboxStateStore cập nhật status =\newline MAINTENANCE, batteryLevel = 12.0\%.\newline 3. Operational DB của EB\_01 vẫn giữ nguyên hoàn toàn (status =\newline AVAILABLE, batteryLevel = 90.0\%, lastUpdatedAt = 10:00:00)\newline (NFR07).\newline 4. Trả về HTTP 200 OK(SyncResult = UPDATED\_SANDBOX). \\
\hline
Kết quả thực tế (Actual /\newline Status) & Pending \\
\hline
\end{longtable}
\vspace{-0.5\baselineskip}
\subsection{Test case cho UC17 – Xác thực qua SSO}
\subsubsection{TCUC1701 – Đăng nhập SSO thành công}
\begin{longtable}{|L{0.24\textwidth}|L{0.70\textwidth}|}
\hline
\textbf{Trường thông tin} & \textbf{Đặc tả kịch bản} \\
\hline
\endfirsthead
\hline
\textbf{Trường thông tin} & \textbf{Đặc tả kịch bản} \\
\hline
\endhead
Test Case ID & TC\_UC17\_01 \\
\hline
Use Case liênquan & UC17 – Xác thựcqua SSO \\
\hline
Mô tả kịch bản & ĐăngnhậpSSO thành công \\
\hline
Tiền điều kiện\newline (Preconditions) & Tài khoản SSO hợp lệ và được ánh xạ quyền trong Hub. \\
\hline
\end{longtable}
\vspace{-0.5\baselineskip}
\begin{longtable}{|L{0.24\textwidth}|L{0.70\textwidth}|}
\hline
\textbf{Trường thông tin} & \textbf{Đặc tả kịch bản} \\
\hline
\endfirsthead
\hline
\textbf{Trường thông tin} & \textbf{Đặc tả kịch bản} \\
\hline
\endhead
Dữ liệu đầu vào (Input\newline Data) & Lần 1: tài khoản sinh viên An.\newline Lần 2: tài khoản nhân viên vận hành.\newline Thao tác: đăngnhậpSSO bằngtừngtài khoản. \\
\hline
Kết quả mong đợi\newline (Expected Result) & 1. Tạo phiên với đúng vai trò.\newline 2. Hiển thị giao diện tương ứng; ghi lịch sử đăng nhập.\newline 3. Không yêu cầu nhậplại thôngtin nếu SSO đã cóphiên hợp lệ. \\
\hline
Kếtquả thực tế / Status & Pending– Chưa thực hiện kiểm thử. \\
\hline
\end{longtable}
\vspace{-0.5\baselineskip}
\subsubsection{TCUC1702 – Hủy đăng nhập}
\begin{longtable}{|L{0.24\textwidth}|L{0.70\textwidth}|}
\hline
\textbf{Trường thông tin} & \textbf{Đặc tả kịch bản} \\
\hline
\endfirsthead
\hline
\textbf{Trường thông tin} & \textbf{Đặc tả kịch bản} \\
\hline
\endhead
Test Case ID & TC\_UC17\_02 \\
\hline
Use Case liênquan & UC17 – Xác thựcqua SSO \\
\hline
Mô tả kịch bản & Hủyđăngnhập \\
\hline
Tiền điều kiện\newline (Preconditions) & Đang ở bước đăng nhập SSO; có thể có lựa chọn xe tạm. \\
\hline
Dữ liệu đầu vào (Input\newline Data) & Thao tác: chọn Hủy hoặc Từ chối tại SSO. \\
\hline
Kết quả mong đợi\newline (Expected Result) & 1. Không tạo phiên mới.\newline 2. Trở về trangchủ; xóa lựa chọn đặt xe tạm nếu có. \\
\hline
Kếtquả thực tế / Status & Pending– Chưa thực hiện kiểm thử. \\
\hline
\end{longtable}
\vspace{-0.5\baselineskip}
\subsubsection{TCUC1703 – Thông tin sai hoặc tài khoản bị khóa}
\begin{longtable}{|L{0.24\textwidth}|L{0.70\textwidth}|}
\hline
\textbf{Trường thông tin} & \textbf{Đặc tả kịch bản} \\
\hline
\endfirsthead
\hline
\textbf{Trường thông tin} & \textbf{Đặc tả kịch bản} \\
\hline
\endhead
Test Case ID & TC\_UC17\_03 \\
\hline
Use Case liênquan & UC17 – Xác thựcqua SSO \\
\hline
Mô tả kịch bản & Thôngtin sai hoặc tài khoản bị khóa \\
\hline
Tiền điều kiện\newline (Preconditions) & SSO được cấu hình trả kết quả từ chối. \\
\hline
Dữ liệu đầu vào (Input\newline Data) & Lần 1: nhập sai thông tin đăng nhập.\newline Lần 2: dùngtài khoản bị khóa. \\
\hline
Kết quả mong đợi & 1. Hiển thị lỗi phù hợp. \\
\hline
(Expected Result) & 2. Khôngtạophiên Hub; khôngcấp quyền truycập. \\
\hline
Kếtquả thực tế / Status & Pending– Chưa thực hiện kiểm thử. \\
\hline
\end{longtable}
\vspace{-0.5\baselineskip}
\subsubsection{TCUC1704 – SSO không phản hồi}
\begin{longtable}{|L{0.24\textwidth}|L{0.70\textwidth}|}
\hline
\textbf{Trường thông tin} & \textbf{Đặc tả kịch bản} \\
\hline
\endfirsthead
\hline
\textbf{Trường thông tin} & \textbf{Đặc tả kịch bản} \\
\hline
\endhead
Test Case ID & TC\_UC17\_04 \\
\hline
Use Case liênquan & UC17 – Xác thựcqua SSO \\
\hline
Mô tả kịch bản & SSO không phản hồi \\
\hline
Tiền điều kiện\newline (Preconditions) & Mô phỏng SSO mất kết nối hoặc không trả lời trong thời gian chờ đã\newline cấu hình. \\
\hline
Dữ liệu đầu vào (Input\newline Data) & Thao tác: bắt đầu đăng nhập hoặc gửi kết quả để xác minh. \\
\hline
Kết quả mong đợi\newline (Expected Result) & 1. Thông báo tạm thời không thể đăng nhập và cho phép thử lại.\newline 2. Khôngtạophiên mới khi chưa xác minh thành công. \\
\hline
Kếtquả thực tế / Status & Pending– Chưa thực hiện kiểm thử. \\
\hline
\end{longtable}
\vspace{-0.5\baselineskip}
\normalsize